\chapter{Presidente Hayes}\label{dept:Presidente Hayes}
\mapaparaguai{15}{Presidente Hayes}
\vfill\clearpage
\section{Presidente Hayes}

\subsection{Ficha climática de Villa Hayes (capital
departamental)}

\textbf{Irradiância solar global --- (kWh/\allowbreak m²/\allowbreak dia):}

{\def\LTcaptype{none} % do not increment counter
\begin{longtable}[]{@{}
  >{\raggedright\arraybackslash}p{(\linewidth - 22\tabcolsep) * \real{0.0833}}
  >{\raggedright\arraybackslash}p{(\linewidth - 22\tabcolsep) * \real{0.0833}}
  >{\raggedright\arraybackslash}p{(\linewidth - 22\tabcolsep) * \real{0.0833}}
  >{\raggedright\arraybackslash}p{(\linewidth - 22\tabcolsep) * \real{0.0833}}
  >{\raggedright\arraybackslash}p{(\linewidth - 22\tabcolsep) * \real{0.0833}}
  >{\raggedright\arraybackslash}p{(\linewidth - 22\tabcolsep) * \real{0.0833}}
  >{\raggedright\arraybackslash}p{(\linewidth - 22\tabcolsep) * \real{0.0833}}
  >{\raggedright\arraybackslash}p{(\linewidth - 22\tabcolsep) * \real{0.0833}}
  >{\raggedright\arraybackslash}p{(\linewidth - 22\tabcolsep) * \real{0.0833}}
  >{\raggedright\arraybackslash}p{(\linewidth - 22\tabcolsep) * \real{0.0833}}
  >{\raggedright\arraybackslash}p{(\linewidth - 22\tabcolsep) * \real{0.0833}}
  >{\raggedright\arraybackslash}p{(\linewidth - 22\tabcolsep) * \real{0.0833}}@{}}
\toprule\noalign{}
\begin{minipage}[b]{\linewidth}\raggedright
Jan
\end{minipage} & \begin{minipage}[b]{\linewidth}\raggedright
Fev
\end{minipage} & \begin{minipage}[b]{\linewidth}\raggedright
Mar
\end{minipage} & \begin{minipage}[b]{\linewidth}\raggedright
Abr
\end{minipage} & \begin{minipage}[b]{\linewidth}\raggedright
Mai
\end{minipage} & \begin{minipage}[b]{\linewidth}\raggedright
Jun
\end{minipage} & \begin{minipage}[b]{\linewidth}\raggedright
Jul
\end{minipage} & \begin{minipage}[b]{\linewidth}\raggedright
Ago
\end{minipage} & \begin{minipage}[b]{\linewidth}\raggedright
Set
\end{minipage} & \begin{minipage}[b]{\linewidth}\raggedright
Out
\end{minipage} & \begin{minipage}[b]{\linewidth}\raggedright
Nov
\end{minipage} & \begin{minipage}[b]{\linewidth}\raggedright
Dez
\end{minipage} \\
\midrule\noalign{}
\endhead
\bottomrule\noalign{}
\endlastfoot
6.76 & 6.20 & 5.42 & 4.44 & 3.36 & 2.83 & 3.21 & 3.93 & 4.61 & 5.49 &
6.42 & 6.78 \\
\end{longtable}
}

Média anual: 4.95 kWh/\allowbreak m²/\allowbreak dia. Inclinação recomendada: 25° Norte (anual);
35° inverno; 15° verão.

\textbf{Precipitação --- (mm/\allowbreak mês):}

{\def\LTcaptype{none} % do not increment counter
\begin{longtable}[]{@{}
  >{\raggedright\arraybackslash}p{(\linewidth - 22\tabcolsep) * \real{0.0833}}
  >{\raggedright\arraybackslash}p{(\linewidth - 22\tabcolsep) * \real{0.0833}}
  >{\raggedright\arraybackslash}p{(\linewidth - 22\tabcolsep) * \real{0.0833}}
  >{\raggedright\arraybackslash}p{(\linewidth - 22\tabcolsep) * \real{0.0833}}
  >{\raggedright\arraybackslash}p{(\linewidth - 22\tabcolsep) * \real{0.0833}}
  >{\raggedright\arraybackslash}p{(\linewidth - 22\tabcolsep) * \real{0.0833}}
  >{\raggedright\arraybackslash}p{(\linewidth - 22\tabcolsep) * \real{0.0833}}
  >{\raggedright\arraybackslash}p{(\linewidth - 22\tabcolsep) * \real{0.0833}}
  >{\raggedright\arraybackslash}p{(\linewidth - 22\tabcolsep) * \real{0.0833}}
  >{\raggedright\arraybackslash}p{(\linewidth - 22\tabcolsep) * \real{0.0833}}
  >{\raggedright\arraybackslash}p{(\linewidth - 22\tabcolsep) * \real{0.0833}}
  >{\raggedright\arraybackslash}p{(\linewidth - 22\tabcolsep) * \real{0.0833}}@{}}
\toprule\noalign{}
\begin{minipage}[b]{\linewidth}\raggedright
Jan
\end{minipage} & \begin{minipage}[b]{\linewidth}\raggedright
Fev
\end{minipage} & \begin{minipage}[b]{\linewidth}\raggedright
Mar
\end{minipage} & \begin{minipage}[b]{\linewidth}\raggedright
Abr
\end{minipage} & \begin{minipage}[b]{\linewidth}\raggedright
Mai
\end{minipage} & \begin{minipage}[b]{\linewidth}\raggedright
Jun
\end{minipage} & \begin{minipage}[b]{\linewidth}\raggedright
Jul
\end{minipage} & \begin{minipage}[b]{\linewidth}\raggedright
Ago
\end{minipage} & \begin{minipage}[b]{\linewidth}\raggedright
Set
\end{minipage} & \begin{minipage}[b]{\linewidth}\raggedright
Out
\end{minipage} & \begin{minipage}[b]{\linewidth}\raggedright
Nov
\end{minipage} & \begin{minipage}[b]{\linewidth}\raggedright
Dez
\end{minipage} \\
\midrule\noalign{}
\endhead
\bottomrule\noalign{}
\endlastfoot
123 & 149 & 130 & 143 & 143 & 71 & 59 & 33 & 70 & 161 & 197 & 182 \\
\end{longtable}
}

Média anual: \textasciitilde1.460 mm/\allowbreak ano. Estação chuvosa Out-Nov; seca
Jul-Ago.
\clearpage
\subsection*{Dados Departamentais}
\subsubsection{Desenvolvimento Humano e
Social}

\subsection{IDH Departamental}

{\def\LTcaptype{none} % do not increment counter
\begin{longtable}[]{@{}ll@{}}
\toprule\noalign{}
Indicador & Valor \\
\midrule\noalign{}
\endhead
\bottomrule\noalign{}
\endlastfoot
IDH & 0,693 \\
Classificação IDH & médio \\
Ranking nacional (1=melhor) & 13/\allowbreak 18 \\
Esperança de vida & 71,8 anos \\
Anos de escolaridade média & 7,8 anos \\
RNB per capita (USD PPA) & 7.900 \\
\end{longtable}
}

\subsection{Indicadores de Pobreza}

{\def\LTcaptype{none} % do not increment counter
\begin{longtable}[]{@{}ll@{}}
\toprule\noalign{}
Indicador & Valor \\
\midrule\noalign{}
\endhead
\bottomrule\noalign{}
\endlastfoot
Taxa de pobreza total (\%) & 29,4\% \\
Taxa de pobreza extrema (\%) & 7,2\% \\
Índice de Gini & 0,453 \\
\end{longtable}
}

\subsection{Infraestrutura Básica}

{\def\LTcaptype{none} % do not increment counter
\begin{longtable}[]{@{}ll@{}}
\toprule\noalign{}
Indicador & Valor \\
\midrule\noalign{}
\endhead
\bottomrule\noalign{}
\endlastfoot
Acesso a água potável (\%) & 71,3\% \\
Acesso a saneamento (\%) & 68,9\% \\
\end{longtable}
}

\subsubsection{Segurança Pública}

\subsection{Indicadores}

{\def\LTcaptype{none} % do not increment counter
\begin{longtable}[]{@{}
  >{\raggedright\arraybackslash}p{(\linewidth - 2\tabcolsep) * \real{0.6111}}
  >{\raggedright\arraybackslash}p{(\linewidth - 2\tabcolsep) * \real{0.3889}}@{}}
\toprule\noalign{}
\begin{minipage}[b]{\linewidth}\raggedright
Indicador
\end{minipage} & \begin{minipage}[b]{\linewidth}\raggedright
Valor
\end{minipage} \\
\midrule\noalign{}
\endhead
\bottomrule\noalign{}
\endlastfoot
Taxa de homicídios (por 100k hab) & 5,5--5,6 (2019--2020, fonte INE\footnote{\fineINE} ---
próxima à média nacional) \\
Crimes registrados (total/\allowbreak ano) & baixo em números absolutos (população
esparsa) \\
Número de comisarías & \textasciitilde8 (estimativa --- vasto território
com baixa densidade) \\
Índice segurança relativo & médio \\
\end{longtable}
}

\subsubsection{Mercado de Terra Rural}

\subsection{Preços por Tipo de Terra}

{\def\LTcaptype{none} % do not increment counter
\begin{longtable}[]{@{}llll@{}}
\toprule\noalign{}
Tipo & Preço (USD/\allowbreak ha) & Faixa & Tendência \\
\midrule\noalign{}
\endhead
\bottomrule\noalign{}
\endlastfoot
Agrícola alta produtividade & 2.500 & 1.500 -- 4.000 & ↑ \\
Pastagem (campo nativo) & 1.200 & 700 -- 2.000 & ↑ \\
Terra bruta /\allowbreak  arbustiva & 500 & 250 -- 900 & → \\
\end{longtable}
}

\subsection{Características do
Mercado}

{\def\LTcaptype{none} % do not increment counter
\begin{longtable}[]{@{}ll@{}}
\toprule\noalign{}
Parâmetro & Valor \\
\midrule\noalign{}
\endhead
\bottomrule\noalign{}
\endlastfoot
Valorização anual média (\%) & 5--8\% \\
Lote mínimo rural (ha) & 20 (zona de chaco central) \\
Imposto territorial rural (\%) & 1\% sobre valor fiscal \\
Liquidez do mercado & média \\
\end{longtable}
}
\clearpage
\newpage
\secaoDiagnostico{Benjamin Aceval}{\mapaDistrito{1502}}{Benjamin Aceval (Presidente Hayes) apresenta perfil operacional
consolidado para comparacao territorial, com leitura conjunta de
infraestrutura, risco hidroclimatico e capacidade de continuidade de
servicos essenciais.

\subsubsection{Por que sim}

\subsubsection{Por que não}

\subsubsection{Pre-condicoes Minimas de
Decisao}

\begin{itemize}
\tightlist
\item
  Atualizar indicadores criticos em ciclo trimestral.
\item
  Confirmar trafegabilidade e contingencia hidrica/\allowbreak energetica local.
\item
  Validar checklist de resiliencia domiciliar/\allowbreak logistica antes de decisao
  final.
\end{itemize}

\subsubsection{Recomendação
Técnica}

Uso recomendado como candidato em analise multicriterio, condicionado ao
cumprimento das pre-condicoes acima. Referência atual de score: GSS 8.2;
risco predominante: baixo.

\begin{itemize}
\tightlist
\item
  \begin{enumerate}
  \def\labelenumi{\arabic{enumi})}
  \tightlist
  \item
    Delimitacao territorial validada por georreferencia institucional.
  \end{enumerate}
\item
  \begin{enumerate}
  \def\labelenumi{\arabic{enumi})}
  \setcounter{enumi}{1}
  \tightlist
  \item
    População municipal/\allowbreak distrital referenciada por base censitaria
    nacional.
  \end{enumerate}
\item
  \begin{enumerate}
  \def\labelenumi{\arabic{enumi})}
  \setcounter{enumi}{2}
  \tightlist
  \item
    Comparabilidade intra-departamental apoiada em indicadores
    distritais oficiais.
  \end{enumerate}
\item
  \begin{enumerate}
  \def\labelenumi{\arabic{enumi})}
  \setcounter{enumi}{3}
  \tightlist
  \item
    Estado macro de conectividade viaria ancorado em informacoes de
    rodovias.
  \end{enumerate}
\item
  \begin{enumerate}
  \def\labelenumi{\arabic{enumi})}
  \setcounter{enumi}{4}
  \tightlist
  \item
    Exposicao hidrometeorologica monitorada por avisos oficiais.
  \end{enumerate}
\item
  \begin{enumerate}
  \def\labelenumi{\arabic{enumi})}
  \setcounter{enumi}{5}
  \tightlist
  \item
    Historico de contexto meteorologico apoiado em publicacoes tecnicas.
  \end{enumerate}
\item
  \begin{enumerate}
  \def\labelenumi{\arabic{enumi})}
  \setcounter{enumi}{6}
  \tightlist
  \item
    Capacidade institucional de resposta a eventos apoiada em acoes da
    SEN\footnote{\fineSEN}.
  \end{enumerate}
\item
  \begin{enumerate}
  \def\labelenumi{\arabic{enumi})}
  \setcounter{enumi}{7}
  \tightlist
  \item
    Infraestrutura energetica analisada via operador nacional.
  \end{enumerate}
\item
  \begin{enumerate}
  \def\labelenumi{\arabic{enumi})}
  \setcounter{enumi}{8}
  \tightlist
  \item
    Referências de obras e logistica nacional verificadas no MOPC.
  \end{enumerate}
\item
  \begin{enumerate}
  \def\labelenumi{\arabic{enumi})}
  \setcounter{enumi}{9}
  \tightlist
  \item
    Contexto sociopolitico institucional referenciado por autoridade
    eleitoral.
  \end{enumerate}
\item
  \begin{enumerate}
  \def\labelenumi{\arabic{enumi})}
  \setcounter{enumi}{10}
  \tightlist
  \item
    Camada cartografica de apoio eleitoral/\allowbreak territorial consultada.
  \end{enumerate}
\item
  \begin{enumerate}
  \def\labelenumi{\arabic{enumi})}
  \setcounter{enumi}{11}
  \tightlist
  \item
    Dados publicos complementares para verificacao cruzada.
  \end{enumerate}
\end{itemize}}{dist:15_Presidente_Hayes:Benjamin_Aceval}
\begin{table}[ht!]
\small \begin{tabular}{p{3.0cm}p{0.8cm}p{6.0cm}} \toprule
\textbf{Critério} & \textbf{Nota} & \textbf{Justificativa} \\ \midrule
A - Ameaças Estratégicas & 8.2 & - \\
B - Risco Social & 9.0 & - \\
C - Riscos Naturais & 6.9 & - \\
D - Autossuficiência & 8.0 & - \\
E - Institucional & 8.7 & - \\
\midrule \textbf{GSS FINAL} & \textbf{8.2} & \textbf{Seguro (bom potencial com preparacao).} \\ \bottomrule \end{tabular}
\end{table}
\subsubsection{Vulnerabilidades e
Mitigação}

\begin{itemize}
\tightlist
\item
  Exposicao hidrometeorologica controlada (\textless45/\allowbreak 100).
\item
  Conectividade viaria favoravel (\textgreater=70/\allowbreak 100).
\item
  Estoque de contingencia para 14 dias (agua/\allowbreak alimentos/\allowbreak combustivel),
  priorizado quando risco hidro=39/\allowbreak 100.
\item
  Duplicar rotas/\allowbreak logistica quando conectividade viaria=79/\allowbreak 100 for
  \textless70; manter rota secundaria mapeada e testada trimestralmente.
\item
  Plano de energia resiliente com autonomia minima de 72h
  (geracao/\allowbreak backup), revisao semestral.
\end{itemize}

\subsubsection{Dossiê de Campo}
\begin{description}[style=nextline,leftmargin=0.5cm,font=\bfseries]
\item[Coordenadas]  25°00'S 57°33'W (geocodificado via Nominatim).
\end{description}
\textbf{Fonte climática:} NASA POWER Climatology API (período 2001-2020)

\textbf{Fonte luminosa:} estimativa world atlas

\paragraph{Irradiação Solar
(kWh/\allowbreak m²/\allowbreak dia)}

{\def\LTcaptype{none} % do not increment counter
\begin{longtable}[]{@{}
  >{\raggedright\arraybackslash}p{(\linewidth - 24\tabcolsep) * \real{0.0746}}
  >{\raggedright\arraybackslash}p{(\linewidth - 24\tabcolsep) * \real{0.0746}}
  >{\raggedright\arraybackslash}p{(\linewidth - 24\tabcolsep) * \real{0.0746}}
  >{\raggedright\arraybackslash}p{(\linewidth - 24\tabcolsep) * \real{0.0746}}
  >{\raggedright\arraybackslash}p{(\linewidth - 24\tabcolsep) * \real{0.0746}}
  >{\raggedright\arraybackslash}p{(\linewidth - 24\tabcolsep) * \real{0.0746}}
  >{\raggedright\arraybackslash}p{(\linewidth - 24\tabcolsep) * \real{0.0746}}
  >{\raggedright\arraybackslash}p{(\linewidth - 24\tabcolsep) * \real{0.0746}}
  >{\raggedright\arraybackslash}p{(\linewidth - 24\tabcolsep) * \real{0.0746}}
  >{\raggedright\arraybackslash}p{(\linewidth - 24\tabcolsep) * \real{0.0746}}
  >{\raggedright\arraybackslash}p{(\linewidth - 24\tabcolsep) * \real{0.0746}}
  >{\raggedright\arraybackslash}p{(\linewidth - 24\tabcolsep) * \real{0.0746}}
  >{\raggedright\arraybackslash}p{(\linewidth - 24\tabcolsep) * \real{0.1045}}@{}}
\toprule\noalign{}
\begin{minipage}[b]{\linewidth}\raggedright
Jan
\end{minipage} & \begin{minipage}[b]{\linewidth}\raggedright
Fev
\end{minipage} & \begin{minipage}[b]{\linewidth}\raggedright
Mar
\end{minipage} & \begin{minipage}[b]{\linewidth}\raggedright
Abr
\end{minipage} & \begin{minipage}[b]{\linewidth}\raggedright
Mai
\end{minipage} & \begin{minipage}[b]{\linewidth}\raggedright
Jun
\end{minipage} & \begin{minipage}[b]{\linewidth}\raggedright
Jul
\end{minipage} & \begin{minipage}[b]{\linewidth}\raggedright
Ago
\end{minipage} & \begin{minipage}[b]{\linewidth}\raggedright
Set
\end{minipage} & \begin{minipage}[b]{\linewidth}\raggedright
Out
\end{minipage} & \begin{minipage}[b]{\linewidth}\raggedright
Nov
\end{minipage} & \begin{minipage}[b]{\linewidth}\raggedright
Dez
\end{minipage} & \begin{minipage}[b]{\linewidth}\raggedright
Média
\end{minipage} \\
\midrule\noalign{}
\endhead
\bottomrule\noalign{}
\endlastfoot
6.64 & 6.12 & 5.48 & 4.54 & 3.38 & 2.87 & 3.21 & 3.93 & 4.62 & 5.45 &
6.37 & 6.67 & \textbf{4.94} \\
\end{longtable}
}

\textbf{Inclinação solar recomendada:} 25° N (anual) · 35° N (inverno
jun-ago) · 15° N (verão nov-jan)

\paragraph{Precipitação (mm/\allowbreak mês)}

{\def\LTcaptype{none} % do not increment counter
\begin{longtable}[]{@{}
  >{\raggedright\arraybackslash}p{(\linewidth - 24\tabcolsep) * \real{0.0704}}
  >{\raggedright\arraybackslash}p{(\linewidth - 24\tabcolsep) * \real{0.0704}}
  >{\raggedright\arraybackslash}p{(\linewidth - 24\tabcolsep) * \real{0.0704}}
  >{\raggedright\arraybackslash}p{(\linewidth - 24\tabcolsep) * \real{0.0704}}
  >{\raggedright\arraybackslash}p{(\linewidth - 24\tabcolsep) * \real{0.0704}}
  >{\raggedright\arraybackslash}p{(\linewidth - 24\tabcolsep) * \real{0.0704}}
  >{\raggedright\arraybackslash}p{(\linewidth - 24\tabcolsep) * \real{0.0704}}
  >{\raggedright\arraybackslash}p{(\linewidth - 24\tabcolsep) * \real{0.0704}}
  >{\raggedright\arraybackslash}p{(\linewidth - 24\tabcolsep) * \real{0.0704}}
  >{\raggedright\arraybackslash}p{(\linewidth - 24\tabcolsep) * \real{0.0704}}
  >{\raggedright\arraybackslash}p{(\linewidth - 24\tabcolsep) * \real{0.0704}}
  >{\raggedright\arraybackslash}p{(\linewidth - 24\tabcolsep) * \real{0.0704}}
  >{\raggedright\arraybackslash}p{(\linewidth - 24\tabcolsep) * \real{0.1549}}@{}}
\toprule\noalign{}
\begin{minipage}[b]{\linewidth}\raggedright
Jan
\end{minipage} & \begin{minipage}[b]{\linewidth}\raggedright
Fev
\end{minipage} & \begin{minipage}[b]{\linewidth}\raggedright
Mar
\end{minipage} & \begin{minipage}[b]{\linewidth}\raggedright
Abr
\end{minipage} & \begin{minipage}[b]{\linewidth}\raggedright
Mai
\end{minipage} & \begin{minipage}[b]{\linewidth}\raggedright
Jun
\end{minipage} & \begin{minipage}[b]{\linewidth}\raggedright
Jul
\end{minipage} & \begin{minipage}[b]{\linewidth}\raggedright
Ago
\end{minipage} & \begin{minipage}[b]{\linewidth}\raggedright
Set
\end{minipage} & \begin{minipage}[b]{\linewidth}\raggedright
Out
\end{minipage} & \begin{minipage}[b]{\linewidth}\raggedright
Nov
\end{minipage} & \begin{minipage}[b]{\linewidth}\raggedright
Dez
\end{minipage} & \begin{minipage}[b]{\linewidth}\raggedright
Total/\allowbreak ano
\end{minipage} \\
\midrule\noalign{}
\endhead
\bottomrule\noalign{}
\endlastfoot
123 & 149 & 130 & 143 & 143 & 71 & 59 & 33 & 70 & 161 & 197 & 182 &
\textbf{1465 mm} \\
\end{longtable}
}

\paragraph{Poluição Luminosa}

{\def\LTcaptype{none} % do not increment counter
\begin{longtable}[]{@{}ll@{}}
\toprule\noalign{}
Parâmetro & Valor \\
\midrule\noalign{}
\endhead
\bottomrule\noalign{}
\endlastfoot
Escala Bortle & 2 --- Céu tipicamente escuro \\
Radiância artificial & 0.3 nW/\allowbreak cm²/\allowbreak sr \\
\end{longtable}
}
\paragraph{Solo (SoilGrids 2.0, média ponderada 0--30
cm)}

{\def\LTcaptype{none} % do not increment counter
\begin{longtable}[]{@{}ll@{}}
\toprule\noalign{}
Parâmetro & Valor \\
\midrule\noalign{}
\endhead
\bottomrule\noalign{}
\endlastfoot
pH (H$_2$O) & 5.95 \\
Carbono orgânico (SOC) & 14.02 g/\allowbreak kg \\
Argila & 28.85 \% \\
Areia & 39.68 \% \\
Silte (calc.) & 31.5 \% \\
Densidade aparente & 1.37 g/\allowbreak cm³ \\
\textbf{Aptidão agrícola} & \textbf{Alta} \\
\end{longtable}
}

\subsubsection{Indicadores Sociais}

\textbf{Fontes:} PNUD 2020, INE\footnote{\fineINE} Censo 2022, INE\footnote{\fineINE} EPHC 2023, Ministerio
Público 2024\\
\textbf{Nota:} Valores marcados como \emph{(dept.)} referem-se ao
departamento de Presidente Hayes; valores \emph{(dist.)} são específicos
deste distrito. Dados marcados \emph{(est.)} são estimativas.

{\def\LTcaptype{none} % do not increment counter
\begin{longtable}[]{@{}
  >{\raggedright\arraybackslash}p{(\linewidth - 4\tabcolsep) * \real{0.4231}}
  >{\raggedright\arraybackslash}p{(\linewidth - 4\tabcolsep) * \real{0.2692}}
  >{\raggedright\arraybackslash}p{(\linewidth - 4\tabcolsep) * \real{0.3077}}@{}}
\toprule\noalign{}
\begin{minipage}[b]{\linewidth}\raggedright
Indicador
\end{minipage} & \begin{minipage}[b]{\linewidth}\raggedright
Valor
\end{minipage} & \begin{minipage}[b]{\linewidth}\raggedright
Âmbito
\end{minipage} \\
\midrule\noalign{}
\endhead
\bottomrule\noalign{}
\endlastfoot
IDH (2020) & 0,693 & dept. \\
Ranking IDH nacional & 13/\allowbreak 18 & dept. \\
Esperança de vida & 71,8 anos & dept. \\
Escolaridade média & 9.4 anos & dist. \\
RNB per capita (USD PPA) & 7.900 & dept. \\
Pobreza monetária (\%) & 29,4\% & dept. \\
Pobreza extrema (\%) & 7,2\% & dept. \\
Índice de Gini & 0,453 & dept. \\
Acesso a água potável (\%) & 71,3\% & dept. \\
Acesso a saneamento (\%) & 68,9\% & dept. \\
Taxa de homicídios (est., /\allowbreak 100k hab) & 5,5--5,6 (2019--2020, fonte INE
--- próxima à média nacional) & dept. \\
Índice de segurança & médio & dept. \\
População (Censo 2022) & N/\allowbreak D & dist. \\
Idade mediana & N/\allowbreak D & dist. \\
Taxa de fecundidade & N/\allowbreak D & dist. \\
\end{longtable}
}
\subsubsection{Combustível}

\textbf{Referência:} PETROPAR /\allowbreak  postos locais (2024) \textbf{Tipo de
localidade:} Interior

{\def\LTcaptype{none} % do not increment counter
\begin{longtable}[]{@{}lll@{}}
\toprule\noalign{}
Combustível & USD/\allowbreak litro & Gs/\allowbreak litro (aprox.) \\
\midrule\noalign{}
\endhead
\bottomrule\noalign{}
\endlastfoot
Gasolina 93 oct & 1.05 & 7,770 \\
Gasolina 97 oct (premium) & 1.15 & 8,510 \\
Diesel & 0.98 & 7,252 \\
\end{longtable}
}

\begin{quote}
Preços podem variar ±5\% conforme posto e sazonalidade. Chaco e interior
remoto apresentam maior variação.
\end{quote}

\subsubsection{Cobertura Celular}

\textbf{Fonte:} CONATEL PY /\allowbreak  operadoras (2024)

{\def\LTcaptype{none} % do not increment counter
\begin{longtable}[]{@{}ll@{}}
\toprule\noalign{}
Parâmetro & Valor \\
\midrule\noalign{}
\endhead
\bottomrule\noalign{}
\endlastfoot
Cobertura 4G população (dept.) & 60\% \\
Cobertura 4G área rural & 35\% \\
Melhor operadora & Tigo \\
Qualidade rural & limitada \\
\end{longtable}
}

\begin{quote}
Para áreas rurais fora do núcleo urbano, recomenda-se chip Tigo como
principal e Personal como backup.
\end{quote}

\subsubsection{Internet}

\textbf{Fonte:} CONATEL /\allowbreak  Speedtest Ookla (2024)

{\def\LTcaptype{none} % do not increment counter
\begin{longtable}[]{@{}ll@{}}
\toprule\noalign{}
Parâmetro & Valor \\
\midrule\noalign{}
\endhead
\bottomrule\noalign{}
\endlastfoot
Velocidade média download & 25 Mbps \\
Domicílios com internet (dept.) & 35\% \\
Tecnologia predominante & rádio/\allowbreak satélite \\
Opção rural & Starlink disponível (\textasciitilde USD 44/\allowbreak mês) \\
\end{longtable}
}

\subsubsection{Mercado Imobiliário e Terra
Rural}

\textbf{Fonte:} INDERT /\allowbreak  Clasificados.com.py (2024)

{\def\LTcaptype{none} % do not increment counter
\begin{longtable}[]{@{}ll@{}}
\toprule\noalign{}
Tipo & Referência \\
\midrule\noalign{}
\endhead
\bottomrule\noalign{}
\endlastfoot
Terra agrícola alta prod. (USD/\allowbreak ha) & 2,000 \\
Imóvel urbano (USD/\allowbreak m²) & 600 \\
Aluguel 2 quartos (USD/\allowbreak mês) & 250 \\
\end{longtable}
}

\begin{quote}
Valores de referência departamental. Localidades menores podem ter
preços 20--40\% abaixo da capital departamental.
\end{quote}

\subsubsection{Saúde}

\textbf{Fonte:} MSPBS /\allowbreak  IPS Paraguay (2024-2026), consolidação
departamental e proxy local

{\def\LTcaptype{none} % do not increment counter
\begin{longtable}[]{@{}ll@{}}
\toprule\noalign{}
Serviço & Disponibilidade \\
\midrule\noalign{}
\endhead
\bottomrule\noalign{}
\endlastfoot
USF /\allowbreak  Posto de Saúde & sim \\
Hospital Regional & não \\
IPS (seguro social) & não \\
Farmácia & sim \\
Distância ao hospital de referência & \textasciitilde14 km (Villa
Hayes) \\
\end{longtable}
}

\textbf{Principais estabelecimentos:} USF local; referencia hospitalar
em Villa Hayes

\textbf{Observação para imigrantes:} Atendimento primario local ou em
raio curto; casos de maior complexidade seguem para Villa Hayes.
Cobertura privada continua recomendada para especialidades e urgências
de maior complexidade.
\newpage
\secaoDiagnostico{Campo Aceval}{\mapaDistrito{1510}}{Campo Aceval (Presidente Hayes) apresenta perfil operacional consolidado
para comparacao territorial, com leitura conjunta de infraestrutura,
risco hidroclimatico e capacidade de continuidade de servicos
essenciais.

\subsubsection{Por que sim}

\subsubsection{Por que não}

\subsubsection{Pre-condicoes Minimas de
Decisao}

\begin{itemize}
\tightlist
\item
  Atualizar indicadores criticos em ciclo trimestral.
\item
  Confirmar trafegabilidade e contingencia hidrica/\allowbreak energetica local.
\item
  Validar checklist de resiliencia domiciliar/\allowbreak logistica antes de decisao
  final.
\end{itemize}

\subsubsection{Recomendação
Técnica}

Uso recomendado como candidato em analise multicriterio, condicionado ao
cumprimento das pre-condicoes acima. Referência atual de score: GSS 7.7;
risco predominante: alto.

\begin{itemize}
\tightlist
\item
  \begin{enumerate}
  \def\labelenumi{\arabic{enumi})}
  \tightlist
  \item
    Delimitacao territorial validada por georreferencia institucional.
  \end{enumerate}
\item
  \begin{enumerate}
  \def\labelenumi{\arabic{enumi})}
  \setcounter{enumi}{1}
  \tightlist
  \item
    População municipal/\allowbreak distrital referenciada por base censitaria
    nacional.
  \end{enumerate}
\item
  \begin{enumerate}
  \def\labelenumi{\arabic{enumi})}
  \setcounter{enumi}{2}
  \tightlist
  \item
    Comparabilidade intra-departamental apoiada em indicadores
    distritais oficiais.
  \end{enumerate}
\item
  \begin{enumerate}
  \def\labelenumi{\arabic{enumi})}
  \setcounter{enumi}{3}
  \tightlist
  \item
    Estado macro de conectividade viaria ancorado em informacoes de
    rodovias.
  \end{enumerate}
\item
  \begin{enumerate}
  \def\labelenumi{\arabic{enumi})}
  \setcounter{enumi}{4}
  \tightlist
  \item
    Exposicao hidrometeorologica monitorada por avisos oficiais.
  \end{enumerate}
\item
  \begin{enumerate}
  \def\labelenumi{\arabic{enumi})}
  \setcounter{enumi}{5}
  \tightlist
  \item
    Historico de contexto meteorologico apoiado em publicacoes tecnicas.
  \end{enumerate}
\item
  \begin{enumerate}
  \def\labelenumi{\arabic{enumi})}
  \setcounter{enumi}{6}
  \tightlist
  \item
    Capacidade institucional de resposta a eventos apoiada em acoes da
    SEN\footnote{\fineSEN}.
  \end{enumerate}
\item
  \begin{enumerate}
  \def\labelenumi{\arabic{enumi})}
  \setcounter{enumi}{7}
  \tightlist
  \item
    Infraestrutura energetica analisada via operador nacional.
  \end{enumerate}
\item
  \begin{enumerate}
  \def\labelenumi{\arabic{enumi})}
  \setcounter{enumi}{8}
  \tightlist
  \item
    Referências de obras e logistica nacional verificadas no MOPC.
  \end{enumerate}
\item
  \begin{enumerate}
  \def\labelenumi{\arabic{enumi})}
  \setcounter{enumi}{9}
  \tightlist
  \item
    Contexto sociopolitico institucional referenciado por autoridade
    eleitoral.
  \end{enumerate}
\item
  \begin{enumerate}
  \def\labelenumi{\arabic{enumi})}
  \setcounter{enumi}{10}
  \tightlist
  \item
    Camada cartografica de apoio eleitoral/\allowbreak territorial consultada.
  \end{enumerate}
\item
  \begin{enumerate}
  \def\labelenumi{\arabic{enumi})}
  \setcounter{enumi}{11}
  \tightlist
  \item
    Dados publicos complementares para verificacao cruzada.
  \end{enumerate}
\end{itemize}}{dist:15_Presidente_Hayes:Campo_Aceval}
\begin{table}[ht!]
\small \begin{tabular}{p{3.0cm}p{0.8cm}p{6.0cm}} \toprule
\textbf{Critério} & \textbf{Nota} & \textbf{Justificativa} \\ \midrule
A - Ameaças Estratégicas & 7.8 & - \\
B - Risco Social & 8.8 & - \\
C - Riscos Naturais & 4.6 & - \\
D - Autossuficiência & 8.5 & - \\
E - Institucional & 7.9 & - \\
\midrule \textbf{GSS FINAL} & \textbf{7.7} & \textbf{Seguro (bom potencial com preparacao).} \\ \bottomrule \end{tabular}
\end{table}
\subsubsection{Vulnerabilidades e
Mitigação}

\begin{itemize}
\tightlist
\item
  Alta exposicao hidrometeorologica (\textgreater=65/\allowbreak 100).
\item
  Conectividade viaria favoravel (\textgreater=70/\allowbreak 100).
\item
  Estoque de contingencia para 14 dias (agua/\allowbreak alimentos/\allowbreak combustivel),
  priorizado quando risco hidro=67/\allowbreak 100.
\item
  Duplicar rotas/\allowbreak logistica quando conectividade viaria=93/\allowbreak 100 for
  \textless70; manter rota secundaria mapeada e testada trimestralmente.
\item
  Plano de energia resiliente com autonomia minima de 72h
  (geracao/\allowbreak backup), revisao semestral.
\end{itemize}

\subsubsection{Dossiê de Campo}
\begin{description}[style=nextline,leftmargin=0.5cm,font=\bfseries]
\item[Coordenadas]  23°12'S 59°46'W (geocodificado via Nominatim).
\end{description}
\textbf{Fonte climática:} NASA POWER Climatology API (período 2001-2020)

\textbf{Fonte luminosa:} estimativa world atlas

\paragraph{Irradiação Solar
(kWh/\allowbreak m²/\allowbreak dia)}

{\def\LTcaptype{none} % do not increment counter
\begin{longtable}[]{@{}
  >{\raggedright\arraybackslash}p{(\linewidth - 24\tabcolsep) * \real{0.0746}}
  >{\raggedright\arraybackslash}p{(\linewidth - 24\tabcolsep) * \real{0.0746}}
  >{\raggedright\arraybackslash}p{(\linewidth - 24\tabcolsep) * \real{0.0746}}
  >{\raggedright\arraybackslash}p{(\linewidth - 24\tabcolsep) * \real{0.0746}}
  >{\raggedright\arraybackslash}p{(\linewidth - 24\tabcolsep) * \real{0.0746}}
  >{\raggedright\arraybackslash}p{(\linewidth - 24\tabcolsep) * \real{0.0746}}
  >{\raggedright\arraybackslash}p{(\linewidth - 24\tabcolsep) * \real{0.0746}}
  >{\raggedright\arraybackslash}p{(\linewidth - 24\tabcolsep) * \real{0.0746}}
  >{\raggedright\arraybackslash}p{(\linewidth - 24\tabcolsep) * \real{0.0746}}
  >{\raggedright\arraybackslash}p{(\linewidth - 24\tabcolsep) * \real{0.0746}}
  >{\raggedright\arraybackslash}p{(\linewidth - 24\tabcolsep) * \real{0.0746}}
  >{\raggedright\arraybackslash}p{(\linewidth - 24\tabcolsep) * \real{0.0746}}
  >{\raggedright\arraybackslash}p{(\linewidth - 24\tabcolsep) * \real{0.1045}}@{}}
\toprule\noalign{}
\begin{minipage}[b]{\linewidth}\raggedright
Jan
\end{minipage} & \begin{minipage}[b]{\linewidth}\raggedright
Fev
\end{minipage} & \begin{minipage}[b]{\linewidth}\raggedright
Mar
\end{minipage} & \begin{minipage}[b]{\linewidth}\raggedright
Abr
\end{minipage} & \begin{minipage}[b]{\linewidth}\raggedright
Mai
\end{minipage} & \begin{minipage}[b]{\linewidth}\raggedright
Jun
\end{minipage} & \begin{minipage}[b]{\linewidth}\raggedright
Jul
\end{minipage} & \begin{minipage}[b]{\linewidth}\raggedright
Ago
\end{minipage} & \begin{minipage}[b]{\linewidth}\raggedright
Set
\end{minipage} & \begin{minipage}[b]{\linewidth}\raggedright
Out
\end{minipage} & \begin{minipage}[b]{\linewidth}\raggedright
Nov
\end{minipage} & \begin{minipage}[b]{\linewidth}\raggedright
Dez
\end{minipage} & \begin{minipage}[b]{\linewidth}\raggedright
Média
\end{minipage} \\
\midrule\noalign{}
\endhead
\bottomrule\noalign{}
\endlastfoot
6.68 & 6.08 & 5.46 & 4.44 & 3.41 & 2.94 & 3.33 & 4.15 & 4.84 & 5.56 &
6.36 & 6.59 & \textbf{4.99} \\
\end{longtable}
}

\textbf{Inclinação solar recomendada:} 23° N (anual) · 33° N (inverno
jun-ago) · 13° N (verão nov-jan)

\paragraph{Precipitação (mm/\allowbreak mês)}

{\def\LTcaptype{none} % do not increment counter
\begin{longtable}[]{@{}
  >{\raggedright\arraybackslash}p{(\linewidth - 24\tabcolsep) * \real{0.0704}}
  >{\raggedright\arraybackslash}p{(\linewidth - 24\tabcolsep) * \real{0.0704}}
  >{\raggedright\arraybackslash}p{(\linewidth - 24\tabcolsep) * \real{0.0704}}
  >{\raggedright\arraybackslash}p{(\linewidth - 24\tabcolsep) * \real{0.0704}}
  >{\raggedright\arraybackslash}p{(\linewidth - 24\tabcolsep) * \real{0.0704}}
  >{\raggedright\arraybackslash}p{(\linewidth - 24\tabcolsep) * \real{0.0704}}
  >{\raggedright\arraybackslash}p{(\linewidth - 24\tabcolsep) * \real{0.0704}}
  >{\raggedright\arraybackslash}p{(\linewidth - 24\tabcolsep) * \real{0.0704}}
  >{\raggedright\arraybackslash}p{(\linewidth - 24\tabcolsep) * \real{0.0704}}
  >{\raggedright\arraybackslash}p{(\linewidth - 24\tabcolsep) * \real{0.0704}}
  >{\raggedright\arraybackslash}p{(\linewidth - 24\tabcolsep) * \real{0.0704}}
  >{\raggedright\arraybackslash}p{(\linewidth - 24\tabcolsep) * \real{0.0704}}
  >{\raggedright\arraybackslash}p{(\linewidth - 24\tabcolsep) * \real{0.1549}}@{}}
\toprule\noalign{}
\begin{minipage}[b]{\linewidth}\raggedright
Jan
\end{minipage} & \begin{minipage}[b]{\linewidth}\raggedright
Fev
\end{minipage} & \begin{minipage}[b]{\linewidth}\raggedright
Mar
\end{minipage} & \begin{minipage}[b]{\linewidth}\raggedright
Abr
\end{minipage} & \begin{minipage}[b]{\linewidth}\raggedright
Mai
\end{minipage} & \begin{minipage}[b]{\linewidth}\raggedright
Jun
\end{minipage} & \begin{minipage}[b]{\linewidth}\raggedright
Jul
\end{minipage} & \begin{minipage}[b]{\linewidth}\raggedright
Ago
\end{minipage} & \begin{minipage}[b]{\linewidth}\raggedright
Set
\end{minipage} & \begin{minipage}[b]{\linewidth}\raggedright
Out
\end{minipage} & \begin{minipage}[b]{\linewidth}\raggedright
Nov
\end{minipage} & \begin{minipage}[b]{\linewidth}\raggedright
Dez
\end{minipage} & \begin{minipage}[b]{\linewidth}\raggedright
Total/\allowbreak ano
\end{minipage} \\
\midrule\noalign{}
\endhead
\bottomrule\noalign{}
\endlastfoot
104 & 122 & 110 & 79 & 58 & 21 & 16 & 11 & 19 & 69 & 97 & 114 &
\textbf{825 mm} \\
\end{longtable}
}

\paragraph{Poluição Luminosa}

{\def\LTcaptype{none} % do not increment counter
\begin{longtable}[]{@{}ll@{}}
\toprule\noalign{}
Parâmetro & Valor \\
\midrule\noalign{}
\endhead
\bottomrule\noalign{}
\endlastfoot
Escala Bortle & 2 --- Céu tipicamente escuro \\
Radiância artificial & 0.3 nW/\allowbreak cm²/\allowbreak sr \\
\end{longtable}
}
\paragraph{Solo (SoilGrids 2.0, média ponderada 0--30
cm)}

{\def\LTcaptype{none} % do not increment counter
\begin{longtable}[]{@{}ll@{}}
\toprule\noalign{}
Parâmetro & Valor \\
\midrule\noalign{}
\endhead
\bottomrule\noalign{}
\endlastfoot
pH (H$_2$O) & 6.85 \\
Carbono orgânico (SOC) & 15.27 g/\allowbreak kg \\
Argila & 33.08 \% \\
Areia & 30.82 \% \\
Silte (calc.) & 36.1 \% \\
Densidade aparente & 1.4 g/\allowbreak cm³ \\
\textbf{Aptidão agrícola} & \textbf{Alta} \\
\end{longtable}
}

\subsubsection{Indicadores Sociais}

\textbf{Fontes:} PNUD 2020, INE\footnote{\fineINE} Censo 2022, INE\footnote{\fineINE} EPHC 2023, Ministerio
Público 2024\\
\textbf{Nota:} Valores marcados como \emph{(dept.)} referem-se ao
departamento de Presidente Hayes; valores \emph{(dist.)} são específicos
deste distrito. Dados marcados \emph{(est.)} são estimativas.

{\def\LTcaptype{none} % do not increment counter
\begin{longtable}[]{@{}
  >{\raggedright\arraybackslash}p{(\linewidth - 4\tabcolsep) * \real{0.4231}}
  >{\raggedright\arraybackslash}p{(\linewidth - 4\tabcolsep) * \real{0.2692}}
  >{\raggedright\arraybackslash}p{(\linewidth - 4\tabcolsep) * \real{0.3077}}@{}}
\toprule\noalign{}
\begin{minipage}[b]{\linewidth}\raggedright
Indicador
\end{minipage} & \begin{minipage}[b]{\linewidth}\raggedright
Valor
\end{minipage} & \begin{minipage}[b]{\linewidth}\raggedright
Âmbito
\end{minipage} \\
\midrule\noalign{}
\endhead
\bottomrule\noalign{}
\endlastfoot
IDH (2020) & 0,693 & dept. \\
Ranking IDH nacional & 13/\allowbreak 18 & dept. \\
Esperança de vida & 71,8 anos & dept. \\
Escolaridade média & 7.6 anos & dist. \\
RNB per capita (USD PPA) & 7.900 & dept. \\
Pobreza monetária (\%) & 29,4\% & dept. \\
Pobreza extrema (\%) & 7,2\% & dept. \\
Índice de Gini & 0,453 & dept. \\
Acesso a água potável (\%) & 71,3\% & dept. \\
Acesso a saneamento (\%) & 68,9\% & dept. \\
Taxa de homicídios (est., /\allowbreak 100k hab) & 5,5--5,6 (2019--2020, fonte INE
--- próxima à média nacional) & dept. \\
Índice de segurança & médio & dept. \\
População (Censo 2022) & N/\allowbreak D & dist. \\
Idade mediana & N/\allowbreak D & dist. \\
Taxa de fecundidade & N/\allowbreak D & dist. \\
\end{longtable}
}
\subsubsection{Combustível}

\textbf{Referência:} PETROPAR /\allowbreak  postos locais (2024) \textbf{Tipo de
localidade:} Interior

{\def\LTcaptype{none} % do not increment counter
\begin{longtable}[]{@{}lll@{}}
\toprule\noalign{}
Combustível & USD/\allowbreak litro & Gs/\allowbreak litro (aprox.) \\
\midrule\noalign{}
\endhead
\bottomrule\noalign{}
\endlastfoot
Gasolina 93 oct & 1.05 & 7,770 \\
Gasolina 97 oct (premium) & 1.15 & 8,510 \\
Diesel & 0.98 & 7,252 \\
\end{longtable}
}

\begin{quote}
Preços podem variar ±5\% conforme posto e sazonalidade. Chaco e interior
remoto apresentam maior variação.
\end{quote}

\subsubsection{Cobertura Celular}

\textbf{Fonte:} CONATEL PY /\allowbreak  operadoras (2024)

{\def\LTcaptype{none} % do not increment counter
\begin{longtable}[]{@{}ll@{}}
\toprule\noalign{}
Parâmetro & Valor \\
\midrule\noalign{}
\endhead
\bottomrule\noalign{}
\endlastfoot
Cobertura 4G população (dept.) & 60\% \\
Cobertura 4G área rural & 35\% \\
Melhor operadora & Tigo \\
Qualidade rural & limitada \\
\end{longtable}
}

\begin{quote}
Para áreas rurais fora do núcleo urbano, recomenda-se chip Tigo como
principal e Personal como backup.
\end{quote}

\subsubsection{Internet}

\textbf{Fonte:} CONATEL /\allowbreak  Speedtest Ookla (2024)

{\def\LTcaptype{none} % do not increment counter
\begin{longtable}[]{@{}ll@{}}
\toprule\noalign{}
Parâmetro & Valor \\
\midrule\noalign{}
\endhead
\bottomrule\noalign{}
\endlastfoot
Velocidade média download & 25 Mbps \\
Domicílios com internet (dept.) & 35\% \\
Tecnologia predominante & rádio/\allowbreak satélite \\
Opção rural & Starlink disponível (\textasciitilde USD 44/\allowbreak mês) \\
\end{longtable}
}

\subsubsection{Mercado Imobiliário e Terra
Rural}

\textbf{Fonte:} INDERT /\allowbreak  Clasificados.com.py (2024)

{\def\LTcaptype{none} % do not increment counter
\begin{longtable}[]{@{}ll@{}}
\toprule\noalign{}
Tipo & Referência \\
\midrule\noalign{}
\endhead
\bottomrule\noalign{}
\endlastfoot
Terra agrícola alta prod. (USD/\allowbreak ha) & 2,000 \\
Imóvel urbano (USD/\allowbreak m²) & 600 \\
Aluguel 2 quartos (USD/\allowbreak mês) & 250 \\
\end{longtable}
}

\begin{quote}
Valores de referência departamental. Localidades menores podem ter
preços 20--40\% abaixo da capital departamental.
\end{quote}

\subsubsection{Saúde}

\textbf{Fonte:} MSPBS /\allowbreak  IPS Paraguay (2024-2026), consolidação
departamental e proxy local

{\def\LTcaptype{none} % do not increment counter
\begin{longtable}[]{@{}ll@{}}
\toprule\noalign{}
Serviço & Disponibilidade \\
\midrule\noalign{}
\endhead
\bottomrule\noalign{}
\endlastfoot
USF /\allowbreak  Posto de Saúde & sim \\
Hospital Regional & não \\
IPS (seguro social) & não \\
Farmácia & sim \\
Distância ao hospital de referência & \textasciitilde387 km (Villa
Hayes) \\
\end{longtable}
}

\textbf{Principais estabelecimentos:} USF local; referencia hospitalar
em Villa Hayes

\textbf{Observação para imigrantes:} Atendimento primario local ou em
raio curto; casos de maior complexidade seguem para Villa Hayes.
Cobertura privada continua recomendada para especialidades e urgências
de maior complexidade.
\newpage
\secaoDiagnostico{General Bruguez}{\mapaConteudo{-25.1000}{-57.5333}}{General Bruguez (Presidente Hayes) apresenta perfil operacional
consolidado para comparacao territorial, com leitura conjunta de
infraestrutura, risco hidroclimatico e capacidade de continuidade de
servicos essenciais.

\subsubsection{Por que sim}

\subsubsection{Por que não}

\subsubsection{Pre-condicoes Minimas de
Decisao}

\begin{itemize}
\tightlist
\item
  Atualizar indicadores criticos em ciclo trimestral.
\item
  Confirmar trafegabilidade e contingencia hidrica/\allowbreak energetica local.
\item
  Validar checklist de resiliencia domiciliar/\allowbreak logistica antes de decisao
  final.
\end{itemize}

\subsubsection{Recomendação
Técnica}

Uso recomendado como candidato em analise multicriterio, condicionado ao
cumprimento das pre-condicoes acima. Referência atual de score: GSS 7.2;
risco predominante: alto.

\begin{itemize}
\tightlist
\item
  \begin{enumerate}
  \def\labelenumi{\arabic{enumi})}
  \tightlist
  \item
    Delimitacao territorial validada por georreferencia institucional.
  \end{enumerate}
\item
  \begin{enumerate}
  \def\labelenumi{\arabic{enumi})}
  \setcounter{enumi}{1}
  \tightlist
  \item
    População municipal/\allowbreak distrital referenciada por base censitaria
    nacional.
  \end{enumerate}
\item
  \begin{enumerate}
  \def\labelenumi{\arabic{enumi})}
  \setcounter{enumi}{2}
  \tightlist
  \item
    Comparabilidade intra-departamental apoiada em indicadores
    distritais oficiais.
  \end{enumerate}
\item
  \begin{enumerate}
  \def\labelenumi{\arabic{enumi})}
  \setcounter{enumi}{3}
  \tightlist
  \item
    Estado macro de conectividade viaria ancorado em informacoes de
    rodovias.
  \end{enumerate}
\item
  \begin{enumerate}
  \def\labelenumi{\arabic{enumi})}
  \setcounter{enumi}{4}
  \tightlist
  \item
    Exposicao hidrometeorologica monitorada por avisos oficiais.
  \end{enumerate}
\item
  \begin{enumerate}
  \def\labelenumi{\arabic{enumi})}
  \setcounter{enumi}{5}
  \tightlist
  \item
    Historico de contexto meteorologico apoiado em publicacoes tecnicas.
  \end{enumerate}
\item
  \begin{enumerate}
  \def\labelenumi{\arabic{enumi})}
  \setcounter{enumi}{6}
  \tightlist
  \item
    Capacidade institucional de resposta a eventos apoiada em acoes da
    SEN\footnote{\fineSEN}.
  \end{enumerate}
\item
  \begin{enumerate}
  \def\labelenumi{\arabic{enumi})}
  \setcounter{enumi}{7}
  \tightlist
  \item
    Infraestrutura energetica analisada via operador nacional.
  \end{enumerate}
\item
  \begin{enumerate}
  \def\labelenumi{\arabic{enumi})}
  \setcounter{enumi}{8}
  \tightlist
  \item
    Referências de obras e logistica nacional verificadas no MOPC.
  \end{enumerate}
\item
  \begin{enumerate}
  \def\labelenumi{\arabic{enumi})}
  \setcounter{enumi}{9}
  \tightlist
  \item
    Contexto sociopolitico institucional referenciado por autoridade
    eleitoral.
  \end{enumerate}
\item
  \begin{enumerate}
  \def\labelenumi{\arabic{enumi})}
  \setcounter{enumi}{10}
  \tightlist
  \item
    Camada cartografica de apoio eleitoral/\allowbreak territorial consultada.
  \end{enumerate}
\item
  \begin{enumerate}
  \def\labelenumi{\arabic{enumi})}
  \setcounter{enumi}{11}
  \tightlist
  \item
    Dados publicos complementares para verificacao cruzada.
  \end{enumerate}
\end{itemize}}{dist:15_Presidente_Hayes:General_Bruguez}
\begin{table}[ht!]
\small \begin{tabular}{p{3.0cm}p{0.8cm}p{6.0cm}} \toprule
\textbf{Critério} & \textbf{Nota} & \textbf{Justificativa} \\ \midrule
A - Ameaças Estratégicas & 7.4 & - \\
B - Risco Social & 8.3 & - \\
C - Riscos Naturais & 4.8 & - \\
D - Autossuficiência & 7.3 & - \\
E - Institucional & 7.7 & - \\
\midrule \textbf{GSS FINAL} & \textbf{7.2} & \textbf{Seguro (bom potencial com preparacao).} \\ \bottomrule \end{tabular}
\end{table}
\subsubsection{Vulnerabilidades e
Mitigação}

\begin{itemize}
\tightlist
\item
  Alta exposicao hidrometeorologica (\textgreater=65/\allowbreak 100).
\item
  Conectividade viaria favoravel (\textgreater=70/\allowbreak 100).
\item
  Estoque de contingencia para 14 dias (agua/\allowbreak alimentos/\allowbreak combustivel),
  priorizado quando risco hidro=65/\allowbreak 100.
\item
  Duplicar rotas/\allowbreak logistica quando conectividade viaria=78/\allowbreak 100 for
  \textless70; manter rota secundaria mapeada e testada trimestralmente.
\item
  Plano de energia resiliente com autonomia minima de 72h
  (geracao/\allowbreak backup), revisao semestral.
\end{itemize}

\subsubsection{Dossiê de Campo}
\begin{description}[style=nextline,leftmargin=0.5cm,font=\bfseries]
\item[Coordenadas]  25°06'S 57°32'W (geocodificado via Nominatim).
\end{description}
\textbf{Fonte climática:} NASA POWER Climatology API (período 2001-2020)

\textbf{Fonte luminosa:} estimativa world atlas

\paragraph{Irradiação Solar
(kWh/\allowbreak m²/\allowbreak dia)}

{\def\LTcaptype{none} % do not increment counter
\begin{longtable}[]{@{}
  >{\raggedright\arraybackslash}p{(\linewidth - 24\tabcolsep) * \real{0.0746}}
  >{\raggedright\arraybackslash}p{(\linewidth - 24\tabcolsep) * \real{0.0746}}
  >{\raggedright\arraybackslash}p{(\linewidth - 24\tabcolsep) * \real{0.0746}}
  >{\raggedright\arraybackslash}p{(\linewidth - 24\tabcolsep) * \real{0.0746}}
  >{\raggedright\arraybackslash}p{(\linewidth - 24\tabcolsep) * \real{0.0746}}
  >{\raggedright\arraybackslash}p{(\linewidth - 24\tabcolsep) * \real{0.0746}}
  >{\raggedright\arraybackslash}p{(\linewidth - 24\tabcolsep) * \real{0.0746}}
  >{\raggedright\arraybackslash}p{(\linewidth - 24\tabcolsep) * \real{0.0746}}
  >{\raggedright\arraybackslash}p{(\linewidth - 24\tabcolsep) * \real{0.0746}}
  >{\raggedright\arraybackslash}p{(\linewidth - 24\tabcolsep) * \real{0.0746}}
  >{\raggedright\arraybackslash}p{(\linewidth - 24\tabcolsep) * \real{0.0746}}
  >{\raggedright\arraybackslash}p{(\linewidth - 24\tabcolsep) * \real{0.0746}}
  >{\raggedright\arraybackslash}p{(\linewidth - 24\tabcolsep) * \real{0.1045}}@{}}
\toprule\noalign{}
\begin{minipage}[b]{\linewidth}\raggedright
Jan
\end{minipage} & \begin{minipage}[b]{\linewidth}\raggedright
Fev
\end{minipage} & \begin{minipage}[b]{\linewidth}\raggedright
Mar
\end{minipage} & \begin{minipage}[b]{\linewidth}\raggedright
Abr
\end{minipage} & \begin{minipage}[b]{\linewidth}\raggedright
Mai
\end{minipage} & \begin{minipage}[b]{\linewidth}\raggedright
Jun
\end{minipage} & \begin{minipage}[b]{\linewidth}\raggedright
Jul
\end{minipage} & \begin{minipage}[b]{\linewidth}\raggedright
Ago
\end{minipage} & \begin{minipage}[b]{\linewidth}\raggedright
Set
\end{minipage} & \begin{minipage}[b]{\linewidth}\raggedright
Out
\end{minipage} & \begin{minipage}[b]{\linewidth}\raggedright
Nov
\end{minipage} & \begin{minipage}[b]{\linewidth}\raggedright
Dez
\end{minipage} & \begin{minipage}[b]{\linewidth}\raggedright
Média
\end{minipage} \\
\midrule\noalign{}
\endhead
\bottomrule\noalign{}
\endlastfoot
6.76 & 6.20 & 5.42 & 4.44 & 3.36 & 2.83 & 3.21 & 3.93 & 4.61 & 5.49 &
6.42 & 6.78 & \textbf{4.95} \\
\end{longtable}
}

\textbf{Inclinação solar recomendada:} 25° N (anual) · 35° N (inverno
jun-ago) · 15° N (verão nov-jan)

\paragraph{Precipitação (mm/\allowbreak mês)}

{\def\LTcaptype{none} % do not increment counter
\begin{longtable}[]{@{}
  >{\raggedright\arraybackslash}p{(\linewidth - 24\tabcolsep) * \real{0.0704}}
  >{\raggedright\arraybackslash}p{(\linewidth - 24\tabcolsep) * \real{0.0704}}
  >{\raggedright\arraybackslash}p{(\linewidth - 24\tabcolsep) * \real{0.0704}}
  >{\raggedright\arraybackslash}p{(\linewidth - 24\tabcolsep) * \real{0.0704}}
  >{\raggedright\arraybackslash}p{(\linewidth - 24\tabcolsep) * \real{0.0704}}
  >{\raggedright\arraybackslash}p{(\linewidth - 24\tabcolsep) * \real{0.0704}}
  >{\raggedright\arraybackslash}p{(\linewidth - 24\tabcolsep) * \real{0.0704}}
  >{\raggedright\arraybackslash}p{(\linewidth - 24\tabcolsep) * \real{0.0704}}
  >{\raggedright\arraybackslash}p{(\linewidth - 24\tabcolsep) * \real{0.0704}}
  >{\raggedright\arraybackslash}p{(\linewidth - 24\tabcolsep) * \real{0.0704}}
  >{\raggedright\arraybackslash}p{(\linewidth - 24\tabcolsep) * \real{0.0704}}
  >{\raggedright\arraybackslash}p{(\linewidth - 24\tabcolsep) * \real{0.0704}}
  >{\raggedright\arraybackslash}p{(\linewidth - 24\tabcolsep) * \real{0.1549}}@{}}
\toprule\noalign{}
\begin{minipage}[b]{\linewidth}\raggedright
Jan
\end{minipage} & \begin{minipage}[b]{\linewidth}\raggedright
Fev
\end{minipage} & \begin{minipage}[b]{\linewidth}\raggedright
Mar
\end{minipage} & \begin{minipage}[b]{\linewidth}\raggedright
Abr
\end{minipage} & \begin{minipage}[b]{\linewidth}\raggedright
Mai
\end{minipage} & \begin{minipage}[b]{\linewidth}\raggedright
Jun
\end{minipage} & \begin{minipage}[b]{\linewidth}\raggedright
Jul
\end{minipage} & \begin{minipage}[b]{\linewidth}\raggedright
Ago
\end{minipage} & \begin{minipage}[b]{\linewidth}\raggedright
Set
\end{minipage} & \begin{minipage}[b]{\linewidth}\raggedright
Out
\end{minipage} & \begin{minipage}[b]{\linewidth}\raggedright
Nov
\end{minipage} & \begin{minipage}[b]{\linewidth}\raggedright
Dez
\end{minipage} & \begin{minipage}[b]{\linewidth}\raggedright
Total/\allowbreak ano
\end{minipage} \\
\midrule\noalign{}
\endhead
\bottomrule\noalign{}
\endlastfoot
123 & 149 & 130 & 143 & 143 & 71 & 59 & 33 & 70 & 161 & 197 & 182 &
\textbf{1465 mm} \\
\end{longtable}
}

\paragraph{Poluição Luminosa}

{\def\LTcaptype{none} % do not increment counter
\begin{longtable}[]{@{}ll@{}}
\toprule\noalign{}
Parâmetro & Valor \\
\midrule\noalign{}
\endhead
\bottomrule\noalign{}
\endlastfoot
Escala Bortle & 2 --- Céu tipicamente escuro \\
Radiância artificial & 0.3 nW/\allowbreak cm²/\allowbreak sr \\
\end{longtable}
}
\paragraph{Solo (SoilGrids 2.0, média ponderada 0--30
cm)}

{\def\LTcaptype{none} % do not increment counter
\begin{longtable}[]{@{}ll@{}}
\toprule\noalign{}
Parâmetro & Valor \\
\midrule\noalign{}
\endhead
\bottomrule\noalign{}
\endlastfoot
pH (H$_2$O) & 5.98 \\
Carbono orgânico (SOC) & 14.32 g/\allowbreak kg \\
Argila & 24.73 \% \\
Areia & 43.82 \% \\
Silte (calc.) & 31.4 \% \\
Densidade aparente & 1.37 g/\allowbreak cm³ \\
\textbf{Aptidão agrícola} & \textbf{Alta} \\
\end{longtable}
}

\subsubsection{Indicadores Sociais}

\textbf{Fontes:} PNUD 2020, INE\footnote{\fineINE} Censo 2022, INE\footnote{\fineINE} EPHC 2023, Ministerio
Público 2024\\
\textbf{Nota:} Valores marcados como \emph{(dept.)} referem-se ao
departamento de Presidente Hayes; valores \emph{(dist.)} são específicos
deste distrito. Dados marcados \emph{(est.)} são estimativas.

{\def\LTcaptype{none} % do not increment counter
\begin{longtable}[]{@{}
  >{\raggedright\arraybackslash}p{(\linewidth - 4\tabcolsep) * \real{0.4231}}
  >{\raggedright\arraybackslash}p{(\linewidth - 4\tabcolsep) * \real{0.2692}}
  >{\raggedright\arraybackslash}p{(\linewidth - 4\tabcolsep) * \real{0.3077}}@{}}
\toprule\noalign{}
\begin{minipage}[b]{\linewidth}\raggedright
Indicador
\end{minipage} & \begin{minipage}[b]{\linewidth}\raggedright
Valor
\end{minipage} & \begin{minipage}[b]{\linewidth}\raggedright
Âmbito
\end{minipage} \\
\midrule\noalign{}
\endhead
\bottomrule\noalign{}
\endlastfoot
IDH (2020) & 0,693 & dept. \\
Ranking IDH nacional & 13/\allowbreak 18 & dept. \\
Esperança de vida & 71,8 anos & dept. \\
Escolaridade média & 7,8 anos & dept. \\
RNB per capita (USD PPA) & 7.900 & dept. \\
Pobreza monetária (\%) & 29,4\% & dept. \\
Pobreza extrema (\%) & 7,2\% & dept. \\
Índice de Gini & 0,453 & dept. \\
Acesso a água potável (\%) & 71,3\% & dept. \\
Acesso a saneamento (\%) & 68,9\% & dept. \\
Taxa de homicídios (est., /\allowbreak 100k hab) & 5,5--5,6 (2019--2020, fonte INE
--- próxima à média nacional) & dept. \\
Índice de segurança & médio & dept. \\
População (Censo 2022) & N/\allowbreak D & dist. \\
Idade mediana & N/\allowbreak D & dist. \\
Taxa de fecundidade & N/\allowbreak D & dist. \\
\end{longtable}
}
\subsubsection{Combustível}

\textbf{Referência:} PETROPAR /\allowbreak  postos locais (2024) \textbf{Tipo de
localidade:} Interior

{\def\LTcaptype{none} % do not increment counter
\begin{longtable}[]{@{}lll@{}}
\toprule\noalign{}
Combustível & USD/\allowbreak litro & Gs/\allowbreak litro (aprox.) \\
\midrule\noalign{}
\endhead
\bottomrule\noalign{}
\endlastfoot
Gasolina 93 oct & 1.05 & 7,770 \\
Gasolina 97 oct (premium) & 1.15 & 8,510 \\
Diesel & 0.98 & 7,252 \\
\end{longtable}
}

\begin{quote}
Preços podem variar ±5\% conforme posto e sazonalidade. Chaco e interior
remoto apresentam maior variação.
\end{quote}

\subsubsection{Cobertura Celular}

\textbf{Fonte:} CONATEL PY /\allowbreak  operadoras (2024)

{\def\LTcaptype{none} % do not increment counter
\begin{longtable}[]{@{}ll@{}}
\toprule\noalign{}
Parâmetro & Valor \\
\midrule\noalign{}
\endhead
\bottomrule\noalign{}
\endlastfoot
Cobertura 4G população (dept.) & 60\% \\
Cobertura 4G área rural & 35\% \\
Melhor operadora & Tigo \\
Qualidade rural & limitada \\
\end{longtable}
}

\begin{quote}
Para áreas rurais fora do núcleo urbano, recomenda-se chip Tigo como
principal e Personal como backup.
\end{quote}

\subsubsection{Internet}

\textbf{Fonte:} CONATEL /\allowbreak  Speedtest Ookla (2024)

{\def\LTcaptype{none} % do not increment counter
\begin{longtable}[]{@{}ll@{}}
\toprule\noalign{}
Parâmetro & Valor \\
\midrule\noalign{}
\endhead
\bottomrule\noalign{}
\endlastfoot
Velocidade média download & 25 Mbps \\
Domicílios com internet (dept.) & 35\% \\
Tecnologia predominante & rádio/\allowbreak satélite \\
Opção rural & Starlink disponível (\textasciitilde USD 44/\allowbreak mês) \\
\end{longtable}
}

\subsubsection{Mercado Imobiliário e Terra
Rural}

\textbf{Fonte:} INDERT /\allowbreak  Clasificados.com.py (2024)

{\def\LTcaptype{none} % do not increment counter
\begin{longtable}[]{@{}ll@{}}
\toprule\noalign{}
Tipo & Referência \\
\midrule\noalign{}
\endhead
\bottomrule\noalign{}
\endlastfoot
Terra agrícola alta prod. (USD/\allowbreak ha) & 2,000 \\
Imóvel urbano (USD/\allowbreak m²) & 600 \\
Aluguel 2 quartos (USD/\allowbreak mês) & 250 \\
\end{longtable}
}

\begin{quote}
Valores de referência departamental. Localidades menores podem ter
preços 20--40\% abaixo da capital departamental.
\end{quote}

\subsubsection{Saúde}

\textbf{Fonte:} MSPBS /\allowbreak  IPS Paraguay (2024-2026), consolidação
departamental e proxy local

{\def\LTcaptype{none} % do not increment counter
\begin{longtable}[]{@{}ll@{}}
\toprule\noalign{}
Serviço & Disponibilidade \\
\midrule\noalign{}
\endhead
\bottomrule\noalign{}
\endlastfoot
USF /\allowbreak  Posto de Saúde & sim \\
Hospital Regional & não \\
IPS (seguro social) & não \\
Farmácia & sim \\
Distância ao hospital de referência & \textasciitilde1 km (Villa
Hayes) \\
\end{longtable}
}

\textbf{Principais estabelecimentos:} USF local; referencia hospitalar
em Villa Hayes

\textbf{Observação para imigrantes:} Atendimento primario local ou em
raio curto; casos de maior complexidade seguem para Villa Hayes.
Cobertura privada continua recomendada para especialidades e urgências
de maior complexidade.
\newpage
\secaoDiagnostico{Jose Falcon}{\mapaDistrito{1506}}{Jose Falcon (Presidente Hayes) apresenta perfil operacional consolidado
para comparacao territorial, com leitura conjunta de infraestrutura,
risco hidroclimatico e capacidade de continuidade de servicos
essenciais.

\subsubsection{Por que sim}

\subsubsection{Por que não}

\subsubsection{Pre-condicoes Minimas de
Decisao}

\begin{itemize}
\tightlist
\item
  Atualizar indicadores criticos em ciclo trimestral.
\item
  Confirmar trafegabilidade e contingencia hidrica/\allowbreak energetica local.
\item
  Validar checklist de resiliencia domiciliar/\allowbreak logistica antes de decisao
  final.
\end{itemize}

\subsubsection{Recomendação
Técnica}

Uso recomendado como candidato em analise multicriterio, condicionado ao
cumprimento das pre-condicoes acima. Referência atual de score: GSS 7.2;
risco predominante: moderado.

\begin{itemize}
\tightlist
\item
  \begin{enumerate}
  \def\labelenumi{\arabic{enumi})}
  \tightlist
  \item
    Delimitacao territorial validada por georreferencia institucional.
  \end{enumerate}
\item
  \begin{enumerate}
  \def\labelenumi{\arabic{enumi})}
  \setcounter{enumi}{1}
  \tightlist
  \item
    População municipal/\allowbreak distrital referenciada por base censitaria
    nacional.
  \end{enumerate}
\item
  \begin{enumerate}
  \def\labelenumi{\arabic{enumi})}
  \setcounter{enumi}{2}
  \tightlist
  \item
    Comparabilidade intra-departamental apoiada em indicadores
    distritais oficiais.
  \end{enumerate}
\item
  \begin{enumerate}
  \def\labelenumi{\arabic{enumi})}
  \setcounter{enumi}{3}
  \tightlist
  \item
    Estado macro de conectividade viaria ancorado em informacoes de
    rodovias.
  \end{enumerate}
\item
  \begin{enumerate}
  \def\labelenumi{\arabic{enumi})}
  \setcounter{enumi}{4}
  \tightlist
  \item
    Exposicao hidrometeorologica monitorada por avisos oficiais.
  \end{enumerate}
\item
  \begin{enumerate}
  \def\labelenumi{\arabic{enumi})}
  \setcounter{enumi}{5}
  \tightlist
  \item
    Historico de contexto meteorologico apoiado em publicacoes tecnicas.
  \end{enumerate}
\item
  \begin{enumerate}
  \def\labelenumi{\arabic{enumi})}
  \setcounter{enumi}{6}
  \tightlist
  \item
    Capacidade institucional de resposta a eventos apoiada em acoes da
    SEN\footnote{\fineSEN}.
  \end{enumerate}
\item
  \begin{enumerate}
  \def\labelenumi{\arabic{enumi})}
  \setcounter{enumi}{7}
  \tightlist
  \item
    Infraestrutura energetica analisada via operador nacional.
  \end{enumerate}
\item
  \begin{enumerate}
  \def\labelenumi{\arabic{enumi})}
  \setcounter{enumi}{8}
  \tightlist
  \item
    Referências de obras e logistica nacional verificadas no MOPC.
  \end{enumerate}
\item
  \begin{enumerate}
  \def\labelenumi{\arabic{enumi})}
  \setcounter{enumi}{9}
  \tightlist
  \item
    Contexto sociopolitico institucional referenciado por autoridade
    eleitoral.
  \end{enumerate}
\item
  \begin{enumerate}
  \def\labelenumi{\arabic{enumi})}
  \setcounter{enumi}{10}
  \tightlist
  \item
    Camada cartografica de apoio eleitoral/\allowbreak territorial consultada.
  \end{enumerate}
\item
  \begin{enumerate}
  \def\labelenumi{\arabic{enumi})}
  \setcounter{enumi}{11}
  \tightlist
  \item
    Dados publicos complementares para verificacao cruzada.
  \end{enumerate}
\end{itemize}}{dist:15_Presidente_Hayes:Jose_Falcon}
\begin{table}[ht!]
\small \begin{tabular}{p{3.0cm}p{0.8cm}p{6.0cm}} \toprule
\textbf{Critério} & \textbf{Nota} & \textbf{Justificativa} \\ \midrule
A - Ameaças Estratégicas & 7.9 & - \\
B - Risco Social & 7.3 & - \\
C - Riscos Naturais & 5.0 & - \\
D - Autossuficiência & 8.2 & - \\
E - Institucional & 7.0 & - \\
\midrule \textbf{GSS FINAL} & \textbf{7.2} & \textbf{Seguro (bom potencial com preparacao).} \\ \bottomrule \end{tabular}
\end{table}
\subsubsection{Vulnerabilidades e
Mitigação}

\begin{itemize}
\tightlist
\item
  Exposicao hidrometeorologica moderada (45-64/\allowbreak 100).
\item
  Conectividade viaria favoravel (\textgreater=70/\allowbreak 100).
\item
  Estoque de contingencia para 14 dias (agua/\allowbreak alimentos/\allowbreak combustivel),
  priorizado quando risco hidro=63/\allowbreak 100.
\item
  Duplicar rotas/\allowbreak logistica quando conectividade viaria=94/\allowbreak 100 for
  \textless70; manter rota secundaria mapeada e testada trimestralmente.
\item
  Plano de energia resiliente com autonomia minima de 72h
  (geracao/\allowbreak backup), revisao semestral.
\end{itemize}

\subsubsection{Dossiê de Campo}
\begin{description}[style=nextline,leftmargin=0.5cm,font=\bfseries]
\item[Coordenadas]  25°13'S 57°43'W (geocodificado via Nominatim).
\end{description}
\textbf{Fonte climática:} NASA POWER Climatology API (período 2001-2020)

\textbf{Fonte luminosa:} estimativa world atlas

\paragraph{Irradiação Solar
(kWh/\allowbreak m²/\allowbreak dia)}

{\def\LTcaptype{none} % do not increment counter
\begin{longtable}[]{@{}
  >{\raggedright\arraybackslash}p{(\linewidth - 24\tabcolsep) * \real{0.0746}}
  >{\raggedright\arraybackslash}p{(\linewidth - 24\tabcolsep) * \real{0.0746}}
  >{\raggedright\arraybackslash}p{(\linewidth - 24\tabcolsep) * \real{0.0746}}
  >{\raggedright\arraybackslash}p{(\linewidth - 24\tabcolsep) * \real{0.0746}}
  >{\raggedright\arraybackslash}p{(\linewidth - 24\tabcolsep) * \real{0.0746}}
  >{\raggedright\arraybackslash}p{(\linewidth - 24\tabcolsep) * \real{0.0746}}
  >{\raggedright\arraybackslash}p{(\linewidth - 24\tabcolsep) * \real{0.0746}}
  >{\raggedright\arraybackslash}p{(\linewidth - 24\tabcolsep) * \real{0.0746}}
  >{\raggedright\arraybackslash}p{(\linewidth - 24\tabcolsep) * \real{0.0746}}
  >{\raggedright\arraybackslash}p{(\linewidth - 24\tabcolsep) * \real{0.0746}}
  >{\raggedright\arraybackslash}p{(\linewidth - 24\tabcolsep) * \real{0.0746}}
  >{\raggedright\arraybackslash}p{(\linewidth - 24\tabcolsep) * \real{0.0746}}
  >{\raggedright\arraybackslash}p{(\linewidth - 24\tabcolsep) * \real{0.1045}}@{}}
\toprule\noalign{}
\begin{minipage}[b]{\linewidth}\raggedright
Jan
\end{minipage} & \begin{minipage}[b]{\linewidth}\raggedright
Fev
\end{minipage} & \begin{minipage}[b]{\linewidth}\raggedright
Mar
\end{minipage} & \begin{minipage}[b]{\linewidth}\raggedright
Abr
\end{minipage} & \begin{minipage}[b]{\linewidth}\raggedright
Mai
\end{minipage} & \begin{minipage}[b]{\linewidth}\raggedright
Jun
\end{minipage} & \begin{minipage}[b]{\linewidth}\raggedright
Jul
\end{minipage} & \begin{minipage}[b]{\linewidth}\raggedright
Ago
\end{minipage} & \begin{minipage}[b]{\linewidth}\raggedright
Set
\end{minipage} & \begin{minipage}[b]{\linewidth}\raggedright
Out
\end{minipage} & \begin{minipage}[b]{\linewidth}\raggedright
Nov
\end{minipage} & \begin{minipage}[b]{\linewidth}\raggedright
Dez
\end{minipage} & \begin{minipage}[b]{\linewidth}\raggedright
Média
\end{minipage} \\
\midrule\noalign{}
\endhead
\bottomrule\noalign{}
\endlastfoot
6.76 & 6.20 & 5.42 & 4.44 & 3.36 & 2.83 & 3.21 & 3.93 & 4.61 & 5.49 &
6.42 & 6.78 & \textbf{4.95} \\
\end{longtable}
}

\textbf{Inclinação solar recomendada:} 25° N (anual) · 35° N (inverno
jun-ago) · 15° N (verão nov-jan)

\paragraph{Precipitação (mm/\allowbreak mês)}

{\def\LTcaptype{none} % do not increment counter
\begin{longtable}[]{@{}
  >{\raggedright\arraybackslash}p{(\linewidth - 24\tabcolsep) * \real{0.0704}}
  >{\raggedright\arraybackslash}p{(\linewidth - 24\tabcolsep) * \real{0.0704}}
  >{\raggedright\arraybackslash}p{(\linewidth - 24\tabcolsep) * \real{0.0704}}
  >{\raggedright\arraybackslash}p{(\linewidth - 24\tabcolsep) * \real{0.0704}}
  >{\raggedright\arraybackslash}p{(\linewidth - 24\tabcolsep) * \real{0.0704}}
  >{\raggedright\arraybackslash}p{(\linewidth - 24\tabcolsep) * \real{0.0704}}
  >{\raggedright\arraybackslash}p{(\linewidth - 24\tabcolsep) * \real{0.0704}}
  >{\raggedright\arraybackslash}p{(\linewidth - 24\tabcolsep) * \real{0.0704}}
  >{\raggedright\arraybackslash}p{(\linewidth - 24\tabcolsep) * \real{0.0704}}
  >{\raggedright\arraybackslash}p{(\linewidth - 24\tabcolsep) * \real{0.0704}}
  >{\raggedright\arraybackslash}p{(\linewidth - 24\tabcolsep) * \real{0.0704}}
  >{\raggedright\arraybackslash}p{(\linewidth - 24\tabcolsep) * \real{0.0704}}
  >{\raggedright\arraybackslash}p{(\linewidth - 24\tabcolsep) * \real{0.1549}}@{}}
\toprule\noalign{}
\begin{minipage}[b]{\linewidth}\raggedright
Jan
\end{minipage} & \begin{minipage}[b]{\linewidth}\raggedright
Fev
\end{minipage} & \begin{minipage}[b]{\linewidth}\raggedright
Mar
\end{minipage} & \begin{minipage}[b]{\linewidth}\raggedright
Abr
\end{minipage} & \begin{minipage}[b]{\linewidth}\raggedright
Mai
\end{minipage} & \begin{minipage}[b]{\linewidth}\raggedright
Jun
\end{minipage} & \begin{minipage}[b]{\linewidth}\raggedright
Jul
\end{minipage} & \begin{minipage}[b]{\linewidth}\raggedright
Ago
\end{minipage} & \begin{minipage}[b]{\linewidth}\raggedright
Set
\end{minipage} & \begin{minipage}[b]{\linewidth}\raggedright
Out
\end{minipage} & \begin{minipage}[b]{\linewidth}\raggedright
Nov
\end{minipage} & \begin{minipage}[b]{\linewidth}\raggedright
Dez
\end{minipage} & \begin{minipage}[b]{\linewidth}\raggedright
Total/\allowbreak ano
\end{minipage} \\
\midrule\noalign{}
\endhead
\bottomrule\noalign{}
\endlastfoot
123 & 149 & 130 & 143 & 143 & 71 & 59 & 33 & 70 & 161 & 197 & 182 &
\textbf{1465 mm} \\
\end{longtable}
}

\paragraph{Poluição Luminosa}

{\def\LTcaptype{none} % do not increment counter
\begin{longtable}[]{@{}ll@{}}
\toprule\noalign{}
Parâmetro & Valor \\
\midrule\noalign{}
\endhead
\bottomrule\noalign{}
\endlastfoot
Escala Bortle & 2 --- Céu tipicamente escuro \\
Radiância artificial & 0.3 nW/\allowbreak cm²/\allowbreak sr \\
\end{longtable}
}
\paragraph{Solo (SoilGrids 2.0, média ponderada 0--30
cm)}

{\def\LTcaptype{none} % do not increment counter
\begin{longtable}[]{@{}ll@{}}
\toprule\noalign{}
Parâmetro & Valor \\
\midrule\noalign{}
\endhead
\bottomrule\noalign{}
\endlastfoot
pH (H$_2$O) & 6.05 \\
Carbono orgânico (SOC) & 14.09 g/\allowbreak kg \\
Argila & 28.1 \% \\
Areia & 38.17 \% \\
Silte (calc.) & 33.7 \% \\
Densidade aparente & 1.37 g/\allowbreak cm³ \\
\textbf{Aptidão agrícola} & \textbf{Alta} \\
\end{longtable}
}

\subsubsection{Indicadores Sociais}

\textbf{Fontes:} PNUD 2020, INE\footnote{\fineINE} Censo 2022, INE\footnote{\fineINE} EPHC 2023, Ministerio
Público 2024\\
\textbf{Nota:} Valores marcados como \emph{(dept.)} referem-se ao
departamento de Presidente Hayes; valores \emph{(dist.)} são específicos
deste distrito. Dados marcados \emph{(est.)} são estimativas.

{\def\LTcaptype{none} % do not increment counter
\begin{longtable}[]{@{}
  >{\raggedright\arraybackslash}p{(\linewidth - 4\tabcolsep) * \real{0.4231}}
  >{\raggedright\arraybackslash}p{(\linewidth - 4\tabcolsep) * \real{0.2692}}
  >{\raggedright\arraybackslash}p{(\linewidth - 4\tabcolsep) * \real{0.3077}}@{}}
\toprule\noalign{}
\begin{minipage}[b]{\linewidth}\raggedright
Indicador
\end{minipage} & \begin{minipage}[b]{\linewidth}\raggedright
Valor
\end{minipage} & \begin{minipage}[b]{\linewidth}\raggedright
Âmbito
\end{minipage} \\
\midrule\noalign{}
\endhead
\bottomrule\noalign{}
\endlastfoot
IDH (2020) & 0,693 & dept. \\
Ranking IDH nacional & 13/\allowbreak 18 & dept. \\
Esperança de vida & 71,8 anos & dept. \\
Escolaridade média & 9.2 anos & dist. \\
RNB per capita (USD PPA) & 7.900 & dept. \\
Pobreza monetária (\%) & 29,4\% & dept. \\
Pobreza extrema (\%) & 7,2\% & dept. \\
Índice de Gini & 0,453 & dept. \\
Acesso a água potável (\%) & 71,3\% & dept. \\
Acesso a saneamento (\%) & 68,9\% & dept. \\
Taxa de homicídios (est., /\allowbreak 100k hab) & 5,5--5,6 (2019--2020, fonte INE
--- próxima à média nacional) & dept. \\
Índice de segurança & médio & dept. \\
População (Censo 2022) & N/\allowbreak D & dist. \\
Idade mediana & N/\allowbreak D & dist. \\
Taxa de fecundidade & N/\allowbreak D & dist. \\
\end{longtable}
}
\subsubsection{Combustível}

\textbf{Referência:} PETROPAR /\allowbreak  postos locais (2024) \textbf{Tipo de
localidade:} Interior

{\def\LTcaptype{none} % do not increment counter
\begin{longtable}[]{@{}lll@{}}
\toprule\noalign{}
Combustível & USD/\allowbreak litro & Gs/\allowbreak litro (aprox.) \\
\midrule\noalign{}
\endhead
\bottomrule\noalign{}
\endlastfoot
Gasolina 93 oct & 1.05 & 7,770 \\
Gasolina 97 oct (premium) & 1.15 & 8,510 \\
Diesel & 0.98 & 7,252 \\
\end{longtable}
}

\begin{quote}
Preços podem variar ±5\% conforme posto e sazonalidade. Chaco e interior
remoto apresentam maior variação.
\end{quote}

\subsubsection{Cobertura Celular}

\textbf{Fonte:} CONATEL PY /\allowbreak  operadoras (2024)

{\def\LTcaptype{none} % do not increment counter
\begin{longtable}[]{@{}ll@{}}
\toprule\noalign{}
Parâmetro & Valor \\
\midrule\noalign{}
\endhead
\bottomrule\noalign{}
\endlastfoot
Cobertura 4G população (dept.) & 60\% \\
Cobertura 4G área rural & 35\% \\
Melhor operadora & Tigo \\
Qualidade rural & limitada \\
\end{longtable}
}

\begin{quote}
Para áreas rurais fora do núcleo urbano, recomenda-se chip Tigo como
principal e Personal como backup.
\end{quote}

\subsubsection{Internet}

\textbf{Fonte:} CONATEL /\allowbreak  Speedtest Ookla (2024)

{\def\LTcaptype{none} % do not increment counter
\begin{longtable}[]{@{}ll@{}}
\toprule\noalign{}
Parâmetro & Valor \\
\midrule\noalign{}
\endhead
\bottomrule\noalign{}
\endlastfoot
Velocidade média download & 25 Mbps \\
Domicílios com internet (dept.) & 35\% \\
Tecnologia predominante & rádio/\allowbreak satélite \\
Opção rural & Starlink disponível (\textasciitilde USD 44/\allowbreak mês) \\
\end{longtable}
}

\subsubsection{Mercado Imobiliário e Terra
Rural}

\textbf{Fonte:} INDERT /\allowbreak  Clasificados.com.py (2024)

{\def\LTcaptype{none} % do not increment counter
\begin{longtable}[]{@{}ll@{}}
\toprule\noalign{}
Tipo & Referência \\
\midrule\noalign{}
\endhead
\bottomrule\noalign{}
\endlastfoot
Terra agrícola alta prod. (USD/\allowbreak ha) & 2,000 \\
Imóvel urbano (USD/\allowbreak m²) & 600 \\
Aluguel 2 quartos (USD/\allowbreak mês) & 250 \\
\end{longtable}
}

\begin{quote}
Valores de referência departamental. Localidades menores podem ter
preços 20--40\% abaixo da capital departamental.
\end{quote}

\subsubsection{Saúde}

\textbf{Fonte:} MSPBS /\allowbreak  IPS Paraguay (2024-2026), consolidação
departamental e proxy local

{\def\LTcaptype{none} % do not increment counter
\begin{longtable}[]{@{}ll@{}}
\toprule\noalign{}
Serviço & Disponibilidade \\
\midrule\noalign{}
\endhead
\bottomrule\noalign{}
\endlastfoot
USF /\allowbreak  Posto de Saúde & sim \\
Hospital Regional & não \\
IPS (seguro social) & não \\
Farmácia & sim \\
Distância ao hospital de referência & \textasciitilde29 km (Villa
Hayes) \\
\end{longtable}
}

\textbf{Principais estabelecimentos:} USF local; referencia hospitalar
em Villa Hayes

\textbf{Observação para imigrantes:} Atendimento primario local ou em
raio curto; casos de maior complexidade seguem para Villa Hayes.
Cobertura privada continua recomendada para especialidades e urgências
de maior complexidade.
\newpage
\secaoDiagnostico{Nanawa}{\mapaDistrito{1505}}{Nanawa é o nó comercial de maior tensão no Bajo Chaco. Oferece acesso
privilegiado a recursos e suprimentos transfronteiriços, agora reforçado
pela conectividade imediata com a capital via Ponte Héroes del Chaco.
Sua força reside no poder do comércio e na abundância hídrica. Contudo,
sua extrema vulnerabilidade hidrológica e a alta visibilidade como alvo
tático fronteiriço a tornam uma localidade de risco elevado em cenários
de instabilidade nacional. É o local ideal para bases de suprimento e
logística de apoio tático, exigindo do ocupante autonomia energética e
planos de evacuação fluvial em caso de colapso terrestre.

\subsubsection{Por que sim}

\subsubsection{Por que não}

\subsubsection{Recomendação
Técnica}

Indicado estritamente como hub logístico de apoio ou base de
suprimentos. Não recomendado para residência permanente de refúgio
tático isolado. O GSS de 5.8 reflete o poder dos recursos equilibrado
pela fragilidade ambiental e tática da localidade. Referência atual de
score: GSS 5.8.

\begin{itemize}
\tightlist
\item
  \begin{enumerate}
  \def\labelenumi{\arabic{enumi})}
  \tightlist
  \item
    Dados censitários validados (INE\footnote{\fineINE} 2022).
  \end{enumerate}
\item
  \begin{enumerate}
  \def\labelenumi{\arabic{enumi})}
  \setcounter{enumi}{1}
  \tightlist
  \item
    Relatórios de monitoramento de cheias do Rio Pilcomayo (SEN\footnote{\fineSEN}/\allowbreak DMH).
  \end{enumerate}
\item
  \begin{enumerate}
  \def\labelenumi{\arabic{enumi})}
  \setcounter{enumi}{2}
  \tightlist
  \item
    Mapa de impacto socioeconômico da Ponte Héroes del Chaco (MOPC).
  \end{enumerate}
\item
  \begin{enumerate}
  \def\labelenumi{\arabic{enumi})}
  \setcounter{enumi}{3}
  \tightlist
  \item
    Registro de segurança pública e governança fronteiriça.
  \end{enumerate}
\end{itemize}}{dist:15_Presidente_Hayes:Nanawa}
\begin{table}[ht!]
\small \begin{tabular}{p{3.0cm}p{0.8cm}p{6.0cm}} \toprule
\textbf{Critério} & \textbf{Nota} & \textbf{Justificativa} \\ \midrule
A - Ameaças Estratégicas & 4.5 & Ponto de fronteira internacional; alvo logístico de alta visibilidade e sensibilidade geopolítica \\
B - Risco Social & 5.5 & Densidade urbana em área reduzida; risco social e de agitação urbana significativo \\
C - Riscos Naturais & 4.0 & Risco hidrológico extremo; nota penalizada pela dependência vital de infraestrutura de defesa \\
D - Autossuficiência & 8.5 & Excepcionais recursos comerciais e de suprimentos; abundância hídrica fluvial \\
E - Institucional & 6.5 & Forte segurança institucional de fronteira e presença aduaneira funcional \\
\midrule \textbf{GSS FINAL} & \textbf{5.8} & \textbf{Moderadamente Seguro (Indicado exclusivamente como hub comercial e logístico de apoio, não recomendado para refúgio tático isolado).} \\ \bottomrule \end{tabular}
\end{table}
\subsubsection{Vulnerabilidades e
Mitigação}

\begin{itemize}
\tightlist
\item
  \textbf{Centralização Política:} Alvo óbvio em crises institucionais
  ou conflitos (Mitigação: manter residência secundária em zona
  rural/\allowbreak periférica).
\item
  \textbf{Dependência de Serviços:} Colapso de rede elétrica ou água
  atinge o núcleo urbano rapidamente (Mitigação: geradores solares e
  cisternas).
\item
  \textbf{Risco de Inundação:} Evitar áreas de baixa cota próximas ao
  rio.
\item
  \textbf{Agitação Social:} Centro de manifestações nacionais
  (Mitigação: localização fora do eixo central/\allowbreak administrativo).
\end{itemize}

\subsubsection{Dossiê de Campo}
\begin{description}[style=nextline,leftmargin=0.5cm,font=\bfseries]
\item[Coordenadas]  25°16'S 57°38'W.
\item[Topografia]  Colinas suaves (geologicamente estável, risco sísmico desprezível). Margeada pelo Rio Paraguai.
\item[Alvos Estratégicos]  Palácio de los López (Governo), Porto Sajonia (Marinha), Base Ñu Guasu (Aeronáutica/\allowbreak Exército), Aeroporto Silvio Pettirossi (Luque), Hub elétrico de transmissão de Itaipu.
\item[Fallout]  Ventos predominantes: Nordeste (verão - Sirocco, quente) e Sul (inverno - Pampero, frio). Baixo risco de fallout global, mas vulnerável se houver ataques a centros de comando nacionais.
\item[Mídia local]  Principais meios incluem ABC Color, Última Hora, NPY, SNT e rádios AM/\allowbreak FM locais. Cobertura 4G ampla na área urbana.
\item[População]  464.185 em 2025 (INE\footnote{\fineINE}). Censo 2022 Final: 462.241. Área metropolitana (Gran Asunción) com \textasciitilde 2.200.000.
\item[Densidade]  \textasciitilde 3.950 hab/\allowbreak km² (Núcleo urbano denso).
\item[Vias de Acesso]  Hub central do país. Rutas PY01, PY02, PY03. Gargalos em pontes sobre o Rio Paraguai (Remanso, Héroes del Chaco).
\item[Serviços]  Melhor infraestrutura hospitalar e educacional do país. Investimentos maciços em novos centros de saúde.
\item[Custo de Vida]  US\$ 600-1.550/\allowbreak mês (dependendo do padrão). Terrenos urbanos: US\$ 400-650/\allowbreak m².
\item[Saúde de alta complexidade]  IMT, INCAN, INE\footnote{\fineINE}RAM, Hospital General Barrio Obrero, Hospital del Trauma, Hospital Materno Infantil San Pablo, Hospital Materno Infantil Santísima Trinidad, Hospital Materno Infantil Loma Pyta, Hospital Psiquiátrico e Hospital del Indígena concentram oferta especializada na capital.
\item[Universidades habilitadas]  Asunción concentra oferta do CONES, com instituições como Universidad Americana, Universidad Autónoma de Asunción e Universidad Adventista del Paraguay.
\item[Conexões logísticas]  Porto de Asunción, sistema 911 e Aeroporto Internacional Silvio Pettirossi reforçam a função de nó administrativo, portuário e metropolitano.
\item[Idade mediana]  33 anos em 2025, segundo projeções do INE\footnote{\fineINE}.
\item[Fecundidade]  1,64 filho por mulher em 2022 e 1,68 em 2025 nas projeções do INE\footnote{\fineINE}.
\item[Escolaridade]  11,7 anos de estudo em homens e 11,4 em mulheres na faixa de 15+ anos; a cidade concentra o maior capital educacional do país.
\item[Assistência escolar]  os indicadores nacionais da EPHC 2025 mostram 39,4\% de frequência geral e 97,7\% de assistência entre 5 e 9 anos.
\item[Hidrologia]  Risco de Inundação Recorrente. Áreas ribeirinhas (Bañados) sofrem cheias periódicas ligadas ao El Niño.
\item[Clima]  Tropical/\allowbreak Subtropical. Verões com temperaturas >40°C. Tempestades severas com rajadas de vento de até 160 km/\allowbreak h.
\item[Sismicidade]  Muito baixa. Região tectonicamente estável.
\item[Inclinação recomendada para placas solares]  25° voltado para o Norte (anual); 35° no inverno (Jun-Ago); 15° no verão (Nov-Jan). Base: latitude local -25.28°S.
\item[Energia]  Dependência total da rede nacional (Itaipu), mas com a energia mais barata da região (\textasciitilde US\$ 0,05/\allowbreak kWh).
\item[Água]  Abundante (Rio Paraguai), tratamento centralizado pela ESSAP\footnote{\fineESSAP}.
\item[Solo]  Solo fértil na Região Oriental, mas área urbana saturada. Potencial agrícola nos arredores (Central/\allowbreak Chaco'i).
\item[Autossuficiência]  Baixa no núcleo urbano; alta dependência de suprimentos logísticos nacionais e internacionais.
\item[Internet de alta velocidade]  110 pontos de WiFi libre foram habilitados em Asunción pelo MITIC; 320 novos pontos chegaram ao país em 2023, com cobertura nacional ampliada.
\item[Rede celular]  a tecnologia 5G foi lançada oficialmente no país em dezembro de 2025 e teve ato em Asunción, no Paseo La Galería.
\item[Poços e águas subterrâneas]  não há profundidade média oficial consolidada para poços artesianos em Asunción; a regulação permanece com MADES e Lei 3239/\allowbreak 07.
\item[Combustível e gás]  não foi identificado inventário oficial único por distrito para postos, eletropostos e distribuição de gás encanado.
\item[Serviço público digital]  a municipalidade de Asunción iniciou processo de transformação digital com apoio do MITIC em 2025.
\item[Segurança]  Taxa de homicídios baixa para padrões latinos (7.1/\allowbreak 100k). Criminalidade focada em delitos contra o patrimônio. Áreas nobres e corporativas são consideradas seguras.
\item[Leis Local]  Sede do governo central. Direitos de propriedade bem definidos. Incentivos fiscais (Lei 60/\allowbreak 90).
\item[Delegacias]  Dirección de Policía de Asunción com 23 comisarías jurisdicionais listadas pela Policía Nacional.
\item[Presídios]  Tacumbú concentra 2.151 PPL em censo recente do Ministério de Justicia; o conjunto penitenciário de Asunción também inclui UPIE e Granja Ko'e Pyahu, com mais de 2.200 PPL em visita institucional de 2025.
\item[Correios e entregas]  Correos del Paraguay opera cobertura nacional e atende a capital; não há tabela oficial única de desempenho local por bairro.
\item[Comércio urbano]  o MIC publicou o REPSE do 1º quadrimestre de 2025 com Asunción concentrando 32,2\% dos novos prestadores de serviços do país, o que funciona como proxy oficial de dinamismo comercial formal; o inventário centralizado de cinemas, shoppings e atacados segue sem tabela única.
\item[Escolaridade e ensino superior]  INE\footnote{\fineINE} e CONES são as bases centrais para alfabetização, escolaridade e oferta universitária; a extração fina por recorte territorial permanece parcial.
\item[Demografia social]  idade mediana de 33 anos e fecundidade de 1,68 filho por mulher em 2025; o item já saiu da zona de lacuna aberta.
\item[Poluição luminosa]  Bortle 8 (céu urbano), radiância artificial estimada em 100 nW/\allowbreak cm²/\allowbreak sr. Fonte: World Atlas of Artificial Night Sky Brightness.
\end{description}
\textbf{Fonte climática:} NASA POWER Climatology API (período 2001-2020)

\textbf{Fonte luminosa:} estimativa world atlas

\paragraph{Irradiação Solar
(kWh/\allowbreak m²/\allowbreak dia)}

{\def\LTcaptype{none} % do not increment counter
\begin{longtable}[]{@{}
  >{\raggedright\arraybackslash}p{(\linewidth - 24\tabcolsep) * \real{0.0746}}
  >{\raggedright\arraybackslash}p{(\linewidth - 24\tabcolsep) * \real{0.0746}}
  >{\raggedright\arraybackslash}p{(\linewidth - 24\tabcolsep) * \real{0.0746}}
  >{\raggedright\arraybackslash}p{(\linewidth - 24\tabcolsep) * \real{0.0746}}
  >{\raggedright\arraybackslash}p{(\linewidth - 24\tabcolsep) * \real{0.0746}}
  >{\raggedright\arraybackslash}p{(\linewidth - 24\tabcolsep) * \real{0.0746}}
  >{\raggedright\arraybackslash}p{(\linewidth - 24\tabcolsep) * \real{0.0746}}
  >{\raggedright\arraybackslash}p{(\linewidth - 24\tabcolsep) * \real{0.0746}}
  >{\raggedright\arraybackslash}p{(\linewidth - 24\tabcolsep) * \real{0.0746}}
  >{\raggedright\arraybackslash}p{(\linewidth - 24\tabcolsep) * \real{0.0746}}
  >{\raggedright\arraybackslash}p{(\linewidth - 24\tabcolsep) * \real{0.0746}}
  >{\raggedright\arraybackslash}p{(\linewidth - 24\tabcolsep) * \real{0.0746}}
  >{\raggedright\arraybackslash}p{(\linewidth - 24\tabcolsep) * \real{0.1045}}@{}}
\toprule\noalign{}
\begin{minipage}[b]{\linewidth}\raggedright
Jan
\end{minipage} & \begin{minipage}[b]{\linewidth}\raggedright
Fev
\end{minipage} & \begin{minipage}[b]{\linewidth}\raggedright
Mar
\end{minipage} & \begin{minipage}[b]{\linewidth}\raggedright
Abr
\end{minipage} & \begin{minipage}[b]{\linewidth}\raggedright
Mai
\end{minipage} & \begin{minipage}[b]{\linewidth}\raggedright
Jun
\end{minipage} & \begin{minipage}[b]{\linewidth}\raggedright
Jul
\end{minipage} & \begin{minipage}[b]{\linewidth}\raggedright
Ago
\end{minipage} & \begin{minipage}[b]{\linewidth}\raggedright
Set
\end{minipage} & \begin{minipage}[b]{\linewidth}\raggedright
Out
\end{minipage} & \begin{minipage}[b]{\linewidth}\raggedright
Nov
\end{minipage} & \begin{minipage}[b]{\linewidth}\raggedright
Dez
\end{minipage} & \begin{minipage}[b]{\linewidth}\raggedright
Média
\end{minipage} \\
\midrule\noalign{}
\endhead
\bottomrule\noalign{}
\endlastfoot
6.76 & 6.20 & 5.42 & 4.44 & 3.36 & 2.83 & 3.21 & 3.93 & 4.61 & 5.49 &
6.42 & 6.78 & \textbf{4.95} \\
\end{longtable}
}

\textbf{Inclinação solar recomendada:} 25° N (anual) · 35° N (inverno
jun-ago) · 15° N (verão nov-jan)

\paragraph{Precipitação (mm/\allowbreak mês)}

{\def\LTcaptype{none} % do not increment counter
\begin{longtable}[]{@{}
  >{\raggedright\arraybackslash}p{(\linewidth - 24\tabcolsep) * \real{0.0704}}
  >{\raggedright\arraybackslash}p{(\linewidth - 24\tabcolsep) * \real{0.0704}}
  >{\raggedright\arraybackslash}p{(\linewidth - 24\tabcolsep) * \real{0.0704}}
  >{\raggedright\arraybackslash}p{(\linewidth - 24\tabcolsep) * \real{0.0704}}
  >{\raggedright\arraybackslash}p{(\linewidth - 24\tabcolsep) * \real{0.0704}}
  >{\raggedright\arraybackslash}p{(\linewidth - 24\tabcolsep) * \real{0.0704}}
  >{\raggedright\arraybackslash}p{(\linewidth - 24\tabcolsep) * \real{0.0704}}
  >{\raggedright\arraybackslash}p{(\linewidth - 24\tabcolsep) * \real{0.0704}}
  >{\raggedright\arraybackslash}p{(\linewidth - 24\tabcolsep) * \real{0.0704}}
  >{\raggedright\arraybackslash}p{(\linewidth - 24\tabcolsep) * \real{0.0704}}
  >{\raggedright\arraybackslash}p{(\linewidth - 24\tabcolsep) * \real{0.0704}}
  >{\raggedright\arraybackslash}p{(\linewidth - 24\tabcolsep) * \real{0.0704}}
  >{\raggedright\arraybackslash}p{(\linewidth - 24\tabcolsep) * \real{0.1549}}@{}}
\toprule\noalign{}
\begin{minipage}[b]{\linewidth}\raggedright
Jan
\end{minipage} & \begin{minipage}[b]{\linewidth}\raggedright
Fev
\end{minipage} & \begin{minipage}[b]{\linewidth}\raggedright
Mar
\end{minipage} & \begin{minipage}[b]{\linewidth}\raggedright
Abr
\end{minipage} & \begin{minipage}[b]{\linewidth}\raggedright
Mai
\end{minipage} & \begin{minipage}[b]{\linewidth}\raggedright
Jun
\end{minipage} & \begin{minipage}[b]{\linewidth}\raggedright
Jul
\end{minipage} & \begin{minipage}[b]{\linewidth}\raggedright
Ago
\end{minipage} & \begin{minipage}[b]{\linewidth}\raggedright
Set
\end{minipage} & \begin{minipage}[b]{\linewidth}\raggedright
Out
\end{minipage} & \begin{minipage}[b]{\linewidth}\raggedright
Nov
\end{minipage} & \begin{minipage}[b]{\linewidth}\raggedright
Dez
\end{minipage} & \begin{minipage}[b]{\linewidth}\raggedright
Total/\allowbreak ano
\end{minipage} \\
\midrule\noalign{}
\endhead
\bottomrule\noalign{}
\endlastfoot
126 & 139 & 135 & 149 & 143 & 74 & 63 & 38 & 73 & 163 & 190 & 179 &
\textbf{1477 mm} \\
\end{longtable}
}

\paragraph{Poluição Luminosa}

{\def\LTcaptype{none} % do not increment counter
\begin{longtable}[]{@{}ll@{}}
\toprule\noalign{}
Parâmetro & Valor \\
\midrule\noalign{}
\endhead
\bottomrule\noalign{}
\endlastfoot
Escala Bortle & 8 --- Céu urbano \\
Radiância artificial & 100.0 nW/\allowbreak cm²/\allowbreak sr \\
\end{longtable}
}
\paragraph{Solo (SoilGrids 2.0, média ponderada 0--30
cm)}

{\def\LTcaptype{none} % do not increment counter
\begin{longtable}[]{@{}ll@{}}
\toprule\noalign{}
Parâmetro & Valor \\
\midrule\noalign{}
\endhead
\bottomrule\noalign{}
\endlastfoot
pH (H$_2$O) & 5.72 \\
Carbono orgânico (SOC) & 24.33 g/\allowbreak kg \\
Argila & 19.82 \% \\
Areia & 56.72 \% \\
Silte (calc.) & 23.5 \% \\
Densidade aparente & 1.21 g/\allowbreak cm³ \\
\textbf{Aptidão agrícola} & \textbf{Alta} \\
\end{longtable}
}

\begin{itemize}
\tightlist
\item
  \textbf{Autossuficiência:} Baixa no núcleo urbano; alta dependência de
  suprimentos logísticos nacionais e internacionais.
\item
  \textbf{Internet de alta velocidade:} 110 pontos de WiFi libre foram
  habilitados em Asunción pelo MITIC; 320 novos pontos chegaram ao país
  em 2023, com cobertura nacional ampliada.
\item
  \textbf{Rede celular:} a tecnologia 5G foi lançada oficialmente no
  país em dezembro de 2025 e teve ato em Asunción, no Paseo La Galería.
\item
  \textbf{Poços e águas subterrâneas:} não há profundidade média oficial
  consolidada para poços artesianos em Asunción; a regulação permanece
  com MADES e Lei 3239/\allowbreak 07.
\item
  \textbf{Combustível e gás:} não foi identificado inventário oficial
  único por distrito para postos, eletropostos e distribuição de gás
  encanado.
\item
  \textbf{Serviço público digital:} a municipalidade de Asunción iniciou
  processo de transformação digital com apoio do MITIC em 2025.
\end{itemize}
\subsubsection{Indicadores Sociais}

\textbf{Fontes:} PNUD 2020, INE\footnote{\fineINE} Censo 2022, INE\footnote{\fineINE} EPHC 2023, Ministerio
Público 2024\\
\textbf{Nota:} Valores marcados como \emph{(dept.)} referem-se ao
departamento de Distrito Capital; valores \emph{(dist.)} são específicos
deste distrito. Dados marcados \emph{(est.)} são estimativas.

{\def\LTcaptype{none} % do not increment counter
\begin{longtable}[]{@{}
  >{\raggedright\arraybackslash}p{(\linewidth - 4\tabcolsep) * \real{0.4231}}
  >{\raggedright\arraybackslash}p{(\linewidth - 4\tabcolsep) * \real{0.2692}}
  >{\raggedright\arraybackslash}p{(\linewidth - 4\tabcolsep) * \real{0.3077}}@{}}
\toprule\noalign{}
\begin{minipage}[b]{\linewidth}\raggedright
Indicador
\end{minipage} & \begin{minipage}[b]{\linewidth}\raggedright
Valor
\end{minipage} & \begin{minipage}[b]{\linewidth}\raggedright
Âmbito
\end{minipage} \\
\midrule\noalign{}
\endhead
\bottomrule\noalign{}
\endlastfoot
IDH (2020) & 0,835 & dept. \\
Ranking IDH nacional & 1/\allowbreak 18 & dept. \\
Esperança de vida & 76,2 anos & dept. \\
Escolaridade média & 11,6 anos (média homens 11,7 /\allowbreak  mulheres 11,4) &
dist. \\
RNB per capita (USD PPA) & 18.200 & dept. \\
Pobreza monetária (\%) & 7,3\% & dept. \\
Pobreza extrema (\%) & 1,2\% & dept. \\
Índice de Gini & 0,435 & dept. \\
Acesso a água potável (\%) & 98,5\% & dept. \\
Acesso a saneamento (\%) & 99,1\% & dept. \\
Taxa de homicídios (est., /\allowbreak 100k hab) & \textasciitilde4,5 (estimativa) &
dept. \\
Índice de segurança & médio & dept. \\
População (Censo 2022) & 462.241 hab. & dist. \\
Idade mediana & 33 anos & dist. \\
Taxa de fecundidade & 1,64 filhos/\allowbreak mulher & dist. \\
\end{longtable}
}
\subsubsection{Combustível}

\textbf{Referência:} PETROPAR /\allowbreak  postos locais (2024) \textbf{Tipo de
localidade:} Capital departamental

{\def\LTcaptype{none} % do not increment counter
\begin{longtable}[]{@{}lll@{}}
\toprule\noalign{}
Combustível & USD/\allowbreak litro & Gs/\allowbreak litro (aprox.) \\
\midrule\noalign{}
\endhead
\bottomrule\noalign{}
\endlastfoot
Gasolina 93 oct & 0.95 & 7,030 \\
Gasolina 97 oct (premium) & 1.05 & 7,770 \\
Diesel & 0.88 & 6,512 \\
\end{longtable}
}

\begin{quote}
Preços podem variar ±5\% conforme posto e sazonalidade. Chaco e interior
remoto apresentam maior variação.
\end{quote}

\subsubsection{Cobertura Celular}

\textbf{Fonte:} CONATEL PY /\allowbreak  operadoras (2024)

{\def\LTcaptype{none} % do not increment counter
\begin{longtable}[]{@{}ll@{}}
\toprule\noalign{}
Parâmetro & Valor \\
\midrule\noalign{}
\endhead
\bottomrule\noalign{}
\endlastfoot
Cobertura 4G população (dept.) & 98\% \\
Cobertura 4G área rural & 95\% \\
Melhor operadora & Tigo/\allowbreak Personal \\
Qualidade rural & boa \\
\end{longtable}
}

\begin{quote}
Para áreas rurais fora do núcleo urbano, recomenda-se chip Tigo como
principal e Personal como backup.
\end{quote}

\subsubsection{Internet}

\textbf{Fonte:} CONATEL /\allowbreak  Speedtest Ookla (2024)

{\def\LTcaptype{none} % do not increment counter
\begin{longtable}[]{@{}ll@{}}
\toprule\noalign{}
Parâmetro & Valor \\
\midrule\noalign{}
\endhead
\bottomrule\noalign{}
\endlastfoot
Velocidade média download & 92 Mbps \\
Domicílios com internet (dept.) & 89\% \\
Tecnologia predominante & fibra/\allowbreak cabo \\
Opção rural & fibra óptica disponível \\
\end{longtable}
}

\subsubsection{Mercado Imobiliário e Terra
Rural}

\textbf{Fonte:} INDERT /\allowbreak  Clasificados.com.py (2024)

{\def\LTcaptype{none} % do not increment counter
\begin{longtable}[]{@{}ll@{}}
\toprule\noalign{}
Tipo & Referência \\
\midrule\noalign{}
\endhead
\bottomrule\noalign{}
\endlastfoot
Terra agrícola alta prod. (USD/\allowbreak ha) & 36,000 \\
Imóvel urbano (USD/\allowbreak m²) & 1,724 \\
Aluguel 2 quartos (USD/\allowbreak mês) & 840 \\
\end{longtable}
}

\begin{quote}
Valores de referência departamental. Localidades menores podem ter
preços 20--40\% abaixo da capital departamental.
\end{quote}

\subsubsection{Saúde}

\textbf{Fonte:} MSPBS /\allowbreak  IPS Paraguay (2024)

{\def\LTcaptype{none} % do not increment counter
\begin{longtable}[]{@{}ll@{}}
\toprule\noalign{}
Serviço & Disponibilidade \\
\midrule\noalign{}
\endhead
\bottomrule\noalign{}
\endlastfoot
USF /\allowbreak  Posto de Saúde & sim \\
Hospital Regional & sim \\
IPS (seguro social) & sim \\
Farmácia & sim \\
Distância ao hospital de referência & local \\
\end{longtable}
}

\textbf{Principais estabelecimentos:} Hospital Central do IPS (Santo
Domingo), Hospital de Clínicas (UNA), IMT (Instituto de Medicina
Tropical), INCAN (oncologia), INE\footnote{\fineINE}RAM (pneumologia), Hospital General
Barrio Obrero, Hospital del Trauma, Hospital Materno Infantil San Pablo,
Hospital Materno Infantil Santísima Trinidad, Hospital Materno Infantil
Loma Pyta, Hospital Psiquiátrico

\textbf{Observação para imigrantes:} Asunción concentra a oferta de alta
complexidade do país. IPS acessível para trabalhadores formais
registrados. Cobertura privada recomendada para quem busca agilidade;
planos nacionais (Tesãi, Medicall, etc.) disponíveis. Seguro saúde
complementar desejável.
\newpage
\secaoDiagnostico{Nueva Asunción}{\mapaDistrito{1511}}{Nueva Asunción (Chaco'i) é o novo epicentro do desenvolvimento
metropolitano paraguaio. Oferece uma infraestrutura logística de elite e
acesso imediato aos recursos da capital, funcionando como uma extensão
moderna e planejada de Assunção no Chaco. Sua força reside na
excepcional conectividade e no suporte institucional/\allowbreak privado. Contudo,
sua proximidade geográfica com o centro de poder nacional a torna
vulnerável a fallout tático e instabilidade social urbana direta. É o
local ideal para bases administrativas de alto padrão, exigindo do
ocupante autonomia técnica absoluta em termos de segurança ambiental
(NBQ) e hídrica.

\subsubsection{Por que sim}

\subsubsection{Por que não}

\subsubsection{Recomendação
Técnica}

Altamente indicado como base logística ou administrativa de alto padrão.
Para moradia permanente, recomenda-se a aquisição de propriedades em
condomínios com infraestrutura de defesa civil integrada. O GSS de 6.1
reflete a excelência dos recursos equilibrada pela vulnerabilidade
tática e ambiental da localidade. Referência atual de score: GSS 6.1.

\begin{itemize}
\tightlist
\item
  \begin{enumerate}
  \def\labelenumi{\arabic{enumi})}
  \tightlist
  \item
    Dados censitários validados (INE\footnote{\fineINE} 2022).
  \end{enumerate}
\item
  \begin{enumerate}
  \def\labelenumi{\arabic{enumi})}
  \setcounter{enumi}{1}
  \tightlist
  \item
    Relatórios de engenharia e impacto da Ponte Héroes del Chaco (MOPC).
  \end{enumerate}
\item
  \begin{enumerate}
  \def\labelenumi{\arabic{enumi})}
  \setcounter{enumi}{2}
  \tightlist
  \item
    Mapa de áreas inundáveis e cotas de refulado (SEN\footnote{\fineSEN}/\allowbreak MOPC).
  \end{enumerate}
\item
  \begin{enumerate}
  \def\labelenumi{\arabic{enumi})}
  \setcounter{enumi}{3}
  \tightlist
  \item
    Registro de segurança pública e investimentos privados em Chaco'i.
  \end{enumerate}
\end{itemize}}{dist:15_Presidente_Hayes:Nueva_Asunción}
\begin{table}[ht!]
\small \begin{tabular}{p{3.0cm}p{0.8cm}p{6.0cm}} \toprule
\textbf{Critério} & \textbf{Nota} & \textbf{Justificativa} \\ \midrule
A - Ameaças Estratégicas & 4.5 & Ponto de entrada estratégico para o Chaco; alvo de alta visibilidade e proximidade crítica com a capital \\
B - Risco Social & 6.0 & População pequena mas em rápido crescimento; densidade manejada por condomínios fechados \\
C - Riscos Naturais & 4.5 & Risco hidrológico extremo; nota penalizada pela vulnerabilidade a inundações extraordinárias \\
D - Autossuficiência & 8.5 & Excepcionais recursos logísticos e proximidade imediata com suprimentos metropolitanos \\
E - Institucional & 7.0 & Presença institucional em expansão e alta segurança privada tática no distrito \\
\midrule \textbf{GSS FINAL} & \textbf{6.1} & \textbf{Moderadamente Seguro (Indicado como base logística de elite ou residência de alto padrão, exigindo mitigação técnica rigorosa para riscos de fallout e inundações).} \\ \bottomrule \end{tabular}
\end{table}
\subsubsection{Vulnerabilidades e
Mitigação}

\begin{itemize}
\tightlist
\item
  \textbf{Insegurança Regional:} Conflitos entre insurgentes e forças de
  segurança (Mitigação: residência próxima a perímetros controlados,
  evitar áreas de mata densa desconhecidas, comunicação via rádio).
\item
  \textbf{Isolamento Logístico:} Estradas internas precárias (Mitigação:
  estoque de insumos para 30 dias e veículos 4x4 robustos).
\item
  \textbf{Conflitos de Terra:} Histórico de tensões rurais (Mitigação:
  verificação rigorosa da titulação legal/\allowbreak INDERT antes da aquisição).
\end{itemize}

\subsubsection{Dossiê de Campo}
\begin{description}[style=nextline,leftmargin=0.5cm,font=\bfseries]
\item[Coordenadas]  23°18'S 56°45'W.
\item[Topografia]  Relevo ondulado, área de transição entre planícies e elevações do Norte. Altitude \textasciitilde 150m. Geologicamente estável, risco sísmico desprezível.
\item[Alvos Estratégicos]  Bases da FTC (Força de Tarefa Conjunta) presentes na região para combate à insurgência. Proximidade com a Ruta PY05 (importante eixo Leste-Oeste).
\item[Fallout]  Ventos predominantes do Norte (médias de 8 km/\allowbreak h).
\item[População]  8.282 (Censo 2022 Final). População majoritariamente rural organizada em 7 núcleos de assentamento.
\item[Densidade]  \textasciitilde 9,2 hab/\allowbreak km² (Baixa densidade, oferecendo isolamento relativo em áreas de mata e fazendas).
\item[Vias de Acesso]  Acesso via Ruta PY05. Os caminhos internos entre os núcleos são de terra e tornam-se frequentemente intransitáveis em períodos de chuva severa.
\item[Serviços]  Unidade de Saúde da Família (USF) local para atendimentos básicos. Dependência de Horqueta ou Concepción (78 km) para serviços complexos. Presença do Colégio Técnico Agropecuário no Núcleo 3.
\item[Custo de Vida]  Baixo; economia baseada em agricultura familiar e subsistência.
\item[Preço da Terra]  US\$ 3.000-5.000/\allowbreak ha (áreas mecanizáveis/\allowbreak agricultura familiar).
\item[Hidrologia]  Risco de isolamento por transbordamento de arroios e lamaçais nas vias vicinais. Histórico de falta de água potável em alguns núcleos (ex: Núcleo 2).
\item[Clima]  Tropical Úmido. Precipitação anual entre 1.300 e 1.500 mm. Suscetibilidade a temporais de verão com ventos fortes.
\item[Energia]  Rede nacional (ANDE\footnote{\fineANDE}), sujeita a cortes por infraestrutura rural precária.
\item[Água]  Poços artesianos e arroios locais; sistemas de distribuição comunitária nos núcleos de assentamento.
\item[Qualidade do Solo]  Predomínio de latossolos vermelhos; férteis com correção, ideais para agricultura familiar, grãos e silvicultura.
\item[Recursos Locais]  Agricultura familiar (tomate, mandioca, grãos) e pecuária de pequeno/\allowbreak médio porte. Elevado potencial para autossuficiência alimentar básica em nível comunitário.
\item[Segurança]  Risco Crítico. Arroyito é historicamente uma das principais zonas de influência e operações do grupo guerrilheiro EPP\footnote{\fineEPP} (Exército do Povo Paraguaio). Presença militar massiva da FTC é necessária e constante. Taxa de homicídios regional de \textasciitilde 14/\allowbreak 100k.
\item[Leis Local: Município de criação recente (2016) originado de assentamento camponês. Restrição de Fronteira (Lei 2532/\allowbreak 05)]  Proibição de compra de terras rurais por estrangeiros de países limítrofes na faixa de 50km.
\end{description}
\textbf{Fonte climática:} NASA POWER Climatology API (período 2001-2020)

\textbf{Fonte luminosa:} estimativa world atlas

\paragraph{Irradiação Solar
(kWh/\allowbreak m²/\allowbreak dia)}

{\def\LTcaptype{none} % do not increment counter
\begin{longtable}[]{@{}
  >{\raggedright\arraybackslash}p{(\linewidth - 24\tabcolsep) * \real{0.0746}}
  >{\raggedright\arraybackslash}p{(\linewidth - 24\tabcolsep) * \real{0.0746}}
  >{\raggedright\arraybackslash}p{(\linewidth - 24\tabcolsep) * \real{0.0746}}
  >{\raggedright\arraybackslash}p{(\linewidth - 24\tabcolsep) * \real{0.0746}}
  >{\raggedright\arraybackslash}p{(\linewidth - 24\tabcolsep) * \real{0.0746}}
  >{\raggedright\arraybackslash}p{(\linewidth - 24\tabcolsep) * \real{0.0746}}
  >{\raggedright\arraybackslash}p{(\linewidth - 24\tabcolsep) * \real{0.0746}}
  >{\raggedright\arraybackslash}p{(\linewidth - 24\tabcolsep) * \real{0.0746}}
  >{\raggedright\arraybackslash}p{(\linewidth - 24\tabcolsep) * \real{0.0746}}
  >{\raggedright\arraybackslash}p{(\linewidth - 24\tabcolsep) * \real{0.0746}}
  >{\raggedright\arraybackslash}p{(\linewidth - 24\tabcolsep) * \real{0.0746}}
  >{\raggedright\arraybackslash}p{(\linewidth - 24\tabcolsep) * \real{0.0746}}
  >{\raggedright\arraybackslash}p{(\linewidth - 24\tabcolsep) * \real{0.1045}}@{}}
\toprule\noalign{}
\begin{minipage}[b]{\linewidth}\raggedright
Jan
\end{minipage} & \begin{minipage}[b]{\linewidth}\raggedright
Fev
\end{minipage} & \begin{minipage}[b]{\linewidth}\raggedright
Mar
\end{minipage} & \begin{minipage}[b]{\linewidth}\raggedright
Abr
\end{minipage} & \begin{minipage}[b]{\linewidth}\raggedright
Mai
\end{minipage} & \begin{minipage}[b]{\linewidth}\raggedright
Jun
\end{minipage} & \begin{minipage}[b]{\linewidth}\raggedright
Jul
\end{minipage} & \begin{minipage}[b]{\linewidth}\raggedright
Ago
\end{minipage} & \begin{minipage}[b]{\linewidth}\raggedright
Set
\end{minipage} & \begin{minipage}[b]{\linewidth}\raggedright
Out
\end{minipage} & \begin{minipage}[b]{\linewidth}\raggedright
Nov
\end{minipage} & \begin{minipage}[b]{\linewidth}\raggedright
Dez
\end{minipage} & \begin{minipage}[b]{\linewidth}\raggedright
Média
\end{minipage} \\
\midrule\noalign{}
\endhead
\bottomrule\noalign{}
\endlastfoot
6.60 & 6.02 & 5.62 & 4.65 & 3.48 & 3.07 & 3.35 & 4.15 & 4.72 & 5.46 &
6.29 & 6.53 & \textbf{4.99} \\
\end{longtable}
}

\textbf{Inclinação solar recomendada:} 23° N (anual) · 33° N (inverno
jun-ago) · 13° N (verão nov-jan)

\paragraph{Precipitação (mm/\allowbreak mês)}

{\def\LTcaptype{none} % do not increment counter
\begin{longtable}[]{@{}
  >{\raggedright\arraybackslash}p{(\linewidth - 24\tabcolsep) * \real{0.0704}}
  >{\raggedright\arraybackslash}p{(\linewidth - 24\tabcolsep) * \real{0.0704}}
  >{\raggedright\arraybackslash}p{(\linewidth - 24\tabcolsep) * \real{0.0704}}
  >{\raggedright\arraybackslash}p{(\linewidth - 24\tabcolsep) * \real{0.0704}}
  >{\raggedright\arraybackslash}p{(\linewidth - 24\tabcolsep) * \real{0.0704}}
  >{\raggedright\arraybackslash}p{(\linewidth - 24\tabcolsep) * \real{0.0704}}
  >{\raggedright\arraybackslash}p{(\linewidth - 24\tabcolsep) * \real{0.0704}}
  >{\raggedright\arraybackslash}p{(\linewidth - 24\tabcolsep) * \real{0.0704}}
  >{\raggedright\arraybackslash}p{(\linewidth - 24\tabcolsep) * \real{0.0704}}
  >{\raggedright\arraybackslash}p{(\linewidth - 24\tabcolsep) * \real{0.0704}}
  >{\raggedright\arraybackslash}p{(\linewidth - 24\tabcolsep) * \real{0.0704}}
  >{\raggedright\arraybackslash}p{(\linewidth - 24\tabcolsep) * \real{0.0704}}
  >{\raggedright\arraybackslash}p{(\linewidth - 24\tabcolsep) * \real{0.1549}}@{}}
\toprule\noalign{}
\begin{minipage}[b]{\linewidth}\raggedright
Jan
\end{minipage} & \begin{minipage}[b]{\linewidth}\raggedright
Fev
\end{minipage} & \begin{minipage}[b]{\linewidth}\raggedright
Mar
\end{minipage} & \begin{minipage}[b]{\linewidth}\raggedright
Abr
\end{minipage} & \begin{minipage}[b]{\linewidth}\raggedright
Mai
\end{minipage} & \begin{minipage}[b]{\linewidth}\raggedright
Jun
\end{minipage} & \begin{minipage}[b]{\linewidth}\raggedright
Jul
\end{minipage} & \begin{minipage}[b]{\linewidth}\raggedright
Ago
\end{minipage} & \begin{minipage}[b]{\linewidth}\raggedright
Set
\end{minipage} & \begin{minipage}[b]{\linewidth}\raggedright
Out
\end{minipage} & \begin{minipage}[b]{\linewidth}\raggedright
Nov
\end{minipage} & \begin{minipage}[b]{\linewidth}\raggedright
Dez
\end{minipage} & \begin{minipage}[b]{\linewidth}\raggedright
Total/\allowbreak ano
\end{minipage} \\
\midrule\noalign{}
\endhead
\bottomrule\noalign{}
\endlastfoot
165 & 157 & 117 & 134 & 112 & 71 & 46 & 32 & 73 & 146 & 205 & 165 &
\textbf{1428 mm} \\
\end{longtable}
}

\paragraph{Poluição Luminosa}

{\def\LTcaptype{none} % do not increment counter
\begin{longtable}[]{@{}ll@{}}
\toprule\noalign{}
Parâmetro & Valor \\
\midrule\noalign{}
\endhead
\bottomrule\noalign{}
\endlastfoot
Escala Bortle & 3 --- Céu rural \\
Radiância artificial & 1.5 nW/\allowbreak cm²/\allowbreak sr \\
\end{longtable}
}
\paragraph{Solo (SoilGrids 2.0, média ponderada 0--30
cm)}

{\def\LTcaptype{none} % do not increment counter
\begin{longtable}[]{@{}ll@{}}
\toprule\noalign{}
Parâmetro & Valor \\
\midrule\noalign{}
\endhead
\bottomrule\noalign{}
\endlastfoot
pH (H$_2$O) & 5.5 \\
Carbono orgânico (SOC) & 17.31 g/\allowbreak kg \\
Argila & 17.4 \% \\
Areia & 62.82 \% \\
Silte (calc.) & 19.8 \% \\
Densidade aparente & 1.28 g/\allowbreak cm³ \\
\textbf{Aptidão agrícola} & \textbf{Média} \\
\end{longtable}
}

\subsubsection{Indicadores Sociais}

\textbf{Fontes:} PNUD 2020, INE\footnote{\fineINE} Censo 2022, INE\footnote{\fineINE} EPHC 2023, Ministerio
Público 2024\\
\textbf{Nota:} Valores marcados como \emph{(dept.)} referem-se ao
departamento de Concepción; valores \emph{(dist.)} são específicos deste
distrito. Dados marcados \emph{(est.)} são estimativas.

{\def\LTcaptype{none} % do not increment counter
\begin{longtable}[]{@{}
  >{\raggedright\arraybackslash}p{(\linewidth - 4\tabcolsep) * \real{0.4231}}
  >{\raggedright\arraybackslash}p{(\linewidth - 4\tabcolsep) * \real{0.2692}}
  >{\raggedright\arraybackslash}p{(\linewidth - 4\tabcolsep) * \real{0.3077}}@{}}
\toprule\noalign{}
\begin{minipage}[b]{\linewidth}\raggedright
Indicador
\end{minipage} & \begin{minipage}[b]{\linewidth}\raggedright
Valor
\end{minipage} & \begin{minipage}[b]{\linewidth}\raggedright
Âmbito
\end{minipage} \\
\midrule\noalign{}
\endhead
\bottomrule\noalign{}
\endlastfoot
IDH (2020) & 0.640 & dist. (est.) \\
Ranking IDH nacional & 9/\allowbreak 18 & dept. \\
Esperança de vida & 72,4 anos & dept. \\
Escolaridade média & 7.6 anos & dist. \\
RNB per capita (USD PPA) & 8.400 & dept. \\
Pobreza monetária (\%) & 31,1\% & dept. \\
Pobreza extrema (\%) & 7,8\% & dept. \\
Índice de Gini & 0,460 & dept. \\
Acesso a água potável (\%) & 74,2\% & dept. \\
Acesso a saneamento (\%) & 71,8\% & dept. \\
Taxa de homicídios (est., /\allowbreak 100k hab) & \textasciitilde12--15
(estimativa, acima da média nacional) & dept. \\
Índice de segurança & baixo & dept. \\
População (Censo 2022) & 8.292 hab. & dist. \\
Idade mediana & 28.0 anos & dist. \\
Taxa de fecundidade & 2 & dist. \\
\end{longtable}
}
\subsubsection{Combustível}

\textbf{Referência:} PETROPAR /\allowbreak  postos locais (2024) \textbf{Tipo de
localidade:} Interior

{\def\LTcaptype{none} % do not increment counter
\begin{longtable}[]{@{}lll@{}}
\toprule\noalign{}
Combustível & USD/\allowbreak litro & Gs/\allowbreak litro (aprox.) \\
\midrule\noalign{}
\endhead
\bottomrule\noalign{}
\endlastfoot
Gasolina 93 oct & 1.01 & 7,474 \\
Gasolina 97 oct (premium) & 1.11 & 8,214 \\
Diesel & 0.94 & 6,956 \\
\end{longtable}
}

\begin{quote}
Preços podem variar ±5\% conforme posto e sazonalidade. Chaco e interior
remoto apresentam maior variação.
\end{quote}

\subsubsection{Cobertura Celular}

\textbf{Fonte:} CONATEL PY /\allowbreak  operadoras (2024)

{\def\LTcaptype{none} % do not increment counter
\begin{longtable}[]{@{}ll@{}}
\toprule\noalign{}
Parâmetro & Valor \\
\midrule\noalign{}
\endhead
\bottomrule\noalign{}
\endlastfoot
Cobertura 4G população (dept.) & 78\% \\
Cobertura 4G área rural & 55\% \\
Melhor operadora & Tigo \\
Qualidade rural & moderada \\
\end{longtable}
}

\begin{quote}
Para áreas rurais fora do núcleo urbano, recomenda-se chip Tigo como
principal e Personal como backup.
\end{quote}

\subsubsection{Internet}

\textbf{Fonte:} CONATEL /\allowbreak  Speedtest Ookla (2024)

{\def\LTcaptype{none} % do not increment counter
\begin{longtable}[]{@{}ll@{}}
\toprule\noalign{}
Parâmetro & Valor \\
\midrule\noalign{}
\endhead
\bottomrule\noalign{}
\endlastfoot
Velocidade média download & 38 Mbps \\
Domicílios com internet (dept.) & 48\% \\
Tecnologia predominante & rádio \\
Opção rural & Starlink disponível (\textasciitilde USD 44/\allowbreak mês) \\
\end{longtable}
}

\subsubsection{Mercado Imobiliário e Terra
Rural}

\textbf{Fonte:} INDERT /\allowbreak  Clasificados.com.py (2024)

{\def\LTcaptype{none} % do not increment counter
\begin{longtable}[]{@{}ll@{}}
\toprule\noalign{}
Tipo & Referência \\
\midrule\noalign{}
\endhead
\bottomrule\noalign{}
\endlastfoot
Terra agrícola alta prod. (USD/\allowbreak ha) & 3,500 \\
Imóvel urbano (USD/\allowbreak m²) & 650 \\
Aluguel 2 quartos (USD/\allowbreak mês) & 260 \\
\end{longtable}
}

\begin{quote}
Valores de referência departamental. Localidades menores podem ter
preços 20--40\% abaixo da capital departamental.
\end{quote}

\subsubsection{Saúde}

\textbf{Fonte:} MSPBS /\allowbreak  IPS Paraguay (2024)

{\def\LTcaptype{none} % do not increment counter
\begin{longtable}[]{@{}ll@{}}
\toprule\noalign{}
Serviço & Disponibilidade \\
\midrule\noalign{}
\endhead
\bottomrule\noalign{}
\endlastfoot
USF /\allowbreak  Posto de Saúde & sim \\
Hospital Regional & não \\
IPS (seguro social) & não \\
Farmácia & sim (limitada) \\
Distância ao hospital de referência & \textasciitilde78 km
(Concepción) \\
\end{longtable}
}

\textbf{Principais estabelecimentos:} USF Arroyito (atendimento básico
nos núcleos de assentamento); Hospital de Horqueta como referência
intermediária (\textasciitilde30 km)

\textbf{Observação para imigrantes:} Cobertura privada indispensável.
IPS inexistente localmente. Para emergências, deslocamento até Horqueta
ou Concepción obrigatório; recomenda-se manter veículo 4x4 e kit de
primeiros socorros completo.
\newpage
\secaoDiagnostico{Puerto Pinasco}{\mapaDistrito{1503}}{Puerto Pinasco (Presidente Hayes) apresenta perfil operacional
consolidado para comparacao territorial, com leitura conjunta de
infraestrutura, risco hidroclimatico e capacidade de continuidade de
servicos essenciais.

\subsubsection{Por que sim}

\subsubsection{Por que não}

\subsubsection{Pre-condicoes Minimas de
Decisao}

\begin{itemize}
\tightlist
\item
  Atualizar indicadores criticos em ciclo trimestral.
\item
  Confirmar trafegabilidade e contingencia hidrica/\allowbreak energetica local.
\item
  Validar checklist de resiliencia domiciliar/\allowbreak logistica antes de decisao
  final.
\end{itemize}

\subsubsection{Recomendação
Técnica}

Uso recomendado como candidato em analise multicriterio, condicionado ao
cumprimento das pre-condicoes acima. Referência atual de score: GSS 7.4;
risco predominante: moderado.

\begin{itemize}
\tightlist
\item
  \begin{enumerate}
  \def\labelenumi{\arabic{enumi})}
  \tightlist
  \item
    Delimitacao territorial validada por georreferencia institucional.
  \end{enumerate}
\item
  \begin{enumerate}
  \def\labelenumi{\arabic{enumi})}
  \setcounter{enumi}{1}
  \tightlist
  \item
    População municipal/\allowbreak distrital referenciada por base censitaria
    nacional.
  \end{enumerate}
\item
  \begin{enumerate}
  \def\labelenumi{\arabic{enumi})}
  \setcounter{enumi}{2}
  \tightlist
  \item
    Comparabilidade intra-departamental apoiada em indicadores
    distritais oficiais.
  \end{enumerate}
\item
  \begin{enumerate}
  \def\labelenumi{\arabic{enumi})}
  \setcounter{enumi}{3}
  \tightlist
  \item
    Estado macro de conectividade viaria ancorado em informacoes de
    rodovias.
  \end{enumerate}
\item
  \begin{enumerate}
  \def\labelenumi{\arabic{enumi})}
  \setcounter{enumi}{4}
  \tightlist
  \item
    Exposicao hidrometeorologica monitorada por avisos oficiais.
  \end{enumerate}
\item
  \begin{enumerate}
  \def\labelenumi{\arabic{enumi})}
  \setcounter{enumi}{5}
  \tightlist
  \item
    Historico de contexto meteorologico apoiado em publicacoes tecnicas.
  \end{enumerate}
\item
  \begin{enumerate}
  \def\labelenumi{\arabic{enumi})}
  \setcounter{enumi}{6}
  \tightlist
  \item
    Capacidade institucional de resposta a eventos apoiada em acoes da
    SEN\footnote{\fineSEN}.
  \end{enumerate}
\item
  \begin{enumerate}
  \def\labelenumi{\arabic{enumi})}
  \setcounter{enumi}{7}
  \tightlist
  \item
    Infraestrutura energetica analisada via operador nacional.
  \end{enumerate}
\item
  \begin{enumerate}
  \def\labelenumi{\arabic{enumi})}
  \setcounter{enumi}{8}
  \tightlist
  \item
    Referências de obras e logistica nacional verificadas no MOPC.
  \end{enumerate}
\item
  \begin{enumerate}
  \def\labelenumi{\arabic{enumi})}
  \setcounter{enumi}{9}
  \tightlist
  \item
    Contexto sociopolitico institucional referenciado por autoridade
    eleitoral.
  \end{enumerate}
\item
  \begin{enumerate}
  \def\labelenumi{\arabic{enumi})}
  \setcounter{enumi}{10}
  \tightlist
  \item
    Camada cartografica de apoio eleitoral/\allowbreak territorial consultada.
  \end{enumerate}
\item
  \begin{enumerate}
  \def\labelenumi{\arabic{enumi})}
  \setcounter{enumi}{11}
  \tightlist
  \item
    Dados publicos complementares para verificacao cruzada.
  \end{enumerate}
\end{itemize}}{dist:15_Presidente_Hayes:Puerto_Pinasco}
\begin{table}[ht!]
\small \begin{tabular}{p{3.0cm}p{0.8cm}p{6.0cm}} \toprule
\textbf{Critério} & \textbf{Nota} & \textbf{Justificativa} \\ \midrule
A - Ameaças Estratégicas & 7.4 & - \\
B - Risco Social & 7.8 & - \\
C - Riscos Naturais & 6.0 & - \\
D - Autossuficiência & 7.6 & - \\
E - Institucional & 7.8 & - \\
\midrule \textbf{GSS FINAL} & \textbf{7.4} & \textbf{Seguro (bom potencial com preparacao).} \\ \bottomrule \end{tabular}
\end{table}
\subsubsection{Vulnerabilidades e
Mitigação}

\begin{itemize}
\tightlist
\item
  \textbf{Conflitos Armados:} Hotspot de guerrilha rural e insurgência
  (Mitigação: segurança privada armada em estâncias, residência em áreas
  com proteção da FTC, planos de evacuação aérea/\allowbreak fluvial).
\item
  \textbf{Insegurança Jurídica/\allowbreak Social:} Assentamentos com histórico de
  conflitos agrários (Mitigação: aquisição apenas de terras com título
  definitivo registrado no RGP).
\item
  \textbf{Gargalo Logístico:} Controle total do Cruce Azotey pode isolar
  a movimentação Norte-Sul pela PY08.
\end{itemize}

\subsubsection{Dossiê de Campo}
\begin{description}[style=nextline,leftmargin=0.5cm,font=\bfseries]
\item[Coordenadas]  23°13'S 56°30'W.
\item[Topografia]  Relevo predominantemente plano com leves ondulações. Altitude \textasciitilde 180m. Geologicamente estável, risco sísmico muito baixo.
\item[Alvos Estratégicos]  Cruce Azotey (Entroncamento tático na Ruta PY08); Base de operações da FTC (Força de Tarefa Conjunta); Proximidade com grandes estâncias pecuárias (ex: Paso Itá).
\item[Fallout]  Ventos predominantes do Norte/\allowbreak Nordeste.
\item[População]  \textasciitilde 9.100 habitantes (Estimativa baseada no Censo 2022). Perfil majoritariamente rural.
\item[Densidade]  \textasciitilde 12,2 hab/\allowbreak km² (Baixa densidade, mas com concentração em núcleos de assentamento como Zanja Morotí).
\item[Vias de Acesso]  Ruta PY08 (eixo asfáltico fundamental para o escoamento de grãos e gado). Estradas vicinais de terra com trafegabilidade crítica em dias de chuva.
\item[Serviços]  Unidade de Saúde da Família (USF) para atenção básica. Casos graves dependem de Santa Rosa del Aguaray (San Pedro) ou Concepción (140 km).
\item[Custo de Vida]  Baixo; economia estritamente vinculada ao agronegócio e pequena agricultura.
\item[Preço da Terra]  US\$ 3.500-5.000/\allowbreak ha (áreas mecanizáveis ou pastagens de alta produtividade).
\item[Hidrologia]  Risco de isolamento logístico por transbordamento de pequenos tributários e degradação das vias de terra.
\item[Clima]  Tropical Úmido. Média de 1.300 a 1.600 mm de chuva anuais. Verões intensos com tempestades súbitas.
\item[Energia]  Rede nacional (ANDE\footnote{\fineANDE}), com vulnerabilidades em áreas rurais remotas.
\item[Água]  Poços artesianos; lençol freático acessível.
\item[Qualidade do Solo]  Alfisois e Ultisois; fertilidade média, exige calagem e adubação para agricultura intensiva (soja/\allowbreak milho).
\item[Recursos Locais]  Forte produção de gado de corte e grãos. Elevado potencial para autossuficiência alimentar básica em nível de estância ou comunidade protegida.
\item[Segurança]  Zona de Risco Crítico. Azotey está no epicentro da área de influência do EPP\footnote{\fineEPP} (Ejército del Pueblo Paraguayo). Presença militar permanente da FTC é a norma de segurança. Taxa de homicídios regional entre as maiores do país (\textasciitilde 15/\allowbreak 100k).
\item[Leis Local: Município criado em 2009. Restrição de Fronteira]  Zona limítrofe à faixa de 50km da Lei 2532/\allowbreak 05; exige verificação detalhada da localização das glebas.
\end{description}
\textbf{Fonte climática:} NASA POWER Climatology API (período 2001-2020)

\textbf{Fonte luminosa:} estimativa world atlas

\paragraph{Irradiação Solar
(kWh/\allowbreak m²/\allowbreak dia)}

{\def\LTcaptype{none} % do not increment counter
\begin{longtable}[]{@{}
  >{\raggedright\arraybackslash}p{(\linewidth - 24\tabcolsep) * \real{0.0746}}
  >{\raggedright\arraybackslash}p{(\linewidth - 24\tabcolsep) * \real{0.0746}}
  >{\raggedright\arraybackslash}p{(\linewidth - 24\tabcolsep) * \real{0.0746}}
  >{\raggedright\arraybackslash}p{(\linewidth - 24\tabcolsep) * \real{0.0746}}
  >{\raggedright\arraybackslash}p{(\linewidth - 24\tabcolsep) * \real{0.0746}}
  >{\raggedright\arraybackslash}p{(\linewidth - 24\tabcolsep) * \real{0.0746}}
  >{\raggedright\arraybackslash}p{(\linewidth - 24\tabcolsep) * \real{0.0746}}
  >{\raggedright\arraybackslash}p{(\linewidth - 24\tabcolsep) * \real{0.0746}}
  >{\raggedright\arraybackslash}p{(\linewidth - 24\tabcolsep) * \real{0.0746}}
  >{\raggedright\arraybackslash}p{(\linewidth - 24\tabcolsep) * \real{0.0746}}
  >{\raggedright\arraybackslash}p{(\linewidth - 24\tabcolsep) * \real{0.0746}}
  >{\raggedright\arraybackslash}p{(\linewidth - 24\tabcolsep) * \real{0.0746}}
  >{\raggedright\arraybackslash}p{(\linewidth - 24\tabcolsep) * \real{0.1045}}@{}}
\toprule\noalign{}
\begin{minipage}[b]{\linewidth}\raggedright
Jan
\end{minipage} & \begin{minipage}[b]{\linewidth}\raggedright
Fev
\end{minipage} & \begin{minipage}[b]{\linewidth}\raggedright
Mar
\end{minipage} & \begin{minipage}[b]{\linewidth}\raggedright
Abr
\end{minipage} & \begin{minipage}[b]{\linewidth}\raggedright
Mai
\end{minipage} & \begin{minipage}[b]{\linewidth}\raggedright
Jun
\end{minipage} & \begin{minipage}[b]{\linewidth}\raggedright
Jul
\end{minipage} & \begin{minipage}[b]{\linewidth}\raggedright
Ago
\end{minipage} & \begin{minipage}[b]{\linewidth}\raggedright
Set
\end{minipage} & \begin{minipage}[b]{\linewidth}\raggedright
Out
\end{minipage} & \begin{minipage}[b]{\linewidth}\raggedright
Nov
\end{minipage} & \begin{minipage}[b]{\linewidth}\raggedright
Dez
\end{minipage} & \begin{minipage}[b]{\linewidth}\raggedright
Média
\end{minipage} \\
\midrule\noalign{}
\endhead
\bottomrule\noalign{}
\endlastfoot
6.60 & 6.02 & 5.62 & 4.65 & 3.48 & 3.07 & 3.35 & 4.15 & 4.72 & 5.46 &
6.29 & 6.53 & \textbf{4.99} \\
\end{longtable}
}

\textbf{Inclinação solar recomendada:} 23° N (anual) · 33° N (inverno
jun-ago) · 13° N (verão nov-jan)

\paragraph{Precipitação (mm/\allowbreak mês)}

{\def\LTcaptype{none} % do not increment counter
\begin{longtable}[]{@{}
  >{\raggedright\arraybackslash}p{(\linewidth - 24\tabcolsep) * \real{0.0704}}
  >{\raggedright\arraybackslash}p{(\linewidth - 24\tabcolsep) * \real{0.0704}}
  >{\raggedright\arraybackslash}p{(\linewidth - 24\tabcolsep) * \real{0.0704}}
  >{\raggedright\arraybackslash}p{(\linewidth - 24\tabcolsep) * \real{0.0704}}
  >{\raggedright\arraybackslash}p{(\linewidth - 24\tabcolsep) * \real{0.0704}}
  >{\raggedright\arraybackslash}p{(\linewidth - 24\tabcolsep) * \real{0.0704}}
  >{\raggedright\arraybackslash}p{(\linewidth - 24\tabcolsep) * \real{0.0704}}
  >{\raggedright\arraybackslash}p{(\linewidth - 24\tabcolsep) * \real{0.0704}}
  >{\raggedright\arraybackslash}p{(\linewidth - 24\tabcolsep) * \real{0.0704}}
  >{\raggedright\arraybackslash}p{(\linewidth - 24\tabcolsep) * \real{0.0704}}
  >{\raggedright\arraybackslash}p{(\linewidth - 24\tabcolsep) * \real{0.0704}}
  >{\raggedright\arraybackslash}p{(\linewidth - 24\tabcolsep) * \real{0.0704}}
  >{\raggedright\arraybackslash}p{(\linewidth - 24\tabcolsep) * \real{0.1549}}@{}}
\toprule\noalign{}
\begin{minipage}[b]{\linewidth}\raggedright
Jan
\end{minipage} & \begin{minipage}[b]{\linewidth}\raggedright
Fev
\end{minipage} & \begin{minipage}[b]{\linewidth}\raggedright
Mar
\end{minipage} & \begin{minipage}[b]{\linewidth}\raggedright
Abr
\end{minipage} & \begin{minipage}[b]{\linewidth}\raggedright
Mai
\end{minipage} & \begin{minipage}[b]{\linewidth}\raggedright
Jun
\end{minipage} & \begin{minipage}[b]{\linewidth}\raggedright
Jul
\end{minipage} & \begin{minipage}[b]{\linewidth}\raggedright
Ago
\end{minipage} & \begin{minipage}[b]{\linewidth}\raggedright
Set
\end{minipage} & \begin{minipage}[b]{\linewidth}\raggedright
Out
\end{minipage} & \begin{minipage}[b]{\linewidth}\raggedright
Nov
\end{minipage} & \begin{minipage}[b]{\linewidth}\raggedright
Dez
\end{minipage} & \begin{minipage}[b]{\linewidth}\raggedright
Total/\allowbreak ano
\end{minipage} \\
\midrule\noalign{}
\endhead
\bottomrule\noalign{}
\endlastfoot
156 & 171 & 120 & 138 & 129 & 81 & 54 & 43 & 85 & 166 & 195 & 170 &
\textbf{1513 mm} \\
\end{longtable}
}

\paragraph{Poluição Luminosa}

{\def\LTcaptype{none} % do not increment counter
\begin{longtable}[]{@{}ll@{}}
\toprule\noalign{}
Parâmetro & Valor \\
\midrule\noalign{}
\endhead
\bottomrule\noalign{}
\endlastfoot
Escala Bortle & 3 --- Céu rural \\
Radiância artificial & 1.5 nW/\allowbreak cm²/\allowbreak sr \\
\end{longtable}
}
\paragraph{Solo (SoilGrids 2.0, média ponderada 0--30
cm)}

{\def\LTcaptype{none} % do not increment counter
\begin{longtable}[]{@{}ll@{}}
\toprule\noalign{}
Parâmetro & Valor \\
\midrule\noalign{}
\endhead
\bottomrule\noalign{}
\endlastfoot
pH (H$_2$O) & 5.42 \\
Carbono orgânico (SOC) & 14.85 g/\allowbreak kg \\
Argila & 20.56 \% \\
Areia & 58.92 \% \\
Silte (calc.) & 20.5 \% \\
Densidade aparente & 1.33 g/\allowbreak cm³ \\
\textbf{Aptidão agrícola} & \textbf{Média} \\
\end{longtable}
}

\subsubsection{Indicadores Sociais}

\textbf{Fontes:} PNUD 2020, INE\footnote{\fineINE} Censo 2022, INE\footnote{\fineINE} EPHC 2023, Ministerio
Público 2024\\
\textbf{Nota:} Valores marcados como \emph{(dept.)} referem-se ao
departamento de Concepción; valores \emph{(dist.)} são específicos deste
distrito. Dados marcados \emph{(est.)} são estimativas.

{\def\LTcaptype{none} % do not increment counter
\begin{longtable}[]{@{}
  >{\raggedright\arraybackslash}p{(\linewidth - 4\tabcolsep) * \real{0.4231}}
  >{\raggedright\arraybackslash}p{(\linewidth - 4\tabcolsep) * \real{0.2692}}
  >{\raggedright\arraybackslash}p{(\linewidth - 4\tabcolsep) * \real{0.3077}}@{}}
\toprule\noalign{}
\begin{minipage}[b]{\linewidth}\raggedright
Indicador
\end{minipage} & \begin{minipage}[b]{\linewidth}\raggedright
Valor
\end{minipage} & \begin{minipage}[b]{\linewidth}\raggedright
Âmbito
\end{minipage} \\
\midrule\noalign{}
\endhead
\bottomrule\noalign{}
\endlastfoot
IDH (2020) & 0.635 & dist. (est.) \\
Ranking IDH nacional & 9/\allowbreak 18 & dept. \\
Esperança de vida & 72,4 anos & dept. \\
Escolaridade média & 7 & dist. \\
RNB per capita (USD PPA) & 8.400 & dept. \\
Pobreza monetária (\%) & 31,1\% & dept. \\
Pobreza extrema (\%) & 7,8\% & dept. \\
Índice de Gini & 0,460 & dept. \\
Acesso a água potável (\%) & 74,2\% & dept. \\
Acesso a saneamento (\%) & 71,8\% & dept. \\
Taxa de homicídios (est., /\allowbreak 100k hab) & \textasciitilde12--15
(estimativa, acima da média nacional) & dept. \\
Índice de segurança & baixo & dept. \\
População (Censo 2022) & \textasciitilde9.100 habitantes & dist. \\
Idade mediana & 22 anos & dist. \\
Taxa de fecundidade & 3 & dist. \\
\end{longtable}
}
\subsubsection{Combustível}

\textbf{Referência:} PETROPAR /\allowbreak  postos locais (2024) \textbf{Tipo de
localidade:} Interior

{\def\LTcaptype{none} % do not increment counter
\begin{longtable}[]{@{}lll@{}}
\toprule\noalign{}
Combustível & USD/\allowbreak litro & Gs/\allowbreak litro (aprox.) \\
\midrule\noalign{}
\endhead
\bottomrule\noalign{}
\endlastfoot
Gasolina 93 oct & 1.01 & 7,474 \\
Gasolina 97 oct (premium) & 1.11 & 8,214 \\
Diesel & 0.94 & 6,956 \\
\end{longtable}
}

\begin{quote}
Preços podem variar ±5\% conforme posto e sazonalidade. Chaco e interior
remoto apresentam maior variação.
\end{quote}

\subsubsection{Cobertura Celular}

\textbf{Fonte:} CONATEL PY /\allowbreak  operadoras (2024)

{\def\LTcaptype{none} % do not increment counter
\begin{longtable}[]{@{}ll@{}}
\toprule\noalign{}
Parâmetro & Valor \\
\midrule\noalign{}
\endhead
\bottomrule\noalign{}
\endlastfoot
Cobertura 4G população (dept.) & 78\% \\
Cobertura 4G área rural & 55\% \\
Melhor operadora & Tigo \\
Qualidade rural & moderada \\
\end{longtable}
}

\begin{quote}
Para áreas rurais fora do núcleo urbano, recomenda-se chip Tigo como
principal e Personal como backup.
\end{quote}

\subsubsection{Internet}

\textbf{Fonte:} CONATEL /\allowbreak  Speedtest Ookla (2024)

{\def\LTcaptype{none} % do not increment counter
\begin{longtable}[]{@{}ll@{}}
\toprule\noalign{}
Parâmetro & Valor \\
\midrule\noalign{}
\endhead
\bottomrule\noalign{}
\endlastfoot
Velocidade média download & 38 Mbps \\
Domicílios com internet (dept.) & 48\% \\
Tecnologia predominante & rádio \\
Opção rural & Starlink disponível (\textasciitilde USD 44/\allowbreak mês) \\
\end{longtable}
}

\subsubsection{Mercado Imobiliário e Terra
Rural}

\textbf{Fonte:} INDERT /\allowbreak  Clasificados.com.py (2024)

{\def\LTcaptype{none} % do not increment counter
\begin{longtable}[]{@{}ll@{}}
\toprule\noalign{}
Tipo & Referência \\
\midrule\noalign{}
\endhead
\bottomrule\noalign{}
\endlastfoot
Terra agrícola alta prod. (USD/\allowbreak ha) & 3,500 \\
Imóvel urbano (USD/\allowbreak m²) & 650 \\
Aluguel 2 quartos (USD/\allowbreak mês) & 260 \\
\end{longtable}
}

\begin{quote}
Valores de referência departamental. Localidades menores podem ter
preços 20--40\% abaixo da capital departamental.
\end{quote}

\subsubsection{Saúde}

\textbf{Fonte:} MSPBS /\allowbreak  IPS Paraguay (2024-2026), consolidação
departamental e proxy local

{\def\LTcaptype{none} % do not increment counter
\begin{longtable}[]{@{}ll@{}}
\toprule\noalign{}
Serviço & Disponibilidade \\
\midrule\noalign{}
\endhead
\bottomrule\noalign{}
\endlastfoot
USF /\allowbreak  Posto de Saúde & sim \\
Hospital Regional & não \\
IPS (seguro social) & sim \\
Farmácia & sim \\
Distância ao hospital de referência & \textasciitilde122 km
(Concepción) \\
\end{longtable}
}

\textbf{Principais estabelecimentos:} USF local; referencia hospitalar
em Concepción

\textbf{Observação para imigrantes:} Atendimento primario local ou em
raio curto; casos de maior complexidade seguem para Concepción.
Cobertura privada continua recomendada para especialidades e urgências
de maior complexidade.
\newpage
\secaoDiagnostico{Teniente Esteban Martinez}{\mapaDistrito{1508}}{Teniente Esteban Martínez é a fronteira final da privacidade no
Paraguai. Oferece um nível de isolamento geográfico e demográfico que
transforma o distrito em uma zona virtualmente invisível a crises
globais ou nacionais. Sua força absoluta reside no anonimato garantido
pela imensidão do território e pelo vazio populacional. Contudo, essa
mesma característica impõe desafios brutais de subsistência: a gestão da
água salobra e o isolamento físico imposto pelas águas do Pilcomayo.
Sobreviver aqui exige um perfil de ``pioneiro tático'', capaz de gerir
sua própria energia, saúde e logística com zero expectativa de auxílio
externo, transformando a dureza do Chaco em uma armadura de proteção
definitiva.

\subsubsection{Por que sim}

\subsubsection{Por que não}

\subsubsection{Recomendação
Técnica}

Indicado exclusivamente para estabelecimentos de isolamento tático
extremo que possuam infraestrutura própria de dessalinização, energia
solar e meios de evacuação (pista de pouso). O GSS de 6.3 reflete a
superioridade do isolamento protegida pela hostilidade natural da
região. Referência atual de score: GSS 6.3.

\begin{itemize}
\tightlist
\item
  \begin{enumerate}
  \def\labelenumi{\arabic{enumi})}
  \tightlist
  \item
    Dados censitários validados (INE\footnote{\fineINE} 2022).
  \end{enumerate}
\item
  \begin{enumerate}
  \def\labelenumi{\arabic{enumi})}
  \setcounter{enumi}{1}
  \tightlist
  \item
    Registros de inundações e vazões do Rio Pilcomayo (SEN\footnote{\fineSEN}/\allowbreak DMH).
  \end{enumerate}
\item
  \begin{enumerate}
  \def\labelenumi{\arabic{enumi})}
  \setcounter{enumi}{2}
  \tightlist
  \item
    Mapa de solos e bacias hidrográficas do Chaco (MAG).
  \end{enumerate}
\item
  \begin{enumerate}
  \def\labelenumi{\arabic{enumi})}
  \setcounter{enumi}{3}
  \tightlist
  \item
    Registro de segurança pública departamental.
  \end{enumerate}
\end{itemize}}{dist:15_Presidente_Hayes:Teniente_Esteban_Martinez}
\begin{table}[ht!]
\small \begin{tabular}{p{3.0cm}p{0.8cm}p{6.0cm}} \toprule
\textbf{Critério} & \textbf{Nota} & \textbf{Justificativa} \\ \midrule
A - Ameaças Estratégicas & 8.0 & Isolamento tático absoluto; ausência total de alvos e interesse estratégico militar \\
B - Risco Social & 9.0 & Densidade populacional ínfima; segurança contra riscos de massa absoluta \\
C - Riscos Naturais & 3.5 & Risco hidrológico e climático extremo; nota penalizada pela instabilidade do Pilcomayo \\
D - Autossuficiência & 6.0 & Vastidão de recursos de proteína animal; severamente limitado pela qualidade hídrica e logística \\
E - Institucional & 4.0 & Ambiente social seguro pelo isolamento, mas com suporte institucional quase inexistente \\
\midrule \textbf{GSS FINAL} & \textbf{6.3} & \textbf{Moderadamente Seguro (Indicado apenas para projetos de isolamento tático extremo "off-grid" que possuam autonomia total de recursos e meios de transporte aéreo/\allowbreak pesado).} \\ \bottomrule \end{tabular}
\end{table}
\subsubsection{Vulnerabilidades e
Mitigação}

\begin{itemize}
\tightlist
\item
  \textbf{Dependência Industrial:} Economia local muito vinculada ao
  frigorífico (Mitigação: diversificação para agricultura de
  subsistência e energia solar).
\item
  \textbf{Instabilidade Climática:} Temperaturas de verão exigem
  construção com alta inércia térmica e sistemas de resfriamento
  eficientes.
\item
  \textbf{Segurança Regional:} Monitoramento de movimentos insurgentes
  em distritos vizinhos (Mitigação: residência próxima ao perímetro
  urbano e comunicação satelital).
\end{itemize}

\subsubsection{Dossiê de Campo}
\begin{description}[style=nextline,leftmargin=0.5cm,font=\bfseries]
\item[Coordenadas]  23°27'S 57°15'W (Ponto exato de passagem do Trópico de Capricórnio).
\item[Topografia]  Relevo predominantemente plano com áreas de várzea próximas ao Rio Ypané. Altitude \textasciitilde 80m. Geologicamente estável, risco sísmico desprezível.
\item[Alvos Estratégicos]  Frigorífico Athena Foods (ex-JBS) - motor econômico e polo de suprimento proteico; Cruzamento de rotas para Concepción e San Pedro.
\item[Fallout]  Ventos predominantes do Norte, Leste e Sudeste.
\item[População]  10.605 habitantes (Censo 2022 Final).
\item[Densidade]  \textasciitilde 37,2 hab/\allowbreak km² (Moderada, com núcleo urbano consolidado e histórico).
\item[Vias de Acesso]  Conectividade asfáltica eficiente com Concepción (21 km). Ponto de trânsito regional de carga pesada (gado e grãos).
\item[Serviços]  Centro de saúde local para urgências básicas. Dependência direta do Hospital Regional de Concepción para alta complexidade.
\item[Custo de Vida]  Significativamente baixo (30-40\% menor que o Brasil para alimentação).
\item[Preço da Terra]  US\$ 4.000-10.000/\allowbreak ha (terras de alta produtividade agrícola/\allowbreak pastagem); US\$ 2.500-3.500/\allowbreak ha (áreas de reflorestamento).
\item[Hidrologia]  Risco Moderado de Inundação. Presença do Rio Ypané; áreas ribeirinhas vulneráveis em anos de El Niño intenso. Praias fluviais são ativos de lazer.
\item[Clima]  Tropical. Temperaturas extremas (máximas de 45°C e geadas ocasionais no inverno). Chuvas concentradas entre novembro e janeiro.
\item[Energia]  Rede nacional (Itaipu). Estabilidade moderada para o núcleo industrial.
\item[Água]  Abundante (Rio Ypané e poços artesianos).
\item[Qualidade do Solo]  Solos podzólicos e lateríticos; profundos e ácidos, respondem muito bem a fertilizantes. Ideais para pastagens, erva-mate e fruticultura.
\item[Recursos Locais]  Grande excedente de proteína animal (frigorífico local). Elevadíssimo potencial para autossuficiência e suprimento de longo prazo.
\end{description}
\textbf{Fonte climática:} NASA POWER Climatology API (período 2001-2020)

\textbf{Fonte luminosa:} estimativa world atlas

\paragraph{Irradiação Solar
(kWh/\allowbreak m²/\allowbreak dia)}

{\def\LTcaptype{none} % do not increment counter
\begin{longtable}[]{@{}
  >{\raggedright\arraybackslash}p{(\linewidth - 24\tabcolsep) * \real{0.0746}}
  >{\raggedright\arraybackslash}p{(\linewidth - 24\tabcolsep) * \real{0.0746}}
  >{\raggedright\arraybackslash}p{(\linewidth - 24\tabcolsep) * \real{0.0746}}
  >{\raggedright\arraybackslash}p{(\linewidth - 24\tabcolsep) * \real{0.0746}}
  >{\raggedright\arraybackslash}p{(\linewidth - 24\tabcolsep) * \real{0.0746}}
  >{\raggedright\arraybackslash}p{(\linewidth - 24\tabcolsep) * \real{0.0746}}
  >{\raggedright\arraybackslash}p{(\linewidth - 24\tabcolsep) * \real{0.0746}}
  >{\raggedright\arraybackslash}p{(\linewidth - 24\tabcolsep) * \real{0.0746}}
  >{\raggedright\arraybackslash}p{(\linewidth - 24\tabcolsep) * \real{0.0746}}
  >{\raggedright\arraybackslash}p{(\linewidth - 24\tabcolsep) * \real{0.0746}}
  >{\raggedright\arraybackslash}p{(\linewidth - 24\tabcolsep) * \real{0.0746}}
  >{\raggedright\arraybackslash}p{(\linewidth - 24\tabcolsep) * \real{0.0746}}
  >{\raggedright\arraybackslash}p{(\linewidth - 24\tabcolsep) * \real{0.1045}}@{}}
\toprule\noalign{}
\begin{minipage}[b]{\linewidth}\raggedright
Jan
\end{minipage} & \begin{minipage}[b]{\linewidth}\raggedright
Fev
\end{minipage} & \begin{minipage}[b]{\linewidth}\raggedright
Mar
\end{minipage} & \begin{minipage}[b]{\linewidth}\raggedright
Abr
\end{minipage} & \begin{minipage}[b]{\linewidth}\raggedright
Mai
\end{minipage} & \begin{minipage}[b]{\linewidth}\raggedright
Jun
\end{minipage} & \begin{minipage}[b]{\linewidth}\raggedright
Jul
\end{minipage} & \begin{minipage}[b]{\linewidth}\raggedright
Ago
\end{minipage} & \begin{minipage}[b]{\linewidth}\raggedright
Set
\end{minipage} & \begin{minipage}[b]{\linewidth}\raggedright
Out
\end{minipage} & \begin{minipage}[b]{\linewidth}\raggedright
Nov
\end{minipage} & \begin{minipage}[b]{\linewidth}\raggedright
Dez
\end{minipage} & \begin{minipage}[b]{\linewidth}\raggedright
Média
\end{minipage} \\
\midrule\noalign{}
\endhead
\bottomrule\noalign{}
\endlastfoot
6.68 & 6.07 & 5.58 & 4.63 & 3.47 & 3.00 & 3.30 & 4.09 & 4.68 & 5.49 &
6.34 & 6.61 & \textbf{5.0} \\
\end{longtable}
}

\textbf{Inclinação solar recomendada:} 23° N (anual) · 33° N (inverno
jun-ago) · 13° N (verão nov-jan)

\paragraph{Precipitação (mm/\allowbreak mês)}

{\def\LTcaptype{none} % do not increment counter
\begin{longtable}[]{@{}
  >{\raggedright\arraybackslash}p{(\linewidth - 24\tabcolsep) * \real{0.0704}}
  >{\raggedright\arraybackslash}p{(\linewidth - 24\tabcolsep) * \real{0.0704}}
  >{\raggedright\arraybackslash}p{(\linewidth - 24\tabcolsep) * \real{0.0704}}
  >{\raggedright\arraybackslash}p{(\linewidth - 24\tabcolsep) * \real{0.0704}}
  >{\raggedright\arraybackslash}p{(\linewidth - 24\tabcolsep) * \real{0.0704}}
  >{\raggedright\arraybackslash}p{(\linewidth - 24\tabcolsep) * \real{0.0704}}
  >{\raggedright\arraybackslash}p{(\linewidth - 24\tabcolsep) * \real{0.0704}}
  >{\raggedright\arraybackslash}p{(\linewidth - 24\tabcolsep) * \real{0.0704}}
  >{\raggedright\arraybackslash}p{(\linewidth - 24\tabcolsep) * \real{0.0704}}
  >{\raggedright\arraybackslash}p{(\linewidth - 24\tabcolsep) * \real{0.0704}}
  >{\raggedright\arraybackslash}p{(\linewidth - 24\tabcolsep) * \real{0.0704}}
  >{\raggedright\arraybackslash}p{(\linewidth - 24\tabcolsep) * \real{0.0704}}
  >{\raggedright\arraybackslash}p{(\linewidth - 24\tabcolsep) * \real{0.1549}}@{}}
\toprule\noalign{}
\begin{minipage}[b]{\linewidth}\raggedright
Jan
\end{minipage} & \begin{minipage}[b]{\linewidth}\raggedright
Fev
\end{minipage} & \begin{minipage}[b]{\linewidth}\raggedright
Mar
\end{minipage} & \begin{minipage}[b]{\linewidth}\raggedright
Abr
\end{minipage} & \begin{minipage}[b]{\linewidth}\raggedright
Mai
\end{minipage} & \begin{minipage}[b]{\linewidth}\raggedright
Jun
\end{minipage} & \begin{minipage}[b]{\linewidth}\raggedright
Jul
\end{minipage} & \begin{minipage}[b]{\linewidth}\raggedright
Ago
\end{minipage} & \begin{minipage}[b]{\linewidth}\raggedright
Set
\end{minipage} & \begin{minipage}[b]{\linewidth}\raggedright
Out
\end{minipage} & \begin{minipage}[b]{\linewidth}\raggedright
Nov
\end{minipage} & \begin{minipage}[b]{\linewidth}\raggedright
Dez
\end{minipage} & \begin{minipage}[b]{\linewidth}\raggedright
Total/\allowbreak ano
\end{minipage} \\
\midrule\noalign{}
\endhead
\bottomrule\noalign{}
\endlastfoot
157 & 151 & 120 & 139 & 116 & 71 & 48 & 31 & 73 & 151 & 204 & 167 &
\textbf{1432 mm} \\
\end{longtable}
}

\paragraph{Poluição Luminosa}

{\def\LTcaptype{none} % do not increment counter
\begin{longtable}[]{@{}ll@{}}
\toprule\noalign{}
Parâmetro & Valor \\
\midrule\noalign{}
\endhead
\bottomrule\noalign{}
\endlastfoot
Escala Bortle & 3 --- Céu rural \\
Radiância artificial & 1.5 nW/\allowbreak cm²/\allowbreak sr \\
\end{longtable}
}
\paragraph{Solo (SoilGrids 2.0, média ponderada 0--30
cm)}

{\def\LTcaptype{none} % do not increment counter
\begin{longtable}[]{@{}ll@{}}
\toprule\noalign{}
Parâmetro & Valor \\
\midrule\noalign{}
\endhead
\bottomrule\noalign{}
\endlastfoot
pH (H$_2$O) & 5.7 \\
Carbono orgânico (SOC) & 15.9 g/\allowbreak kg \\
Argila & 16.88 \% \\
Areia & 60.01 \% \\
Silte (calc.) & 23.1 \% \\
Densidade aparente & 1.32 g/\allowbreak cm³ \\
\textbf{Aptidão agrícola} & \textbf{Média} \\
\end{longtable}
}

\subsubsection{Indicadores Sociais}

\textbf{Fontes:} PNUD 2020, INE\footnote{\fineINE} Censo 2022, INE\footnote{\fineINE} EPHC 2023, Ministerio
Público 2024\\
\textbf{Nota:} Valores marcados como \emph{(dept.)} referem-se ao
departamento de Concepción; valores \emph{(dist.)} são específicos deste
distrito. Dados marcados \emph{(est.)} são estimativas.

{\def\LTcaptype{none} % do not increment counter
\begin{longtable}[]{@{}
  >{\raggedright\arraybackslash}p{(\linewidth - 4\tabcolsep) * \real{0.4231}}
  >{\raggedright\arraybackslash}p{(\linewidth - 4\tabcolsep) * \real{0.2692}}
  >{\raggedright\arraybackslash}p{(\linewidth - 4\tabcolsep) * \real{0.3077}}@{}}
\toprule\noalign{}
\begin{minipage}[b]{\linewidth}\raggedright
Indicador
\end{minipage} & \begin{minipage}[b]{\linewidth}\raggedright
Valor
\end{minipage} & \begin{minipage}[b]{\linewidth}\raggedright
Âmbito
\end{minipage} \\
\midrule\noalign{}
\endhead
\bottomrule\noalign{}
\endlastfoot
IDH (2020) & 0.655 & dist. (est.) \\
Ranking IDH nacional & 9/\allowbreak 18 & dept. \\
Esperança de vida & 72,4 anos & dept. \\
Escolaridade média & 8.6 anos & dist. \\
RNB per capita (USD PPA) & 8.400 & dept. \\
Pobreza monetária (\%) & 31,1\% & dept. \\
Pobreza extrema (\%) & 7,8\% & dept. \\
Índice de Gini & 0,460 & dept. \\
Acesso a água potável (\%) & 74,2\% & dept. \\
Acesso a saneamento (\%) & 71,8\% & dept. \\
Taxa de homicídios (est., /\allowbreak 100k hab) & \textasciitilde12--15
(estimativa, acima da média nacional) & dept. \\
Índice de segurança & baixo & dept. \\
População (Censo 2022) & 10.605 hab. & dist. \\
Idade mediana & 28.0 anos & dist. \\
Taxa de fecundidade & 2 & dist. \\
\end{longtable}
}
\subsubsection{Combustível}

\textbf{Referência:} PETROPAR /\allowbreak  postos locais (2024) \textbf{Tipo de
localidade:} Interior

{\def\LTcaptype{none} % do not increment counter
\begin{longtable}[]{@{}lll@{}}
\toprule\noalign{}
Combustível & USD/\allowbreak litro & Gs/\allowbreak litro (aprox.) \\
\midrule\noalign{}
\endhead
\bottomrule\noalign{}
\endlastfoot
Gasolina 93 oct & 1.01 & 7,474 \\
Gasolina 97 oct (premium) & 1.11 & 8,214 \\
Diesel & 0.94 & 6,956 \\
\end{longtable}
}

\begin{quote}
Preços podem variar ±5\% conforme posto e sazonalidade. Chaco e interior
remoto apresentam maior variação.
\end{quote}

\subsubsection{Cobertura Celular}

\textbf{Fonte:} CONATEL PY /\allowbreak  operadoras (2024)

{\def\LTcaptype{none} % do not increment counter
\begin{longtable}[]{@{}ll@{}}
\toprule\noalign{}
Parâmetro & Valor \\
\midrule\noalign{}
\endhead
\bottomrule\noalign{}
\endlastfoot
Cobertura 4G população (dept.) & 78\% \\
Cobertura 4G área rural & 55\% \\
Melhor operadora & Tigo \\
Qualidade rural & moderada \\
\end{longtable}
}

\begin{quote}
Para áreas rurais fora do núcleo urbano, recomenda-se chip Tigo como
principal e Personal como backup.
\end{quote}

\subsubsection{Internet}

\textbf{Fonte:} CONATEL /\allowbreak  Speedtest Ookla (2024)

{\def\LTcaptype{none} % do not increment counter
\begin{longtable}[]{@{}ll@{}}
\toprule\noalign{}
Parâmetro & Valor \\
\midrule\noalign{}
\endhead
\bottomrule\noalign{}
\endlastfoot
Velocidade média download & 38 Mbps \\
Domicílios com internet (dept.) & 48\% \\
Tecnologia predominante & rádio \\
Opção rural & Starlink disponível (\textasciitilde USD 44/\allowbreak mês) \\
\end{longtable}
}

\subsubsection{Mercado Imobiliário e Terra
Rural}

\textbf{Fonte:} INDERT /\allowbreak  Clasificados.com.py (2024)

{\def\LTcaptype{none} % do not increment counter
\begin{longtable}[]{@{}ll@{}}
\toprule\noalign{}
Tipo & Referência \\
\midrule\noalign{}
\endhead
\bottomrule\noalign{}
\endlastfoot
Terra agrícola alta prod. (USD/\allowbreak ha) & 3,500 \\
Imóvel urbano (USD/\allowbreak m²) & 650 \\
Aluguel 2 quartos (USD/\allowbreak mês) & 260 \\
\end{longtable}
}

\begin{quote}
Valores de referência departamental. Localidades menores podem ter
preços 20--40\% abaixo da capital departamental.
\end{quote}

\subsubsection{Saúde}

\textbf{Fonte:} MSPBS /\allowbreak  IPS Paraguay (2024-2026), consolidação
departamental e proxy local

{\def\LTcaptype{none} % do not increment counter
\begin{longtable}[]{@{}ll@{}}
\toprule\noalign{}
Serviço & Disponibilidade \\
\midrule\noalign{}
\endhead
\bottomrule\noalign{}
\endlastfoot
USF /\allowbreak  Posto de Saúde & sim \\
Hospital Regional & sim \\
IPS (seguro social) & sim \\
Farmácia & sim \\
Distância ao hospital de referência & \textasciitilde24 km
(Concepción) \\
\end{longtable}
}

\textbf{Principais estabelecimentos:} USF local; referencia hospitalar
em Concepción

\textbf{Observação para imigrantes:} Atendimento primario local ou em
raio curto; casos de maior complexidade seguem para Concepción.
Cobertura privada continua recomendada para especialidades e urgências
de maior complexidade.
\newpage
\secaoDiagnostico{Teniente Irala Fernandez}{\mapaConteudo{-22.9666}{-59.5166}}{Teniente Irala Fernández é o coração produtivo do Chaco Central em
Presidente Hayes. Oferece um isolamento demográfico que beira a
invulnerabilidade social, protegida pela imensidão de seu território.
Sua força absoluta reside na abundância de proteína animal e na
localização tática sobre a Transchaco. Contudo, a sobrevivência na
localidade é um desafio técnico permanente devido à escassez de água
doce e ao calor extremo. É o local ideal para projetos de
autossuficiência de larga escala (fazendas menonitas ou estâncias
modernas) que possuam capital para investir em segurança hídrica e
autonomia energética radical.

\subsubsection{Por que sim}

\subsubsection{Por que não}

\subsubsection{Recomendação
Técnica}

Altamente indicado para estabelecimentos de autossuficiência
agropecuária que busquem isolamento tático absoluto. Requer investimento
obrigatório em reservatórios de água vedados e sistemas de energia solar
potentes. O GSS de 6.7 reflete a segurança do isolamento e a riqueza de
proteína em contraste com a hostilidade climática e hídrica. Referência
atual de score: GSS 6.7.

\begin{itemize}
\tightlist
\item
  \begin{enumerate}
  \def\labelenumi{\arabic{enumi})}
  \tightlist
  \item
    Dados censitários validados (INE\footnote{\fineINE} 2022).
  \end{enumerate}
\item
  \begin{enumerate}
  \def\labelenumi{\arabic{enumi})}
  \setcounter{enumi}{1}
  \tightlist
  \item
    Relatórios de operação do Aqueduto do Chaco (MOPC/\allowbreak ESSAP\footnote{\fineESSAP}).
  \end{enumerate}
\item
  \begin{enumerate}
  \def\labelenumi{\arabic{enumi})}
  \setcounter{enumi}{2}
  \tightlist
  \item
    Mapa de solos e aptidão pecuária do Chaco Central (MAG).
  \end{enumerate}
\item
  \begin{enumerate}
  \def\labelenumi{\arabic{enumi})}
  \setcounter{enumi}{3}
  \tightlist
  \item
    Registro de segurança pública departamental.
  \end{enumerate}
\end{itemize}}{dist:15_Presidente_Hayes:Teniente_Irala_Fernandez}
\begin{table}[ht!]
\small \begin{tabular}{p{3.0cm}p{0.8cm}p{6.0cm}} \toprule
\textbf{Critério} & \textbf{Nota} & \textbf{Justificativa} \\ \midrule
A - Ameaças Estratégicas & 6.5 & Localização remota em eixo logístico vital; boa invisibilidade tática fora da rodovia \\
B - Risco Social & 8.5 & Densidade populacional extremamente baixa em território vasto; isolamento social superior \\
C - Riscos Naturais & 4.5 & Risco climático extremo; nota penalizada pela seca crônica e inundações repentinas \\
D - Autossuficiência & 7.5 & Recursos de proteína animal e leite massivos; limitado pela escassez de água doce \\
E - Institucional & 6.0 & Ambiente social seguro e estável; suporte institucional funcional para o contexto do Chaco \\
\midrule \textbf{GSS FINAL} & \textbf{6.7} & \textbf{Moderadamente Seguro (Ideal para quem busca isolamento pecuário em larga escala, exigindo alta autonomia hídrica e logística).} \\ \bottomrule \end{tabular}
\end{table}
\subsubsection{Vulnerabilidades e
Mitigação}

\begin{itemize}
\tightlist
\item
  \textbf{Instabilidade Regional:} Presença de grupos armados no
  departamento (Mitigação: residência no perímetro urbano controlado ou
  segurança privada dedicada em zonas rurais).
\item
  \textbf{Risco Hídrico:} Inundações severas (Mitigação: evitar áreas de
  cota baixa, mapear rotas de saída para o Chaco ou Sul).
\item
  \textbf{Pico Logístico:} Dependência de poucas pontes para cruzar o
  Rio Paraguai.
\end{itemize}

\subsubsection{Dossiê de Campo}
\begin{description}[style=nextline,leftmargin=0.5cm,font=\bfseries]
\item[Coordenadas]  23°24'S 57°26'W.
\item[Topografia]  Planície próxima ao Rio Paraguai. Altitude \textasciitilde 70m. Região tectonicamente estável, risco sísmico quase nulo.
\item[Alvos Estratégicos]  Porto de Concepción (logística de grãos e cimento), Aeroporto ANC (em modernização), Sede do CODI (Comando de Operações de Defesa Interna), Regimento de Infantaria N° 10 "Sauce". Planta da Paracel (Celulose) em implantação.
\item[Fallout]  Ventos predominantes do Norte/\allowbreak Nordeste (quentes) e Sul (frios no inverno).
\item[População]  73.360 (Censo 2022 Final).
\item[Densidade]  \textasciitilde 85 hab/\allowbreak km² (distrito total), zona urbana densa.
\item[Vias de Acesso]  Ruta PY05 (Conexão Leste-Oeste), Ponte sobre o Rio Paraguai. Gargalos logísticos devido à dependência de poucas travessias fluviais.
\item[Serviços]  Gran Hospital de Concepción (em construção, referência para o Norte). Sede administrativa regional.
\item[Custo de Vida]  US\$ 800-1.200/\allowbreak mês. Lotes urbanos: US\$ 15-30/\allowbreak m². Terras rurais: US\$ 5.000-8.000/\allowbreak ha.
\item[Hidrologia]  Risco Crítico de Inundação. Bairros ribeirinhos vulneráveis a cheias do Rio Paraguai e transbordamentos súbitos do Rio Ypané.
\item[Clima]  Tropical Úmido. Verões com temperaturas extremas. Tempestades intensas frequentes entre outubro e maio.
\item[Inclinação recomendada para placas solares]  23° voltado para o Norte (anual); 33° no inverno (Jun-Ago); 13° no verão (Nov-Jan). Base: latitude local -23.40°S.
\item[Energia]  Rede nacional (Itaipu), com melhorias recentes mas picos de instabilidade no verão.
\end{description}
\textbf{Fonte climática:} NASA POWER Climatology API (período 2001-2020)

\textbf{Fonte luminosa:} estimativa world atlas

\paragraph{Irradiação Solar
(kWh/\allowbreak m²/\allowbreak dia)}

{\def\LTcaptype{none} % do not increment counter
\begin{longtable}[]{@{}
  >{\raggedright\arraybackslash}p{(\linewidth - 24\tabcolsep) * \real{0.0746}}
  >{\raggedright\arraybackslash}p{(\linewidth - 24\tabcolsep) * \real{0.0746}}
  >{\raggedright\arraybackslash}p{(\linewidth - 24\tabcolsep) * \real{0.0746}}
  >{\raggedright\arraybackslash}p{(\linewidth - 24\tabcolsep) * \real{0.0746}}
  >{\raggedright\arraybackslash}p{(\linewidth - 24\tabcolsep) * \real{0.0746}}
  >{\raggedright\arraybackslash}p{(\linewidth - 24\tabcolsep) * \real{0.0746}}
  >{\raggedright\arraybackslash}p{(\linewidth - 24\tabcolsep) * \real{0.0746}}
  >{\raggedright\arraybackslash}p{(\linewidth - 24\tabcolsep) * \real{0.0746}}
  >{\raggedright\arraybackslash}p{(\linewidth - 24\tabcolsep) * \real{0.0746}}
  >{\raggedright\arraybackslash}p{(\linewidth - 24\tabcolsep) * \real{0.0746}}
  >{\raggedright\arraybackslash}p{(\linewidth - 24\tabcolsep) * \real{0.0746}}
  >{\raggedright\arraybackslash}p{(\linewidth - 24\tabcolsep) * \real{0.0746}}
  >{\raggedright\arraybackslash}p{(\linewidth - 24\tabcolsep) * \real{0.1045}}@{}}
\toprule\noalign{}
\begin{minipage}[b]{\linewidth}\raggedright
Jan
\end{minipage} & \begin{minipage}[b]{\linewidth}\raggedright
Fev
\end{minipage} & \begin{minipage}[b]{\linewidth}\raggedright
Mar
\end{minipage} & \begin{minipage}[b]{\linewidth}\raggedright
Abr
\end{minipage} & \begin{minipage}[b]{\linewidth}\raggedright
Mai
\end{minipage} & \begin{minipage}[b]{\linewidth}\raggedright
Jun
\end{minipage} & \begin{minipage}[b]{\linewidth}\raggedright
Jul
\end{minipage} & \begin{minipage}[b]{\linewidth}\raggedright
Ago
\end{minipage} & \begin{minipage}[b]{\linewidth}\raggedright
Set
\end{minipage} & \begin{minipage}[b]{\linewidth}\raggedright
Out
\end{minipage} & \begin{minipage}[b]{\linewidth}\raggedright
Nov
\end{minipage} & \begin{minipage}[b]{\linewidth}\raggedright
Dez
\end{minipage} & \begin{minipage}[b]{\linewidth}\raggedright
Média
\end{minipage} \\
\midrule\noalign{}
\endhead
\bottomrule\noalign{}
\endlastfoot
6.68 & 6.07 & 5.58 & 4.63 & 3.47 & 3.00 & 3.30 & 4.09 & 4.68 & 5.49 &
6.34 & 6.61 & \textbf{5.0} \\
\end{longtable}
}

\textbf{Inclinação solar recomendada:} 23° N (anual) · 33° N (inverno
jun-ago) · 13° N (verão nov-jan)

\paragraph{Precipitação (mm/\allowbreak mês)}

{\def\LTcaptype{none} % do not increment counter
\begin{longtable}[]{@{}
  >{\raggedright\arraybackslash}p{(\linewidth - 24\tabcolsep) * \real{0.0704}}
  >{\raggedright\arraybackslash}p{(\linewidth - 24\tabcolsep) * \real{0.0704}}
  >{\raggedright\arraybackslash}p{(\linewidth - 24\tabcolsep) * \real{0.0704}}
  >{\raggedright\arraybackslash}p{(\linewidth - 24\tabcolsep) * \real{0.0704}}
  >{\raggedright\arraybackslash}p{(\linewidth - 24\tabcolsep) * \real{0.0704}}
  >{\raggedright\arraybackslash}p{(\linewidth - 24\tabcolsep) * \real{0.0704}}
  >{\raggedright\arraybackslash}p{(\linewidth - 24\tabcolsep) * \real{0.0704}}
  >{\raggedright\arraybackslash}p{(\linewidth - 24\tabcolsep) * \real{0.0704}}
  >{\raggedright\arraybackslash}p{(\linewidth - 24\tabcolsep) * \real{0.0704}}
  >{\raggedright\arraybackslash}p{(\linewidth - 24\tabcolsep) * \real{0.0704}}
  >{\raggedright\arraybackslash}p{(\linewidth - 24\tabcolsep) * \real{0.0704}}
  >{\raggedright\arraybackslash}p{(\linewidth - 24\tabcolsep) * \real{0.0704}}
  >{\raggedright\arraybackslash}p{(\linewidth - 24\tabcolsep) * \real{0.1549}}@{}}
\toprule\noalign{}
\begin{minipage}[b]{\linewidth}\raggedright
Jan
\end{minipage} & \begin{minipage}[b]{\linewidth}\raggedright
Fev
\end{minipage} & \begin{minipage}[b]{\linewidth}\raggedright
Mar
\end{minipage} & \begin{minipage}[b]{\linewidth}\raggedright
Abr
\end{minipage} & \begin{minipage}[b]{\linewidth}\raggedright
Mai
\end{minipage} & \begin{minipage}[b]{\linewidth}\raggedright
Jun
\end{minipage} & \begin{minipage}[b]{\linewidth}\raggedright
Jul
\end{minipage} & \begin{minipage}[b]{\linewidth}\raggedright
Ago
\end{minipage} & \begin{minipage}[b]{\linewidth}\raggedright
Set
\end{minipage} & \begin{minipage}[b]{\linewidth}\raggedright
Out
\end{minipage} & \begin{minipage}[b]{\linewidth}\raggedright
Nov
\end{minipage} & \begin{minipage}[b]{\linewidth}\raggedright
Dez
\end{minipage} & \begin{minipage}[b]{\linewidth}\raggedright
Total/\allowbreak ano
\end{minipage} \\
\midrule\noalign{}
\endhead
\bottomrule\noalign{}
\endlastfoot
157 & 151 & 120 & 139 & 116 & 71 & 48 & 31 & 73 & 151 & 204 & 167 &
\textbf{1432 mm} \\
\end{longtable}
}

\paragraph{Poluição Luminosa}

{\def\LTcaptype{none} % do not increment counter
\begin{longtable}[]{@{}ll@{}}
\toprule\noalign{}
Parâmetro & Valor \\
\midrule\noalign{}
\endhead
\bottomrule\noalign{}
\endlastfoot
Escala Bortle & 4 --- Céu rural-suburbano \\
Radiância artificial & 3.0 nW/\allowbreak cm²/\allowbreak sr \\
\end{longtable}
}
\paragraph{Solo (SoilGrids 2.0, média ponderada 0--30
cm)}

{\def\LTcaptype{none} % do not increment counter
\begin{longtable}[]{@{}ll@{}}
\toprule\noalign{}
Parâmetro & Valor \\
\midrule\noalign{}
\endhead
\bottomrule\noalign{}
\endlastfoot
pH (H$_2$O) & 5.8 \\
Carbono orgânico (SOC) & 14.59 g/\allowbreak kg \\
Argila & 22.31 \% \\
Areia & 52.33 \% \\
Silte (calc.) & 25.4 \% \\
Densidade aparente & 1.35 g/\allowbreak cm³ \\
\textbf{Aptidão agrícola} & \textbf{Alta} \\
\end{longtable}
}

\begin{itemize}
\tightlist
\item
  \textbf{Autossuficiência:} Elevada em proteína animal (indústria
  cárnea) e potencial agrícola de exportação.
\end{itemize}
\subsubsection{Indicadores Sociais}

\textbf{Fontes:} PNUD 2020, INE\footnote{\fineINE} Censo 2022, INE\footnote{\fineINE} EPHC 2023, Ministerio
Público 2024\\
\textbf{Nota:} Valores marcados como \emph{(dept.)} referem-se ao
departamento de Concepción; valores \emph{(dist.)} são específicos deste
distrito. Dados marcados \emph{(est.)} são estimativas.

{\def\LTcaptype{none} % do not increment counter
\begin{longtable}[]{@{}
  >{\raggedright\arraybackslash}p{(\linewidth - 4\tabcolsep) * \real{0.4231}}
  >{\raggedright\arraybackslash}p{(\linewidth - 4\tabcolsep) * \real{0.2692}}
  >{\raggedright\arraybackslash}p{(\linewidth - 4\tabcolsep) * \real{0.3077}}@{}}
\toprule\noalign{}
\begin{minipage}[b]{\linewidth}\raggedright
Indicador
\end{minipage} & \begin{minipage}[b]{\linewidth}\raggedright
Valor
\end{minipage} & \begin{minipage}[b]{\linewidth}\raggedright
Âmbito
\end{minipage} \\
\midrule\noalign{}
\endhead
\bottomrule\noalign{}
\endlastfoot
IDH (2020) & 0.658 & dist. (est.) \\
Ranking IDH nacional & 9/\allowbreak 18 & dept. \\
Esperança de vida & 72,4 anos & dept. \\
Escolaridade média & 10.5 anos & dist. \\
RNB per capita (USD PPA) & 8.400 & dept. \\
Pobreza monetária (\%) & 31,1\% & dept. \\
Pobreza extrema (\%) & 7,8\% & dept. \\
Índice de Gini & 0,460 & dept. \\
Acesso a água potável (\%) & 74,2\% & dept. \\
Acesso a saneamento (\%) & 71,8\% & dept. \\
Taxa de homicídios (est., /\allowbreak 100k hab) & \textasciitilde12--15
(estimativa, acima da média nacional) & dept. \\
Índice de segurança & baixo & dept. \\
População (Censo 2022) & 7.336 hab. & dist. \\
Idade mediana & 27.0 anos & dist. \\
Taxa de fecundidade & 2 & dist. \\
\end{longtable}
}
\subsubsection{Combustível}

\textbf{Referência:} PETROPAR /\allowbreak  postos locais (2024) \textbf{Tipo de
localidade:} Capital departamental

{\def\LTcaptype{none} % do not increment counter
\begin{longtable}[]{@{}lll@{}}
\toprule\noalign{}
Combustível & USD/\allowbreak litro & Gs/\allowbreak litro (aprox.) \\
\midrule\noalign{}
\endhead
\bottomrule\noalign{}
\endlastfoot
Gasolina 93 oct & 1.01 & 7,474 \\
Gasolina 97 oct (premium) & 1.11 & 8,214 \\
Diesel & 0.94 & 6,956 \\
\end{longtable}
}

\begin{quote}
Preços podem variar ±5\% conforme posto e sazonalidade. Chaco e interior
remoto apresentam maior variação.
\end{quote}

\subsubsection{Cobertura Celular}

\textbf{Fonte:} CONATEL PY /\allowbreak  operadoras (2024)

{\def\LTcaptype{none} % do not increment counter
\begin{longtable}[]{@{}ll@{}}
\toprule\noalign{}
Parâmetro & Valor \\
\midrule\noalign{}
\endhead
\bottomrule\noalign{}
\endlastfoot
Cobertura 4G população (dept.) & 78\% \\
Cobertura 4G área rural & 55\% \\
Melhor operadora & Tigo \\
Qualidade rural & moderada \\
\end{longtable}
}

\begin{quote}
Para áreas rurais fora do núcleo urbano, recomenda-se chip Tigo como
principal e Personal como backup.
\end{quote}

\subsubsection{Internet}

\textbf{Fonte:} CONATEL /\allowbreak  Speedtest Ookla (2024)

{\def\LTcaptype{none} % do not increment counter
\begin{longtable}[]{@{}ll@{}}
\toprule\noalign{}
Parâmetro & Valor \\
\midrule\noalign{}
\endhead
\bottomrule\noalign{}
\endlastfoot
Velocidade média download & 38 Mbps \\
Domicílios com internet (dept.) & 48\% \\
Tecnologia predominante & rádio \\
Opção rural & Starlink disponível (\textasciitilde USD 44/\allowbreak mês) \\
\end{longtable}
}

\subsubsection{Mercado Imobiliário e Terra
Rural}

\textbf{Fonte:} INDERT /\allowbreak  Clasificados.com.py (2024)

{\def\LTcaptype{none} % do not increment counter
\begin{longtable}[]{@{}ll@{}}
\toprule\noalign{}
Tipo & Referência \\
\midrule\noalign{}
\endhead
\bottomrule\noalign{}
\endlastfoot
Terra agrícola alta prod. (USD/\allowbreak ha) & 3,500 \\
Imóvel urbano (USD/\allowbreak m²) & 747 \\
Aluguel 2 quartos (USD/\allowbreak mês) & 312 \\
\end{longtable}
}

\begin{quote}
Valores de referência departamental. Localidades menores podem ter
preços 20--40\% abaixo da capital departamental.
\end{quote}

\subsubsection{Saúde}

\textbf{Fonte:} MSPBS /\allowbreak  IPS Paraguay (2024)

{\def\LTcaptype{none} % do not increment counter
\begin{longtable}[]{@{}ll@{}}
\toprule\noalign{}
Serviço & Disponibilidade \\
\midrule\noalign{}
\endhead
\bottomrule\noalign{}
\endlastfoot
USF /\allowbreak  Posto de Saúde & sim \\
Hospital Regional & sim \\
IPS (seguro social) & sim \\
Farmácia & sim \\
Distância ao hospital de referência & local \\
\end{longtable}
}

\textbf{Principais estabelecimentos:} Hospital Regional de Concepción
(referência para todo o Norte), IPS Concepción, clínicas privadas
(CEDIMEC, Sanatorio San Antonio), Gran Hospital de Concepción (em
implantação)

\textbf{Observação para imigrantes:} IPS acessível mediante contrato
formal de trabalho. Cobertura privada recomendada para maior agilidade;
serviços públicos funcionais mas com demanda elevada. Imigrantes sem
vínculo empregatício devem contratar seguro privado.
\newpage
\secaoDiagnostico{Villa Hayes}{\mapaDistrito{1504}}{}{dist:15_Presidente_Hayes:Villa_Hayes}
\begin{table}[ht!]
\small \begin{tabular}{p{3.0cm}p{0.8cm}p{6.0cm}} \toprule
\textbf{Critério} & \textbf{Nota} & \textbf{Justificativa} \\ \midrule
A - Ameaças Estratégicas & 4.5 & Ponto de entrada estratégico para o Chaco; alvos industriais/\allowbreak elétricos de altíssima visibilidade tática \\
B - Risco Social & 6.5 & Capital departamental com urbanização concentrada; risco social em crises metropolitana \\
C - Riscos Naturais & 6.5 & Geologia estável; nota penalizada por risco de inundação e histórico de sismos 3.5-4.6 \\
D - Autossuficiência & 8.5 & Recursos excepcionais: hub siderúrgico, elétrico e abundância de proteína animal \\
E - Institucional & 8.0 & Forte segurança institucional e serviços de saúde de referência regional \\
\midrule \textbf{GSS FINAL} & \textbf{6.7} & \textbf{Moderadamente Seguro (Indicado como hub logístico e operacional de alta resiliência, requerendo planos de contingência para visibilidade estratégica).} \\ \bottomrule \end{tabular}
\end{table}
\subsubsection{Vulnerabilidades e
Mitigação}

\begin{itemize}
\tightlist
\item
  \textbf{Instabilidade Regional:} Atividade de grupos irregulares em
  áreas rurais remotas (Mitigação: residência próxima ao centro urbano
  ou segurança privada em fazendas mecanizadas).
\item
  \textbf{Dependência de Monocultura:} Economia sensível a preços de
  commodities (Mitigação: implementação de hortas de subsistência e
  diversificação de ativos).
\item
  \textbf{Risco Logístico:} Controle militar de rodovias em cenários de
  crise (Mitigação: manutenção de veículos off-road e rotas secundárias
  mapeadas).
\end{itemize}

\subsubsection{Dossiê de Campo}
\begin{description}[style=nextline,leftmargin=0.5cm,font=\bfseries]
\item[Coordenadas]  23°20'S 57°03'W.
\item[Topografia]  Relevo ondulado ("lomadas"), solo profundo e bem drenado. Altitude \textasciitilde 170m. Geologicamente estável, risco sísmico muito baixo.
\item[Alvos Estratégicos]  Cruzamento das Rutas PY05 e rodovias locais; Polo de processamento de Estévia (Ka'a He'e) e óleos vegetais; Sede tática da FTC (Força de Tarefa Conjunta) e base de operações do CODI.
\item[Fallout]  Ventos predominantes do Norte (quentes) e Sul/\allowbreak Sudeste (frios).
\item[População]  39.548 habitantes (Censo 2022 Final). Segunda cidade mais importante do departamento.
\item[Densidade]  \textasciitilde 31,4 hab/\allowbreak km² (Moderada, com núcleo urbano dinâmico e vasta área rural produtiva).
\item[Vias de Acesso]  Conectividade asfáltica de alta qualidade via Ruta PY05 (Concepción-Pedro Juan Caballero). Eixo logístico fundamental para o agronegócio do Norte.
\item[Serviços]  Hospital Distrital de Horqueta, Unidade do IPS e clínicas privadas (CEDIMEC, San Antonio). Boa oferta de comércio e serviços financeiros.
\item[Custo de Vida]  Baixo a Médio; economia baseada no agronegócio exportador e agricultura familiar.
\item[Preço da Terra]  US\$ 4.000-5.500/\allowbreak ha (terras mecanizáveis de "tierra roja").
\item[Hidrologia]  Baixo risco de inundações em comparação com o litoral do Rio Paraguai. Drenagem natural eficiente pelo relevo ondulado.
\item[Clima]  Tropical. Temperaturas extremas no verão (>40°C). Chuvas bem distribuídas, com pico entre novembro e janeiro.
\item[Energia]  Rede nacional (Itaipu). Subestação local com manutenção prioritária devido à agroindústria.
\item[Água]  Abundante via poços artesianos e arroios locais.
\item[Qualidade do Solo]  Latossolo Vermelho ("Tierra Roja"); alta fertilidade e excelente aptidão para agricultura mecanizada (soja, milho, stevia).
\item[Recursos Locais]  Capital nacional da Estévia; grande produção de proteína animal (vizinhança com frigoríficos de Concepción) e grãos. Elevadíssimo potencial de autossuficiência.
\item[Segurança]  Zona de segurança tática. Presença ostensiva da FTC devido ao histórico de insurgência (EPP\footnote{\fineEPP}) na região rural. Taxa de homicídios regional de \textasciitilde 20/\allowbreak 100k, influenciada por conflitos agrários. Perímetro urbano considerado seguro e controlado.
\item[Leis Local: Município consolidado. Restrição de Fronteira]  Aplicação da Lei 2532/\allowbreak 05 exige verificação de distância da linha de fronteira para estrangeiros limítrofes.
\end{description}
\textbf{Fonte climática:} NASA POWER Climatology API (período 2001-2020)

\textbf{Fonte luminosa:} estimativa world atlas

\paragraph{Irradiação Solar
(kWh/\allowbreak m²/\allowbreak dia)}

{\def\LTcaptype{none} % do not increment counter
\begin{longtable}[]{@{}
  >{\raggedright\arraybackslash}p{(\linewidth - 24\tabcolsep) * \real{0.0746}}
  >{\raggedright\arraybackslash}p{(\linewidth - 24\tabcolsep) * \real{0.0746}}
  >{\raggedright\arraybackslash}p{(\linewidth - 24\tabcolsep) * \real{0.0746}}
  >{\raggedright\arraybackslash}p{(\linewidth - 24\tabcolsep) * \real{0.0746}}
  >{\raggedright\arraybackslash}p{(\linewidth - 24\tabcolsep) * \real{0.0746}}
  >{\raggedright\arraybackslash}p{(\linewidth - 24\tabcolsep) * \real{0.0746}}
  >{\raggedright\arraybackslash}p{(\linewidth - 24\tabcolsep) * \real{0.0746}}
  >{\raggedright\arraybackslash}p{(\linewidth - 24\tabcolsep) * \real{0.0746}}
  >{\raggedright\arraybackslash}p{(\linewidth - 24\tabcolsep) * \real{0.0746}}
  >{\raggedright\arraybackslash}p{(\linewidth - 24\tabcolsep) * \real{0.0746}}
  >{\raggedright\arraybackslash}p{(\linewidth - 24\tabcolsep) * \real{0.0746}}
  >{\raggedright\arraybackslash}p{(\linewidth - 24\tabcolsep) * \real{0.0746}}
  >{\raggedright\arraybackslash}p{(\linewidth - 24\tabcolsep) * \real{0.1045}}@{}}
\toprule\noalign{}
\begin{minipage}[b]{\linewidth}\raggedright
Jan
\end{minipage} & \begin{minipage}[b]{\linewidth}\raggedright
Fev
\end{minipage} & \begin{minipage}[b]{\linewidth}\raggedright
Mar
\end{minipage} & \begin{minipage}[b]{\linewidth}\raggedright
Abr
\end{minipage} & \begin{minipage}[b]{\linewidth}\raggedright
Mai
\end{minipage} & \begin{minipage}[b]{\linewidth}\raggedright
Jun
\end{minipage} & \begin{minipage}[b]{\linewidth}\raggedright
Jul
\end{minipage} & \begin{minipage}[b]{\linewidth}\raggedright
Ago
\end{minipage} & \begin{minipage}[b]{\linewidth}\raggedright
Set
\end{minipage} & \begin{minipage}[b]{\linewidth}\raggedright
Out
\end{minipage} & \begin{minipage}[b]{\linewidth}\raggedright
Nov
\end{minipage} & \begin{minipage}[b]{\linewidth}\raggedright
Dez
\end{minipage} & \begin{minipage}[b]{\linewidth}\raggedright
Média
\end{minipage} \\
\midrule\noalign{}
\endhead
\bottomrule\noalign{}
\endlastfoot
6.68 & 6.07 & 5.58 & 4.63 & 3.47 & 3.00 & 3.30 & 4.09 & 4.68 & 5.49 &
6.34 & 6.61 & \textbf{5.0} \\
\end{longtable}
}

\textbf{Inclinação solar recomendada:} 23° N (anual) · 33° N (inverno
jun-ago) · 13° N (verão nov-jan)

\paragraph{Precipitação (mm/\allowbreak mês)}

{\def\LTcaptype{none} % do not increment counter
\begin{longtable}[]{@{}
  >{\raggedright\arraybackslash}p{(\linewidth - 24\tabcolsep) * \real{0.0704}}
  >{\raggedright\arraybackslash}p{(\linewidth - 24\tabcolsep) * \real{0.0704}}
  >{\raggedright\arraybackslash}p{(\linewidth - 24\tabcolsep) * \real{0.0704}}
  >{\raggedright\arraybackslash}p{(\linewidth - 24\tabcolsep) * \real{0.0704}}
  >{\raggedright\arraybackslash}p{(\linewidth - 24\tabcolsep) * \real{0.0704}}
  >{\raggedright\arraybackslash}p{(\linewidth - 24\tabcolsep) * \real{0.0704}}
  >{\raggedright\arraybackslash}p{(\linewidth - 24\tabcolsep) * \real{0.0704}}
  >{\raggedright\arraybackslash}p{(\linewidth - 24\tabcolsep) * \real{0.0704}}
  >{\raggedright\arraybackslash}p{(\linewidth - 24\tabcolsep) * \real{0.0704}}
  >{\raggedright\arraybackslash}p{(\linewidth - 24\tabcolsep) * \real{0.0704}}
  >{\raggedright\arraybackslash}p{(\linewidth - 24\tabcolsep) * \real{0.0704}}
  >{\raggedright\arraybackslash}p{(\linewidth - 24\tabcolsep) * \real{0.0704}}
  >{\raggedright\arraybackslash}p{(\linewidth - 24\tabcolsep) * \real{0.1549}}@{}}
\toprule\noalign{}
\begin{minipage}[b]{\linewidth}\raggedright
Jan
\end{minipage} & \begin{minipage}[b]{\linewidth}\raggedright
Fev
\end{minipage} & \begin{minipage}[b]{\linewidth}\raggedright
Mar
\end{minipage} & \begin{minipage}[b]{\linewidth}\raggedright
Abr
\end{minipage} & \begin{minipage}[b]{\linewidth}\raggedright
Mai
\end{minipage} & \begin{minipage}[b]{\linewidth}\raggedright
Jun
\end{minipage} & \begin{minipage}[b]{\linewidth}\raggedright
Jul
\end{minipage} & \begin{minipage}[b]{\linewidth}\raggedright
Ago
\end{minipage} & \begin{minipage}[b]{\linewidth}\raggedright
Set
\end{minipage} & \begin{minipage}[b]{\linewidth}\raggedright
Out
\end{minipage} & \begin{minipage}[b]{\linewidth}\raggedright
Nov
\end{minipage} & \begin{minipage}[b]{\linewidth}\raggedright
Dez
\end{minipage} & \begin{minipage}[b]{\linewidth}\raggedright
Total/\allowbreak ano
\end{minipage} \\
\midrule\noalign{}
\endhead
\bottomrule\noalign{}
\endlastfoot
156 & 151 & 118 & 140 & 118 & 74 & 50 & 32 & 76 & 154 & 205 & 165 &
\textbf{1446 mm} \\
\end{longtable}
}

\paragraph{Poluição Luminosa}

{\def\LTcaptype{none} % do not increment counter
\begin{longtable}[]{@{}ll@{}}
\toprule\noalign{}
Parâmetro & Valor \\
\midrule\noalign{}
\endhead
\bottomrule\noalign{}
\endlastfoot
Escala Bortle & 3 --- Céu rural \\
Radiância artificial & 1.5 nW/\allowbreak cm²/\allowbreak sr \\
\end{longtable}
}
\paragraph{Solo (SoilGrids 2.0, média ponderada 0--30
cm)}

{\def\LTcaptype{none} % do not increment counter
\begin{longtable}[]{@{}ll@{}}
\toprule\noalign{}
Parâmetro & Valor \\
\midrule\noalign{}
\endhead
\bottomrule\noalign{}
\endlastfoot
pH (H$_2$O) & 5.78 \\
Carbono orgânico (SOC) & 16.63 g/\allowbreak kg \\
Argila & 17.8 \% \\
Areia & 62.77 \% \\
Silte (calc.) & 19.4 \% \\
Densidade aparente & 1.27 g/\allowbreak cm³ \\
\textbf{Aptidão agrícola} & \textbf{Média} \\
\end{longtable}
}

\begin{itemize}
\tightlist
\item
  \textbf{Homicidios/\allowbreak 100k:} \textasciitilde20.0 (Regional/\allowbreak Rural).
\end{itemize}
\subsubsection{Indicadores Sociais}

\textbf{Fontes:} PNUD 2020, INE\footnote{\fineINE} Censo 2022, INE\footnote{\fineINE} EPHC 2023, Ministerio
Público 2024\\
\textbf{Nota:} Valores marcados como \emph{(dept.)} referem-se ao
departamento de Concepción; valores \emph{(dist.)} são específicos deste
distrito. Dados marcados \emph{(est.)} são estimativas.

{\def\LTcaptype{none} % do not increment counter
\begin{longtable}[]{@{}
  >{\raggedright\arraybackslash}p{(\linewidth - 4\tabcolsep) * \real{0.4231}}
  >{\raggedright\arraybackslash}p{(\linewidth - 4\tabcolsep) * \real{0.2692}}
  >{\raggedright\arraybackslash}p{(\linewidth - 4\tabcolsep) * \real{0.3077}}@{}}
\toprule\noalign{}
\begin{minipage}[b]{\linewidth}\raggedright
Indicador
\end{minipage} & \begin{minipage}[b]{\linewidth}\raggedright
Valor
\end{minipage} & \begin{minipage}[b]{\linewidth}\raggedright
Âmbito
\end{minipage} \\
\midrule\noalign{}
\endhead
\bottomrule\noalign{}
\endlastfoot
IDH (2020) & 0.648 & dist. (est.) \\
Ranking IDH nacional & 9/\allowbreak 18 & dept. \\
Esperança de vida & 72,4 anos & dept. \\
Escolaridade média & 8.8 anos & dist. \\
RNB per capita (USD PPA) & 8.400 & dept. \\
Pobreza monetária (\%) & 31,1\% & dept. \\
Pobreza extrema (\%) & 7,8\% & dept. \\
Índice de Gini & 0,460 & dept. \\
Acesso a água potável (\%) & 74,2\% & dept. \\
Acesso a saneamento (\%) & 71,8\% & dept. \\
Taxa de homicídios (est., /\allowbreak 100k hab) & \textasciitilde12--15
(estimativa, acima da média nacional) & dept. \\
Índice de segurança & baixo & dept. \\
População (Censo 2022) & 39.548 hab. & dist. \\
Idade mediana & 28.0 anos & dist. \\
Taxa de fecundidade & 2 & dist. \\
\end{longtable}
}
\subsubsection{Combustível}

\textbf{Referência:} PETROPAR /\allowbreak  postos locais (2024) \textbf{Tipo de
localidade:} Interior

{\def\LTcaptype{none} % do not increment counter
\begin{longtable}[]{@{}lll@{}}
\toprule\noalign{}
Combustível & USD/\allowbreak litro & Gs/\allowbreak litro (aprox.) \\
\midrule\noalign{}
\endhead
\bottomrule\noalign{}
\endlastfoot
Gasolina 93 oct & 1.01 & 7,474 \\
Gasolina 97 oct (premium) & 1.11 & 8,214 \\
Diesel & 0.94 & 6,956 \\
\end{longtable}
}

\begin{quote}
Preços podem variar ±5\% conforme posto e sazonalidade. Chaco e interior
remoto apresentam maior variação.
\end{quote}

\subsubsection{Cobertura Celular}

\textbf{Fonte:} CONATEL PY /\allowbreak  operadoras (2024)

{\def\LTcaptype{none} % do not increment counter
\begin{longtable}[]{@{}ll@{}}
\toprule\noalign{}
Parâmetro & Valor \\
\midrule\noalign{}
\endhead
\bottomrule\noalign{}
\endlastfoot
Cobertura 4G população (dept.) & 78\% \\
Cobertura 4G área rural & 55\% \\
Melhor operadora & Tigo \\
Qualidade rural & moderada \\
\end{longtable}
}

\begin{quote}
Para áreas rurais fora do núcleo urbano, recomenda-se chip Tigo como
principal e Personal como backup.
\end{quote}

\subsubsection{Internet}

\textbf{Fonte:} CONATEL /\allowbreak  Speedtest Ookla (2024)

{\def\LTcaptype{none} % do not increment counter
\begin{longtable}[]{@{}ll@{}}
\toprule\noalign{}
Parâmetro & Valor \\
\midrule\noalign{}
\endhead
\bottomrule\noalign{}
\endlastfoot
Velocidade média download & 38 Mbps \\
Domicílios com internet (dept.) & 48\% \\
Tecnologia predominante & rádio \\
Opção rural & Starlink disponível (\textasciitilde USD 44/\allowbreak mês) \\
\end{longtable}
}

\subsubsection{Mercado Imobiliário e Terra
Rural}

\textbf{Fonte:} INDERT /\allowbreak  Clasificados.com.py (2024)

{\def\LTcaptype{none} % do not increment counter
\begin{longtable}[]{@{}ll@{}}
\toprule\noalign{}
Tipo & Referência \\
\midrule\noalign{}
\endhead
\bottomrule\noalign{}
\endlastfoot
Terra agrícola alta prod. (USD/\allowbreak ha) & 3,500 \\
Imóvel urbano (USD/\allowbreak m²) & 650 \\
Aluguel 2 quartos (USD/\allowbreak mês) & 260 \\
\end{longtable}
}

\begin{quote}
Valores de referência departamental. Localidades menores podem ter
preços 20--40\% abaixo da capital departamental.
\end{quote}

\subsubsection{Saúde}

\textbf{Fonte:} MSPBS /\allowbreak  IPS Paraguay (2024-2026), consolidação
departamental e proxy local

{\def\LTcaptype{none} % do not increment counter
\begin{longtable}[]{@{}ll@{}}
\toprule\noalign{}
Serviço & Disponibilidade \\
\midrule\noalign{}
\endhead
\bottomrule\noalign{}
\endlastfoot
USF /\allowbreak  Posto de Saúde & sim \\
Hospital Regional & sim \\
IPS (seguro social) & sim \\
Farmácia & sim \\
Distância ao hospital de referência & \textasciitilde51 km
(Concepción) \\
\end{longtable}
}

\textbf{Principais estabelecimentos:} USF local; referencia hospitalar
em Concepción

\textbf{Observação para imigrantes:} Atendimento primario local ou em
raio curto; casos de maior complexidade seguem para Concepción.
Cobertura privada continua recomendada para especialidades e urgências
de maior complexidade.
